% Santa Clara University PhD Thesis Template
% Document class: report, 12pt, letter paper, single-sided
%
% Structure: 10 chapters organized by estimator fit order (zeroth/first/
% second), per the fit-order-ladder outline agreed 2026-09-29. This
% supersedes the 6-chapter, per-paper structure in thesis.tex (kept
% alongside, untouched, as a prose source to port from).
%
% Section/subsection headers below are SKELETON ONLY. Each carries a
% %% SOURCE comment pointing at where the real prose should be ported
% from (thesis.tex line range, or a paper file + section name) so a
% later pass can move content mechanically rather than redraft it.
% Figures are placeholder \fbox blocks pending real figure generation.
\documentclass[12pt, letterpaper, oneside]{report}

% -----------------------------------------------------------------------
% Packages
% -----------------------------------------------------------------------
\usepackage[left=1.5in, right=1in, top=1in, bottom=1in, headheight=15pt]{geometry}
\usepackage{setspace}
\usepackage{amsmath, amsfonts, amssymb}
\usepackage{amsthm}
\usepackage{graphicx}
\graphicspath{{figures/}}
\usepackage{array}
\usepackage{caption}
\usepackage{algorithm}
\usepackage{algorithmic}
\usepackage{fancyhdr}
\usepackage{booktabs}
\usepackage{subcaption}
\usepackage{tikz}
\usepackage{url}
\usepackage{textcomp}
\usepackage[hidelinks]{hyperref}

% -----------------------------------------------------------------------
% Theorem environments -- one running counter, tied to chapter number
% -----------------------------------------------------------------------
\newtheorem{theorem}{Theorem}[chapter]
\newtheorem{lemma}[theorem]{Lemma}
\newtheorem{proposition}[theorem]{Proposition}
\newtheorem{corollary}[theorem]{Corollary}
\newtheorem{remark}[theorem]{Remark}
\newtheorem{assumption}[theorem]{Assumption}

% -----------------------------------------------------------------------
% Thesis metadata -- edit these fields
% -----------------------------------------------------------------------
\newcommand{\thesistitle}{Distributed Multirobot Estimation and Navigation in Vector Fields: From Hardware Testbed to Second-Order Field Structure}
\newcommand{\thesisauthor}{Christopher J. Waight}
\newcommand{\thesisdepartment}{Mechanical Engineering}
\newcommand{\thesisdegreefield}{Mechanical Engineering}
\newcommand{\thesisadvisor}{Christopher A. Kitts}
\newcommand{\thesisreaderA}{Reader Name A}
\newcommand{\thesisreaderB}{Reader Name B}
\newcommand{\thesisreaderC}{Reader Name C}
\newcommand{\thesisreaderD}{Reader Name D}
\newcommand{\thesisdeptchair}{Department Chair Name}
\newcommand{\thesisyear}{2026}
\newcommand{\thesismonth}{Month}
\newcommand{\thesisday}{Day}

% -----------------------------------------------------------------------
% Page style setup
% -----------------------------------------------------------------------
\pagestyle{fancy}
\fancyhf{}
\fancyhead[R]{\thepage}
\renewcommand{\headrulewidth}{0pt}

\fancypagestyle{plain}{%
  \fancyhf{}
  \fancyhead[R]{\thepage}
  \renewcommand{\headrulewidth}{0pt}
}

% -----------------------------------------------------------------------
% TOC / LOF / LOT formatting
% -----------------------------------------------------------------------
\renewcommand{\contentsname}{Table of Contents}
\renewcommand{\listfigurename}{List of Figures}
\renewcommand{\listtablename}{List of Tables}

% -----------------------------------------------------------------------
% Hyphenation hints
% -----------------------------------------------------------------------
\hyphenation{multi-robot cluster-space Jacobian separ-atrix}

% -----------------------------------------------------------------------
% Placeholder figure macro
% -----------------------------------------------------------------------
\newcommand{\placeholderfig}[1]{\fbox{\rule{0pt}{2in}\rule{0.7\textwidth}{0pt}}\\[-1.6in]\centering\textit{[placeholder: #1]}\\[1.4in]}

% -----------------------------------------------------------------------
% Document
% -----------------------------------------------------------------------
\begin{document}

% Front matter uses roman numerals
\pagenumbering{roman}
\setcounter{page}{1}

% ===================================================================
% 1. APPROVAL / SIGNATURE PAGE
% ===================================================================
\begin{titlepage}
\begin{center}
  \large
  Santa Clara University \\[6pt]
  Department of \thesisdepartment

  \vspace{1.5in}

  \normalsize
  Date: \thesismonth\ \thesisday, \thesisyear

  \vspace{1in}

  I HEREBY RECOMMEND THAT THE THESIS PREPARED UNDER MY SUPERVISION BY

  \vspace{0.5in}
  \MakeUppercase{\thesisauthor}
  \vspace{0.5in}

  ENTITLED

  \vspace{0.4in}
  \textit{\thesistitle}
  \vspace{0.4in}

  BE ACCEPTED IN PARTIAL FULFILLMENT OF THE REQUIREMENTS FOR THE DEGREE

  \vspace{0.3in}
  \textbf{DOCTOR OF PHILOSOPHY IN \MakeUppercase{\thesisdegreefield}}
\end{center}

\vspace{0.8in}

\noindent
\begin{tabular}{p{3in} p{0.3in} p{3in}}
  \hrulefill & & \hrulefill \\[2pt]
  Thesis Advisor & & Thesis Reader \\[24pt]
  \hrulefill & & \hrulefill \\[2pt]
  Chairman of Department & & Thesis Reader \\[24pt]
  & & \hrulefill \\[2pt]
  & & Thesis Reader \\[24pt]
  & & \hrulefill \\[2pt]
  & & Thesis Reader
\end{tabular}

\end{titlepage}

% ===================================================================
% 2. TITLE PAGE
% ===================================================================
\begin{titlepage}
\begin{center}
  \vspace*{1in}

  {\Large \thesistitle}

  \vspace{0.6in}
  by \\[6pt]
  \thesisauthor

  \vspace{1.2in}
  Dissertation \\[6pt]
  Submitted in Partial Fulfillment of the Requirements \\
  for the Degree of Doctor of Philosophy \\
  in \thesisdegreefield \\[6pt]
  in the School of Engineering at \\
  Santa Clara University, \thesisyear

  \vspace{1in}
  Santa Clara, California \\
  \thesisyear
\end{center}
\end{titlepage}

% ===================================================================
% 3. COPYRIGHT
% ===================================================================
\clearpage
\thispagestyle{plain}
\vspace*{3in}
\begin{center}
  \copyright\ Copyright by \thesisauthor\ \thesisyear \\
  All Rights Reserved
\end{center}
\clearpage

% ===================================================================
% 4. ACKNOWLEDGEMENTS
% ===================================================================
\chapter*{Acknowledgements}
\addcontentsline{toc}{chapter}{Acknowledgements}

% SOURCE: thesis.tex lines 181-196 (ported verbatim, first-pass draft,
% not yet reviewed).

I would like to thank my advisor, Christopher A. Kitts, and the members of
the Santa Clara University Robotic Systems Laboratory for their support
throughout this work.

A portion of this work has been published or submitted in [1] and [2], and
portions are in preparation for future publication.

\clearpage

% ===================================================================
% 5. ABSTRACT
% ===================================================================
\chapter*{Abstract}
\addcontentsline{toc}{chapter}{Abstract}

\begin{center}
  \thesistitle \\[6pt]
  \thesisauthor \\[6pt]
  Department of \thesisdepartment \\
  Santa Clara University \\
  Santa Clara, California \\
  \thesisyear
\end{center}

\vspace{0.4in}

% SOURCE: thesis.tex lines 220-246, ported with fit-order framing added.
% First-pass draft, not yet reviewed.

\noindent
This dissertation develops a progression of tools for multirobot estimation
and navigation in physical fields, in which the order of a local polynomial
fit to cooperative field measurements sets the class of feature the team can
locate. A functional hardware testbed of omnidirectional rovers, printed
HSV-encoded vector field maps, and neural-network sensor calibration
establishes the experimental platform used throughout, together with a
zeroth-order (direction-only) baseline that motivates the need for
higher-order estimation. A first-order estimator is then shown to recover a
two-dimensional vector field's local Jacobian from three robots' simultaneous
local measurements alone, enabling closed-form critical point localization,
eigenvalue-based classification, and attraction and orbital control laws;
three robots is proven both necessary and sufficient for this estimate, and
the method is validated in simulation across eight canonical fields and on
hardware with a 100\% convergence success rate over 157 trials. This result
is then extended to second order: a six-robot team recovers a flow's local
Jacobian, Okubo-Weiss parameter, and Hessian from simultaneous measurements,
enabling two modified-Newton controllers that track dynamic boundaries and
separatrices no single critical point can describe, with a proven stability
guarantee and validation on both a canonical double-gyre flow and real
oceanographic current data. Together, these contributions establish a
coherent methodology for distributed, model-free field estimation that
scales from heading-only baselines through isolated critical points to
extended coherent structures using only local, cooperative sensing.

\clearpage

% ===================================================================
% 6. GLOSSARY OF TERMS
% ===================================================================
\chapter*{Glossary of Terms}
\addcontentsline{toc}{chapter}{Glossary of Terms}

% SOURCE: thesis.tex lines 269-295, ported verbatim.

\begin{table}[h]
\centering
\begin{tabular}{l p{3.5in}}
$\mathbf{p}$ & Position vector in the plane, $\mathbf{p}=(x,y)$ \\
$\mathbf{p}^*$ & Critical point of a vector field, $\mathbf{v}(\mathbf{p}^*)=\mathbf{0}$ \\
$\mathbf{v}(\mathbf{p})$ & Vector field value at $\mathbf{p}$, components $(u,v)$ \\
$\mathbf{J}$ & Field Jacobian, $\mathbf{J} = \left[\begin{smallmatrix} u_x & u_y \\ v_x & v_y \end{smallmatrix}\right]$ \\
$\mathbf{J}_c$ & Kinematic (cluster-space) Jacobian, robot-space to cluster-space velocity map \\
$\alpha_{\text{eig}}$ & Real part of a complex eigenvalue pair of $\mathbf{J}$ (spiral classification) \\
$\alpha_{\text{mom}}$ & Momentum / discretization coefficient, $\alpha_{\text{mom}} = \exp(-\Delta t/\tau)$ \\
$p, q, \beta$ & SAS triangle side lengths and interior angle (Kitts-lab notation, unchanged across chapters) \\
$\mathbf{q}$ (bold) & Cluster-space state vector (pose + shape, dimension $2N$) \\
$r$ & Orbital radial vector / magnitude, $\mathbf{r} = \mathbf{p}^*-\mathbf{p}_c$ \\
$\rho$ & Polar radius (field definitions, Appendix~A) and pentagon ring radius \\
$D$ & Okubo-Weiss parameter, $D = \det(\mathbf{J})$ \\
$\mathbf{H}_D$ & Hessian of $D$ \\
$\boldsymbol{\phi}(\tilde{\mathbf{p}})$ & Quadratic sensing basis vector \\
$\boldsymbol{\Phi}$ & Six-robot formation matrix (capital Phi) \\
$\tau$ & Actuator time constant, 0.28~s on Decabot hardware (after recalibration) \\
$\Delta t$ & Control period, 0.1~s (10~Hz) \\
$v_{\max}$ & Maximum robot speed, 0.3~m/s \\
\end{tabular}
\end{table}

\clearpage

% ===================================================================
% 7. TABLE OF CONTENTS
% ===================================================================
\tableofcontents
\clearpage

% ===================================================================
% 8. LIST OF FIGURES
% ===================================================================
\listoffigures
\addcontentsline{toc}{chapter}{List of Figures}
\clearpage

% ===================================================================
% 9. LIST OF TABLES
% ===================================================================
\listoftables
\addcontentsline{toc}{chapter}{List of Tables}
\clearpage

% ===================================================================
% BODY -- switch to arabic page numbering
% ===================================================================
\pagenumbering{arabic}
\setcounter{page}{1}
\doublespacing

% ===================================================================
% CHAPTER 1: INTRODUCTION
% ===================================================================
\chapter{Introduction}
\label{ch:intro}

\section{Adaptive Navigation in Vector Field Environments}
\label{sec:ch1:motivation}
% Ported from: thesis.tex Ch.1 Motivation (lines 332-346); extended with
% cwaight_systems_paper.tex Sec. I (convergence zones, leaks, eddies) and
% Draft_11.tex Sec. I (coherent structures, transport barriers).

Many environments of scientific and practical interest are organized by an
underlying vector field: an ocean current, an atmospheric flow, a diffusing
plume, a magnetic field. Locating and characterizing the special points and
structures of such a field, its sinks, sources, vortices, saddles, and the
separatrices that divide one flow regime from another, is often more useful
than simply following the field's local direction. A search-and-rescue team
drifting with an ocean current benefits from knowing where a nearby eddy
will trap debris; an environmental monitoring team benefits from knowing
where two water masses meet. In many of these settings no map of the field
exists in advance, and no single sensor can see the whole field at once. A
team of mobile robots, each sampling the field locally and sharing
measurements with its neighbors, can estimate field structure that no
individual robot could determine alone, and can then use that estimate to
steer.

The oceanic and atmospheric environments that motivate this work are
governed by a spatially varying velocity or force, and their transport and
accumulation behavior is organized by the critical points of the field, the
coordinates where the flow vanishes \cite{1,2}. A surface
convergence zone acts as a sink, gathering buoyant debris and drifting
bodies; this makes such zones high probability search regions after a loss
at sea \cite{3,4}. Containing a ruptured wellhead or a harbor
discharge requires reaching the source against its own outward flow
\cite{5,6}. A mesoscale ocean eddy serves as a center, trapping a
body of water, along with the plankton bloom inside it, and carrying both
across a basin for weeks \cite{7,8,9}. Monitoring the bloom
over its lifetime therefore requires a sensor to orbit the drifting core at
a fixed standoff; a ground-fixed station would not stay with the eddy.

Beyond critical points, vector fields have curves that are useful for
navigation. Dynamic ocean flows contain naturally occurring boundaries known
as coherent structures, the strain-dominated curves across which water
masses separate and along which floating material collects
\cite{10,11}. Mapping these curves lets a team deploy resources to a
targeted search area instead of sampling a much larger basin to find where
the structure lies \cite{11,12,13,14}.

The gap this dissertation addresses is locating and characterizing these
features, points and curves alike, from measurements a robot team takes of
the physical field itself, not from an artificial field constructed to
guide the team, and not from a heading alone. A heading toward, or along, a
feature is not the same as knowing where the feature is.

\section{Literature Review}
\label{sec:ch1:literature}

\subsection{Adaptive Navigation in Scalar Fields}
\label{sec:ch1:lit-scalar}
% Ported from: Draft_11.tex Sec. I (lines 94-106), extended with
% cwaight_systems_paper.tex Sec. I (line 51) for the strategy list and
% field-deployment citations.

Adaptive navigation (AN) in scalar fields is well developed. Each robot
contributes one scalar measurement to a policy that reaches and follows a
feature, and formations are sized so that differential measurements across
the cluster recover the gradient and curvature that policy consumes
\cite{15,16}. \"{O}gren, Fiorelli, and Leonard established cooperative
gradient climbing with a formation of mobile sensors \cite{16}.
Bri\~{n}\'{o}n-Arranz, Renzaglia, and Schenato extended it to the Hessian,
five robots on a ring plus one at the center recovering both in closed form
\cite{15}. Kitts, McDonald, and Neumann developed primitives for extrema
finding, contour following, ridge/trench following, and saddle point
station keeping \cite{17}, and verified ridge, trench, saddle, and front
tracking experimentally on an indoor testbed \cite{18,19}.

Additional strategies extend this line of work to gradient-based cluster
space navigation \cite{20}, extended information consensus filters
\cite{21,22}, constrained nonsmooth distributed optimization
\cite{23}, and bio-inspired search \cite{24}, and several such
systems have been field-deployed on ground, surface, and aerial vehicles for
environmental monitoring, pollution tracking, and search and rescue
\cite{20,14,12,25}.

\subsection{Artificial Vector Fields for Guidance}
\label{sec:ch1:lit-artificial}
% Ported from: cwaight_systems_paper.tex Sec. I (line 53), expanded with
% individual citations, including Liddy et al. from thesis.tex, naming the
% examples the source sentence only cited as a group.

Much of the broader vector field literature constructs artificial fields
for path-following and obstacle avoidance, a separate problem from sensing
a real flow, and field campaigns that track features carried by an actual
ocean current show the difference \cite{26}. Representative examples
include vector field path following for miniature air vehicles
\cite{27}, singularity-free guiding vector fields for robot navigation
\cite{28}, obstacle avoidance using complex vector fields
\cite{29}, motion planning and collision avoidance using navigation
vector fields \cite{30}, and adaptive vector field guidance without a
priori knowledge of course dynamics and wind \cite{31}. None of these
fields are sensed; they are specified in advance, and the robot's task is
to follow one, a different problem from the one this dissertation poses.

\subsection{Robotic Sensing of Physical Flows}
\label{sec:ch1:lit-sensing}
% Ported from: Draft_11.tex Sec. I (lines 116-124), extended with
% cwaight_systems_paper.tex Sec. I (lines 53, 57) for the two-camps framing
% and the Knizhnik citation.

Work that senses the physical field directly falls into two camps. One
camp presumes prior knowledge of the environment, a map, a model, or an
assumed functional form, unavailable in the unmapped settings this
dissertation targets \cite{32,33}. The other operates online using
only local measurements. Michini et al. bracket a presumed stable or
unstable manifold with three robots using only local velocity measurements
\cite{34}, an approach inspired by the PIM triple procedure, extended
to $N$ robots \cite{35} and validated experimentally in gyre-like flows
\cite{36}. Knizhnik et al. address a related setting, flow-based
control of marine robots in gyre-like environments \cite{37}. In the
Michini line of work the robots begin on the manifold and are advected by
the flow, and the manifold's identity is supplied in advance; the saddle
point itself is treated by neither the controller nor its analysis, and
\cite{34} reports veering away from the structure near a saddle in
flow-tank experiments. Kularatne and Hsieh confirm the boundary by
computing the local finite-time Lyapunov exponent field on board, at the
cost of a nonzero integration horizon and a commitment to attracting
structures \cite{1}. In every case the robotic formation pre-straddles
the structure and does not acquire a structure it does not already occupy.

\section{Problem Statement}
\label{sec:ch1:problem}
% Ported from: thesis.tex Ch.1 Problem Statement (lines 365-388).

Existing approaches to multirobot navigation in scalar and vector fields
generally fall into two categories. The first requires prior knowledge of
the field, a map, a model, or an assumed functional form, against which a
robot team's local measurements are compared or fit. The second uses only
local information but provides directional guidance alone: climb the
gradient, follow the current, avoid the strain. Neither category directly
answers the question this dissertation poses: given only local, real-time
measurements and no prior model, can a robot team estimate the location and
type of a field's critical points, and, more generally, the field's local
first- and second-order structure, well enough to navigate with respect to
it, not just along it.

This dissertation answers that question constructively. It shows that
three robots are both necessary and sufficient to recover a
two-dimensional vector field's local linearization (its Jacobian) from
simultaneous local measurements alone, and that six robots are both
necessary and (almost always) sufficient to recover the same field's local
quadratic structure (its Jacobian, Okubo-Weiss parameter, and Hessian). In
both cases the estimate is exact for the class of fields the robots'
formation size implies (affine for three robots, quadratic for six),
assembled without any assumption about the field's global shape, and used
directly inside a navigation control law whose stability is established
analytically and validated in simulation and, where hardware permits, on a
physical testbed.

\section{Thesis Statement: Fit Order Sets the Reachable Feature Class}
\label{sec:ch1:thesis}
% New text. States the fit-order ladder, cites cwaight_systems_paper.tex
% Sec. VI-B (vector-sum/vector-to-scalar drift) and Draft_11.tex Sec. I/VII
% (linear-fit-at-one-point limitation, cubic future work) as the motivating
% failures each rung fixes.

This dissertation is organized around one idea: the order of the local
polynomial a robot team fits to a sensed vector field sets which class of
field feature that team can locate. A zeroth-order fit returns a single
vector, the field's local mean. A first-order (affine) fit returns the
field's local Jacobian. A second-order (quadratic) fit returns the field's
local Hessian in addition to the Jacobian. A third-order (cubic) fit would
return the field's local third derivatives. Each order needs more robots,
sampled at one instant, to determine without assuming information the
field has not supplied, and each order unlocks a feature class the order
below cannot represent. Table~\ref{tab:ch1:fit-order-ladder} states the
ladder used throughout this dissertation.

\begin{table}[h]
\centering
\caption{The fit-order ladder used throughout this dissertation.}
\label{tab:ch1:fit-order-ladder}
\small
\begin{tabular}{@{}p{1.3in}cp{1.1in}p{1.7in}c@{}}
\toprule
Fit order & Robots per component & Degeneracy & Features reachable & Chapter \\
\midrule
Zeroth (mean vector, magnitude) & 1+ & none & heading only; drifts on orbits & \ref{ch:zeroth} \\
First (affine, Jacobian) & 3 & collinear robots & critical points: location, type, orbit & \ref{ch:first_order} \\
Second (quadratic) & 6 & robots on a common conic & separatrices, Okubo-Weiss boundaries, saddles as minima & \ref{ch:second_order} \\
Third (cubic, future) & 10 & robots on a common cubic & true Hessian of $D$, transverse curvature & \ref{ch:limitations} \\
\bottomrule
\end{tabular}
\end{table}

Chapter~\ref{ch:zeroth} shows why the ladder starts here. A zeroth-order
primitive steers on the mean of the perceived flow direction across the
formation (the vector-sum primitive) or on the flow's magnitude alone (the
vector-to-scalar primitive \cite{38}, here modified to move tangentially
along the direction of steepest descent). Both supply a heading, not a
location. Orbiting a vortex under either primitive drifts outward in
spiral fashion, because neither has a mechanism to detect or correct
radial error with respect to a specific point, and the vector-sum drift has
been confirmed in hardware \cite{39}. On a saddle the failure is worse:
the field vectors do not all point toward the critical point, so the
vector-sum controller cannot produce an orbital trajectory at all.

A first-order fit removes this failure. Recovering the field's Jacobian at
a single point gives both the location and the type of an isolated critical
point, and supports a control law that converges onto it or holds a
commanded orbit around it (Chapter~\ref{ch:first_order}). But a linear fit
is a property of one point \cite{40}. It cannot by itself locate an
extended curve, such as the separatrix or Okubo-Weiss boundary that
organizes transport in a region no single critical point describes, because
a one-point fit has no gradient of its own to climb. A second-order fit
gives the velocity gradient as a field over the formation rather than a
value at a point, so a quantity built from it, the Jacobian determinant or
the smaller rate-of-strain eigenvalue, comes with its own gradient in
closed form and can be climbed to and held on a curve
(Chapter~\ref{ch:second_order}).

The next rung is not built in this dissertation. Chapter~\ref{ch:limitations}
describes a cubic fit from ten robots as future work: it would recover the
field's third derivatives directly, giving the true Hessian of the
determinant $D$ and the transverse curvature of a rate-of-strain trench,
quantities the second-order fit can only approximate.

\section{Contributions}
\label{sec:ch1:contributions}
% Ported from: thesis.tex Ch.1 Contributions (lines 389-424), rewritten in
% fit-order terms. The old list's second item, a compositional grammar for
% general cluster-space kinematics, is not listed here; it is out of this
% dissertation's outline, which derives only the three-robot and six-robot
% kinematics used by the first- and second-order chapters (Appendices
% \ref{app:kin3} and \ref{app:kin6}).

The main contributions of this dissertation are:
\begin{enumerate}
  \item A hardware and simulation toolkit for multirobot vector-field
  navigation: HSV-encoded printed fields and neural-network sensor
  calibration and recalibration (R\textsuperscript{2}~=~0.96 direction,
  0.91 magnitude) \cite{39}, robot dynamics identification, and a
  matching simulator (Chapter~\ref{ch:tools}). The Decabot platform itself
  is prior work of other members of the Robotic Systems Laboratory; the
  contribution here is the sensing, calibration, and simulation layer built
  on it.
  \item Zeroth-order baseline primitives, vector-sum and vector-to-scalar
  direction-only navigation from prior work \cite{38}, implemented and
  evaluated on the hardware testbed of Chapter~\ref{ch:tools}
  \cite{39}, establishing what direction
  alone can and cannot certify (Chapter~\ref{ch:zeroth}).
  \item A distributed, model-free method for three robots to estimate a
  two-dimensional vector field's local Jacobian from simultaneous local
  measurements alone, with no prior field knowledge, proven minimal at
  exactly three robots, enabling closed-form critical point localization,
  eigenvalue-based type classification, and attraction and orbital control
  laws; validated in simulation (eight fields, 1000 Monte Carlo trials each,
  100\% convergence) and on hardware (157 convergence trials plus 12
  orbital trials, 100\% success, sub-centimeter bias)
  (Chapter~\ref{ch:first_order}).
  \item A distributed, model-free method for six robots (pentagon plus
  center) to estimate a two-dimensional flow's local Jacobian, Okubo-Weiss
  parameter, and Hessian from simultaneous local measurements, proven
  minimal at exactly six robots and almost-always sufficient, ruling out
  the degenerate hexagon and conic trap, enabling two modified-Newton
  controllers, the $D$ tracker and the $s_1$ tracker, that track the
  Okubo-Weiss boundary and the separatrix trench respectively, with a
  proven stability and traversal guarantee and validation on both a
  canonical double-gyre model and real ocean high-frequency-radar current
  data (Chapter~\ref{ch:second_order}).
  \item Findings that hold across fit orders: how estimation noise
  amplifies with fit order, how formation size trades truncation error
  against noise sensitivity, and how objectivity under a rotating observer
  depends on which surrogate field a controller climbs
  (Chapter~\ref{ch:crosscutting}).
\end{enumerate}

\section{Publications}
\label{sec:ch1:publications}
% Ported from: thesis.tex Ch.1 Publications (lines 426-453), restructured
% to the four publications behind this dissertation's fit-order chapters.
% Titles taken verbatim from each source's own title. Statuses left as
% [TBD] where no source states one; per CLAUDE.md, this dissertation's own
% submission is intentionally held until the critical-points paper
% publishes, and that is not a blocker to flag.

Portions of this dissertation have appeared, or are in preparation to
appear, in the following publications.

\begin{table}[h]
\centering
\caption{Mapping of dissertation chapters to publications.}
\label{tab:ch1:publications}
\small
\begin{tabular}{@{}cp{1.9in}p{1.0in}c@{}}
\toprule
Chapter & Working title & Venue & Status \\
\midrule
\ref{ch:tools}, \ref{ch:zeroth} & A Functional Indoor Testbed for Multirobot Adaptive
Navigation in Vector Field Environments \cite{39} & ASME IDETC/CIE
2025, Anaheim (DETC2025-167604) & Published \\
\ref{ch:first_order} & Adaptive Navigation of Multirobot Systems to
Critical Points in 2D Vector Fields in Simulation and Experiment & IEEE
Systems Journal & \textbf{[TBD]} \\
\ref{ch:second_order} & Second-Order Cooperative Field Estimation for
Multirobot Tracking of Coherent Flow Structures & IEEE Systems Journal &
In preparation \\
\bottomrule
\end{tabular}
\end{table}
%% IDETC 2025 testbed paper: published (author confirmed 2026-09-29; final
%% PDF in Paper_Writing/Reference Papers/[33]).
%% The 2024 IDETC submission (DETC2024-142746) was rejected; its early
%% ideas are not part of this dissertation and it is not cited.
%% TODO(status): confirm the critical-points paper's submission status.

\section{Dissertation Organization}
\label{sec:ch1:organization}
% Ported from: thesis.tex Ch.1 Organization (lines 454-474), rewritten for
% the ten-chapter, six-appendix fit-order structure.

Chapter~\ref{ch:tools} describes the verification tools this dissertation
depends on: the Decabot hardware platform (prior work of the Robotic
Systems Laboratory), the HSV-encoded printed vector fields and
neural-network sensor calibration developed for it, which converts a
camera reading to a velocity estimate \cite{39}, robot dynamics
identification, and the matching simulation environment used throughout
the rest of this dissertation.

Chapter~\ref{ch:preliminaries} sets out the mathematics common to every fit
order: the planar vector field, the velocity gradient tensor, critical
point classification by eigenvalue structure, the rotation-strain
decomposition and the Okubo-Weiss partition, coherent structures,
objectivity under a change of observer frame, and the canonical benchmark
fields used for validation.

Chapter~\ref{ch:estimation} presents the cooperative estimation framework
shared by every fit order above zeroth: the multilayer control
architecture, cluster space control for the three-robot and six-robot
formations, local polynomial field estimation and its minimality,
estimator sensitivity to noise and truncation, and the fit-order ladder of
Section~\ref{sec:ch1:thesis} restated with this framework's machinery.

Chapter~\ref{ch:zeroth} presents the zeroth-order baseline: direction-only
navigation and the vector-sum and vector-to-scalar primitives \cite{38},
hardware trials on
fixed and sinking vortices on the testbed of Chapter~\ref{ch:tools}
\cite{39}, and the argument that direction alone is not location.

Chapter~\ref{ch:first_order} presents the first-order result: three robots
recovering a vector field's local Jacobian, and the resulting critical
point estimation, classification, attraction, and orbital control laws,
validated in simulation and on hardware and compared against the
zeroth-order primitives of Chapter~\ref{ch:zeroth}.

Chapter~\ref{ch:second_order} presents the second-order result: six robots
recovering a flow's local Hessian and Okubo-Weiss structure, the $D$
tracker and $s_1$ tracker built on that estimate, their stability and
traversal properties, and their validation on a double-gyre benchmark and
on Santa Barbara Channel radar current data.

Chapter~\ref{ch:crosscutting} collects findings that appear only once
results from more than one fit order are set side by side: how noise
amplifies with fit order, how formation size trades truncation error
against noise sensitivity, formation conditioning versus field
conditioning, instantaneous estimation on a time-varying field, and
objectivity and sign references across controllers.

Chapter~\ref{ch:limitations} states this dissertation's limitations by fit
order and describes future work, including the cubic, ten-robot extension
that would complete the next rung of the ladder.

Chapter~\ref{ch:conclusion} summarizes the contributions and closes with
final thoughts.

Appendix~\ref{app:fields} collects the canonical field definitions used in
simulation throughout this dissertation. Appendix~\ref{app:kin3} derives
the three-robot cluster's kinematics in full. Appendix~\ref{app:kin6}
derives the six-robot pentagon-plus-center cluster's kinematics in full.
Appendix~\ref{app:frame_equivariance} proves the frame equivariance of the
$s_1$ tracker. Appendix~\ref{app:stability} collects the transverse
stability proofs for the controllers of Chapters~\ref{ch:first_order}
and~\ref{ch:second_order}. Appendix~\ref{app:stats} documents the
statistical methods used to report simulation and hardware results.

\section{Notation}
\label{sec:ch1:notation}
% Ported from: thesis.tex Ch.1 "Notation and Definitions" (lines 476-586),
% chapter cross-references updated to the ten-chapter structure and
% trimmed of notation tied to material not in this dissertation's outline
% (the consensus orientation gauge and the composition-tree spine angle,
% both part of the compositional-grammar work excluded per Section
% \ref{sec:ch1:contributions}). The full symbol reference is the Glossary
% of Terms in the front matter.

Throughout this dissertation, bold lowercase letters denote vectors, e.g.,
$\mathbf{p} \in \mathbb{R}^2$, and bold uppercase letters denote matrices,
e.g., $\mathbf{J} \in \mathbb{R}^{m \times n}$. The studies collected in
this dissertation were originally written as independent papers, and their
notations were never cross-checked against one another. Most symbols
already agree; a small number collide and are resolved here, once, for the
whole document. Chapter-specific notation not covered here is introduced
where it is first used.

\subsection*{The bold-\texorpdfstring{$\mathbf{p}$}{p} versus scalar-\texorpdfstring{$p$}{p} rule}

Bold $\mathbf{p}$ always denotes a two-dimensional position,
$\mathbf{p} = (x,y)$, and is never written unbolded when used as a
position. Scalar $p$ always denotes the side-length formation parameter
inherited from the Kitts-lab cluster-space lineage, introduced in
Chapter~\ref{ch:estimation} and used through Appendices~\ref{app:kin3} and
\ref{app:kin6}: the distance from robot 1 to robot 2 in a side-angle-side
(SAS) triangular formation. There is no unbolded $p$ for position and no
bold $\mathbf{p}$ for side length. This rule is carried unchanged from the
source papers and holds throughout every chapter of this dissertation.

\subsection*{The bold-\texorpdfstring{$\mathbf{q}$}{q} versus scalar-\texorpdfstring{$q$}{q} rule}

Chapter~\ref{ch:estimation} introduces bold $\mathbf{q}$ as the general
cluster-space state vector (pose and shape stacked, dimension $2N$ for $N$
robots), while scalar $q$, following the same inherited SAS convention as
scalar $p$, denotes the triangle's second side length (the 2--3 distance).
These two uses never appear in the same equation and are distinguished
exactly as bold $\mathbf{p}$ and scalar $p$ are: there is no unbolded
$\mathbf{q}$ for the state vector and no bold $q$ for the side length.

\subsection*{The alpha collision}

The symbol $\alpha$ is used for three unrelated quantities across the
source materials. Each is disambiguated with a subscript in this
dissertation and must never be written as bare $\alpha$ after this point.

\begin{itemize}
\item $\alpha_{\text{eig}}$: the real part of a complex eigenvalue pair of
  the field Jacobian $\mathbf{J}$, written $\alpha_{\text{eig}} \pm
  i\omega$. This is a property of the vector field's linearization at a
  critical point and classifies spiral-type critical points
  (Chapter~\ref{ch:first_order}).
\item $\alpha_{\text{mom}}$: the discrete first-order momentum filter
  coefficient for robot dynamics, $\alpha_{\text{mom}} = \exp(-\Delta
  t/\tau)$. At 10~Hz with $\tau \approx 0.28$~s, $\alpha_{\text{mom}} \approx
  0.70$. This is characterized in Chapter~\ref{ch:tools} and appears
  identically in the robot dynamics of Chapters~\ref{ch:first_order} and
  \ref{ch:second_order}.
\item $\phi_1, \phi_3$: the two SAS-triangle vertex angles at robots 1 and
  3, used only in the triangle's law-of-cosines derivation
  (Appendix~\ref{app:kin3}). These were originally written as $\alpha$ and
  $\gamma$ in one source paper's appendix; they are renamed here since they
  are otherwise unused and the original choice collided with both
  $\alpha_{\text{eig}}$ and $\alpha_{\text{mom}}$.
\end{itemize}

\subsection*{The \texorpdfstring{$r$}{r} collision}

Three meanings of $r$ appear across the source materials: the orbital
radial vector/magnitude $\mathbf{r} = \mathbf{p}^* - \mathbf{p}_c$
(Chapter~\ref{ch:first_order}), the SAS triangle's third side length (the
1--3 diagonal), and the polar radius used in several canonical field
definitions. This dissertation keeps $r$ for the orbital radial quantity,
since it appears in a named control law. The triangle's third side is
renamed $s$, and the polar radius in field definitions is renamed $\rho$,
matching the pentagon ring radius already used in
Chapter~\ref{ch:second_order} and Appendix~\ref{app:kin6}, so that
``radius of a circular construction'' reads consistently throughout.

\subsection*{The \texorpdfstring{$\mathbf{J}$}{J} collision}

Bold $\mathbf{J}$ denotes the field Jacobian, $\mathbf{J} =
\left[\begin{smallmatrix} u_x & u_y \\ v_x & v_y \end{smallmatrix}\right]$,
throughout this dissertation: it is defined in
Chapter~\ref{ch:preliminaries}, recovered from measurements in
Chapter~\ref{ch:estimation}, and is the central object of
Chapter~\ref{ch:first_order} (first order) and Chapter~\ref{ch:second_order}
(embedded in the second-order estimator). Every kinematic or cluster-space
Jacobian, mapping shape-parameter velocities to individual robot
velocities, is written $\mathbf{J}_c$ instead, whether it is the $6\times
6$ triangle Jacobian or the $12\times 12$ pentagon Jacobian, both presented
in Chapter~\ref{ch:estimation} and derived in full in
Appendices~\ref{app:kin3} and~\ref{app:kin6}.

\subsection*{Other symbols}

The full symbol reference for field and estimation notation, critical
point classification, control laws, robot dynamics, and formation
parameters is the Glossary of Terms in the front matter, extended with the
Chapter~\ref{ch:estimation} and Chapter~\ref{ch:second_order} additions
listed there. Of particular note: capital $\boldsymbol{\Phi}$ (the
six-robot formation matrix, Chapter~\ref{ch:second_order}) is visually
close to lowercase $\phi_1, \phi_3$ (the triangle vertex angles,
Appendix~\ref{app:kin3}) in some fonts; the two never appear in the same
chapter, but the reader should not assume a relationship between them.
Similarly, $\rho$ appears in two guises across this dissertation, the field
definitions' polar radius (Chapter~\ref{ch:preliminaries},
Appendix~\ref{app:fields}) and the pentagon formation's ring radius
(Chapter~\ref{ch:second_order}, Appendix~\ref{app:kin6}).

As in Chapter~\ref{ch:first_order}, this dissertation writes a critical
point of a flow, including a saddle point of the Okubo-Weiss field in
Chapter~\ref{ch:second_order}, as $\mathbf{p}^*$.

% ===================================================================
% CHAPTER 2: VERIFICATION TOOLS
% ===================================================================
\chapter{Verification Tools}
\label{ch:tools}

This chapter presents the hardware and simulation tools used to verify
the estimation and control results of the rest of this dissertation. It
separates what these tools inherit from prior lab work from what this
dissertation adds to them; Section~\ref{sec:ch2:summary} states that
division plainly.

\section{Decabots (Prior Work)}
% Ported from: thesis.tex Ch.2 Introduction (lines 593-622); Systems
% paper Sec. V-A hardware description; bib_sources.tex sys:52, sys:47.
\label{sec:ch2:decabots}

Every hardware result in this dissertation runs on the Decabot testbed,
a platform built by prior members of the Santa Clara University Robotic
Systems Laboratory and reused here without modification to the robots
or the motion capture system. The Decabot is a low-cost omnidirectional
rover developed for hardware verification of cluster-space control
schemes \cite{41}. A prior generation of this testbed encoded
vector field readings directly in RGB color and printed the result on
the floor as a map for the robots to sense, following an indoor testbed
feasibility study for vector-field adaptive navigation \cite{42}.
That RGB approach suffered a persistent limitation: interference between
color channels in the robots' undercarriage-mounted sensors degraded the
system's ability to recover accurate direction and magnitude
information, and the discrete color palette limited the effective
sampling rate to roughly 0.15~Hz, too slow to support the continuous,
closed-loop control laws developed in
Chapters~\ref{ch:first_order} and~\ref{ch:second_order}.

The remainder of this chapter describes what this dissertation adds on
top of this inherited platform: an HSV-based field encoding, a
neural-network sensor calibration that corrects for the channel
interference the RGB approach could not, a characterization of robot
dynamics and velocity response, and the simulation environment used
throughout the rest of this dissertation.

\section{Decabot Testbed and Motion Capture}
% Ported from: thesis.tex Ch.2 "Testbed Overview" (lines 662-671) and
% Control Architecture Overview (lines 645-661); Systems paper Sec. V-A
% for the workspace footprint used in current hardware trials; DETC2025-167604
% (thesis:33) Sec. 3.1 "Testbed Overview" and Sec. 5 "System Limitations"
% for the floor-map area and physical constraint figures.
\label{sec:ch2:testbed}

The physical testbed consists of three omnidirectional Decabot rovers
operating on 2.3~m\textsuperscript{2} multicolored printed floor maps,
an OptiTrack motion capture system providing centimeter-level position
accuracy via reflective markers on each robot's roof, and a central
computer running the control architecture in MATLAB/Simulink
\cite{39}. Robots communicate with the central computer over
WiFi at 10~Hz: at each control cycle the central computer receives
OptiTrack positions and sensor readings, computes the control laws, and
returns velocity commands, which each robot's onboard controller
converts into drive motor commands. Physical constraints include a
maximum 3~m~$\times$~6~m work area and approximately one hour of
battery life per robot \cite{39}; thesis.tex additionally reports
a lower bound of roughly 1~m~$\times$~1~m, a collision-clearance minimum
for a multirobot team. The three-robot hardware trials reported in
Chapter~\ref{ch:first_order} used a 1.6~m~$\times$~1.6~m printed map
within this range.

Each Decabot carries a downward-facing TCS34725 RGB color sensor,
enclosed in a light-shrouded housing to minimize external lighting
effects, whose channels are not fully isolated from one another per its
datasheet, the source of the interference problem
Section~\ref{sec:ch2:decabots} describes. Correcting for this
interference is the subject of Section~\ref{sec:ch2:nn_calibration}.

The robot controller layer, the cluster-space controller layer, and the
adaptive navigation layer that this testbed runs are described in full
in Chapter~\ref{ch:estimation}; this chapter is concerned only with the
physical and sensing tools those layers run on.

\section{Vector Field Encoding in HSV Color Space}
\label{sec:ch2:hsv}

\subsection{Limitations of RGB Encoding}
% Ported from: thesis.tex Ch.2 Introduction (lines 599-607); DETC2025-167604
% (thesis:33) Sec. 1 Introduction and Sec. 3.2 for detail thesis.tex dropped.

Encoding a vector field reading directly in RGB color, as the prior
testbed generation did, ties the sensor's three color channels directly
to the field's own two independent components (direction and
magnitude), with no channel held in reserve to detect or correct for
cross-talk between them. Approaches using red and green to represent
orthogonal vector components suffered from sensor channel interference
on the robots' RGB sensors, leaving only a small portion of the color
spectrum usable once high-interference regions were eliminated
\cite{39}. This interference made accurate decoding of direction
and magnitude information difficult, and was particularly severe for
vectors of small magnitude, where minor sensor value changes could
produce large apparent direction shifts \cite{39}. The TCS34725's
red, green, and blue channels are not fully isolated per its own
specification sheet \cite{39}, and this effect persisted even
after conversion to an HSV color space; traditional calibration methods
neither significantly improved performance nor supported effective
scaling across multiple robots \cite{39}, motivating the
supervised-learning calibration of Section~\ref{sec:ch2:nn_calibration}.
The resulting discrete color palette also limited sampling to roughly
0.15~Hz, too slow for continuous robot motion \cite{39}.

\subsection{Hue as Direction, Saturation as Magnitude}
% Ported from: thesis.tex Ch.2 "Modeling Vector Fields in Color Space"
% (lines 665-671); DETC2025-167604 (thesis:33) Sec. 3.2 for the
% value-versus-saturation reasoning thesis.tex dropped.
\label{sec:ch2:hue_sat}

A vector field reading has two components at each point: a direction
and a magnitude. This dissertation's testbed encodes direction in hue
and magnitude in saturation, allowing a two-dimensional vector field to
be represented as a printed image and sensed by a simple color sensor
rather than requiring active field generation. Hue runs from 0 to
$2\pi$ and represents flow direction; hue's circular topology is a
natural match for vector direction \cite{39}. Saturation, rather
than value, was chosen to encode magnitude, running from 0 to 1, after
early tests showed the value channel to be highly sensitive to printer
ink variation, ambient lighting, and shadow effects \cite{39}.
This separates the field's two physical components onto two different,
independently interpretable channels of the color space, rather than
overloading three RGB channels with two quantities.

\subsection{Printed Field Construction}
% Ported from: DETC2025-167604 (thesis:33) Sec. 3.2-3.3 for the field
% construction and scaling detail; Systems paper Sec. V-A (lines 328-345)
% for the current-work restatement of the same scaling.

This methodology was used to construct two vector fields for the
first-iteration testbed trials of \cite{39}: a fixed vortex and,
by adding a sinking component to the same vortex, a sinking vortex, both
centered on the printed map. Once computed on the print grid, vector
magnitudes were scaled to the 0.3--1.0 saturation range: values below
0.3 do not produce a commanded velocity strong enough to overcome
stiction (Section~\ref{sec:ch2:sat_velocity}), and values below 0.2 were
excluded from calibration training entirely
(Section~\ref{sec:ch2:nn_calibration}) \cite{39}. In the current
critical-points work, saturation is remapped to 0--1 in software, with
saturation 1.0 corresponding to the maximum robot velocity of 0.3~m/s
(the same cap given in Section~\ref{sec:ch2:dynamics}). Together with
the sensor shroud and the sensor's own illumination, this encoding is
minimally affected by overhead lighting variations; a test of turning
the overhead lab lights on and off confirmed no effect on readings from
the calibration palette \cite{39}. The robots' RGB sensors are
calibrated against a printed palette spanning hue 0 to $2\pi$ and
saturation 0 to 1 at 24 intervals in each dimension; each robot records
raw RGB and clear-channel counts over every palette cell. This printed-map
method has imperfections in both printing and sensing, addressed by the
calibration procedure of Section~\ref{sec:ch2:nn_calibration} and
revisited as a simulation noise source in
Section~\ref{sec:ch2:noise_model}.

\begin{figure}[h]
  \centering
  \includegraphics[width=0.75\textwidth]{ch2_measurement_error_comparison.png}
  \caption{Analytical fields, sensor-based reconstructions, and
  reconstruction error for vortex and saddle fields.}
  \label{fig:ch2:hsv_field}
\end{figure}
% FIGURE SOURCE: Paper_Writing/Systems Paper/cwaight_systems_paper/measurement_error_comparison.png

\section{Neural Network Sensor Calibration}
\label{sec:ch2:nn_calibration}
% Ported from: DETC2025-167604, "A Functional Indoor Testbed for
% Multirobot Adaptive Navigation in Vector Field Environments" (Waight
% and Kitts, 2025), cited throughout as thesis:33, Sec. 3.3-3.4 and
% Table 1, for the first-iteration methodology and numbers, verified
% directly against the extracted text at
% PhD Thesis/sources/idetc2025_testbed.txt; Systems paper Sec. V-A
% (lines 328-345) for the current, recalibrated numbers.
%
% Resolved: an earlier pass cited PhD Thesis/sources/idetc_paper.txt
% (DETC2024-142746, a control-primitives paper with no HSV/NN content)
% for this section by mistake; thesis:33 is the correct source and all
% numbers below are now checked directly against it.

\subsection{Calibration Palette and Data Collection}
% Ported from: thesis:33 Sec. 3.3.

The undercarriage color sensor's raw channels, red, green, blue, and a
clear channel, do not map cleanly to hue and saturation because the
channels are not fully isolated from one another. A calibration palette
with hues spanning 0 to $2\pi$ and saturations spanning 0 to 1, spaced at
24 equal intervals in each dimension, was printed, and three robots were
positioned above each palette cell to record RGB and clear-channel (K)
measurements alongside the intended hue and saturation values
\cite{39}. The three-robot dataset was combined and augmented by
adding $\pm 1\%$ noise to the RGB readings for model robustness. Raw
RGBK values were normalized to $[0,1]$, converted to hue, saturation,
and value using standard libraries, and appended to the original data,
giving a seven-dimensional input space $(R,G,B,K,H,S,V)$. Data points
with saturation below 0.2 were removed to improve model performance
(the same threshold noted in Section~\ref{sec:ch2:hsv}), and all inputs
and targets were normalized to $[-1,1]$ before training
\cite{39}.

\subsection{Network Architecture and Training}
% Ported from: thesis:33 Sec. 3.3.

This dissertation's solution is a pair of neural networks, trained per
robot (or, for a generalization test, trained on three robots and
validated on a fourth held out), that map the sensor's raw
seven-dimensional input $(R,G,B,K,H,S,V)$ to corrected hue and saturation
estimates. Both networks use two fully connected hidden layers with Tanh
activations; the saturation network has a single output neuron for
direct magnitude prediction, while the hue network has two output
neurons, sine and cosine of hue, to respect hue's circular topology
\cite{39}. A random search determined hidden layer sizes (2--12
units in the first layer, 4--10 in the second); patience and
augmentation parameters were also randomized, while the
Levenberg-Marquardt optimizer and default learning rate were held fixed.
Data was split 70/15/15 into training, validation, and test sets
\cite{39}. Candidate models were evaluated on eight validation
sets (the calibration and verification plots of all four robots,
Section~\ref{sec:ch2:generalization}); selection used a min-max
criterion, identifying each candidate's lowest $R^2$ score across all
eight sets and keeping the candidate with the highest such minimum, for
each network independently \cite{39}. The current calibration in
Chapter~\ref{ch:first_order}'s critical-points work uses a simplified
pair of feedforward networks, each with two hyperbolic-tangent hidden
layers and a four-dimensional $(R,G,B,K)$ input, with the saturation
network regressing a single output and the hue network outputting sine
and cosine of hue as above.

\subsection{Cross-Robot Generalization}
\label{sec:ch2:generalization}
% Ported from: thesis:33 Sec. 4.1 and Table 1, checked directly against
% the extracted paper text (PhD Thesis/sources/idetc2025_testbed.txt,
% lines 328-416). Replaces an earlier, coarser two-table summary drawn
% from thesis.tex that rounded off the per-dataset spread; the values
% below are the paper's own per-robot, per-dataset breakdown.

Summarizing across all validation sets, the cross-trained model achieves
an overall $R^2$ of 0.96 in direction (hue) and 0.91 in magnitude
(saturation) \cite{39}. The full per-robot, per-dataset breakdown
is given in Table~\ref{tab:ch2:calibration}: robots 1 through 3 supply
both a calibration (Cal) and verification (Ver) validation set each,
used in training the cross-robot model, while robot 4 (R4) is held out
of training entirely and used only to test generalization. Across the
six validation sets from robots 1--3, the cross-trained model's hue
$R^2$ ranges from 0.966 to 0.998 and saturation $R^2$ from 0.935 to
0.998. On the held-out robot 4, the cross-trained model achieves hue
$R^2 = 0.993$ (RMSE 0.024) and saturation $R^2 = 0.908$ (RMSE 0.073) on
its calibration set, and hue $R^2 = 0.964$ (RMSE 0.054) and saturation
$R^2 = 0.925$ (RMSE 0.079) on its verification set \cite{39}.
Individually tuned per-robot models improve further, typically to hue
$R^2 > 0.99$ and saturation $R^2 > 0.95$, at the cost of several hours
of per-robot calibration that the cross-trained model avoids
\cite{39}; individual models are not available for the held-out
robot 4 by construction.

\begin{table}[h]
\centering
\caption{First-iteration calibration model performance, cross-trained
versus individually tuned, on each robot's calibration (Cal) and
verification (Ver) validation set. Robot 4 was held out of training
entirely.}
\label{tab:ch2:calibration}
\begin{tabular}{llcccc}
\toprule
Robot-set & Model & Hue $R^2$ & Hue RMSE & Sat.\ $R^2$ & Sat.\ RMSE \\
\midrule
R1-Cal & Cross & 0.966 & 0.017 & 0.950 & 0.054 \\
       & Indiv & 0.999 & 0.009 & 0.972 & 0.040 \\
R1-Ver & Cross & 0.969 & 0.050 & 0.950 & 0.064 \\
       & Indiv & 0.995 & 0.020 & 0.963 & 0.055 \\
R2-Cal & Cross & 0.998 & 0.013 & 0.935 & 0.061 \\
       & Indiv & 0.996 & 0.018 & 0.968 & 0.043 \\
R2-Ver & Cross & 0.967 & 0.052 & 0.943 & 0.069 \\
       & Indiv & 0.991 & 0.027 & 0.952 & 0.063 \\
R3-Cal & Cross & 0.998 & 0.014 & 0.948 & 0.055 \\
       & Indiv & 0.999 & 0.007 & 0.980 & 0.034 \\
R3-Ver & Cross & 0.991 & 0.027 & 0.953 & 0.062 \\
       & Indiv & 0.997 & 0.015 & 0.976 & 0.044 \\
R4-Cal & Cross & 0.993 & 0.024 & 0.908 & 0.073 \\
       & Indiv & -- & -- & -- & -- \\
R4-Ver & Cross & 0.964 & 0.054 & 0.925 & 0.079 \\
       & Indiv & -- & -- & -- & -- \\
\bottomrule
\end{tabular}
\end{table}

\subsection{Recalibration After Maintenance}
% Ported from: Systems paper Sec. V-A (lines 328-345), the current,
% authoritative recalibration.

All robots were recalibrated after receiving preventive maintenance and
servicing for the critical-points work of Chapter~\ref{ch:first_order}.
This recalibration is the authoritative, current numbers referenced
throughout the rest of this dissertation; the numbers of
Table~\ref{tab:ch2:calibration} above are the first-iteration figures
they superseded. Using two feedforward networks trained on the
four-dimensional $(R,G,B,K)$ feature vector, with training pooled across
robots so a single model transfers to an uncalibrated unit without
per-robot tuning, the current calibration achieves hue RMSE of 0.17~rad
(9.7\textdegree, $R^2 = 0.991$) and saturation RMSE of 4.4\%
($R^2 = 0.951$) across the operational range
(Table~\ref{tab:ch2:recalibration}).

\begin{table}[h]
\centering
\caption{Sensor calibration accuracy, first iteration versus current
recalibration.}
\label{tab:ch2:recalibration}
\begin{tabular}{lcc}
\toprule
 & Hue & Saturation \\
\midrule
First iteration, cross-robot model overall & $R^2 = 0.96$ & $R^2 = 0.91$ \\
Current recalibration & RMSE 0.17 rad (9.7\textdegree), $R^2 = 0.991$ & RMSE 4.4\%, $R^2 = 0.951$ \\
\bottomrule
\end{tabular}
\end{table}

\begin{figure}[h]
  \centering
  \includegraphics[width=\textwidth]{ch2_color_sensor_hsv_predictions.png}
  \caption{Neural network predictions for HSV field encoding on
  unit-normalized axes. Hue RMSE = 0.027 (0.17~rad, 9.7\textdegree),
  saturation RMSE = 0.044 (4.4\%).}
  \label{fig:ch2:hue_predictions}
\end{figure}
% FIGURE SOURCE: Paper_Writing/Systems Paper/cwaight_systems_paper/color_sensor_hsv_predictions.png

\section{Saturation-to-Velocity Characterization}
% Ported from: thesis.tex Ch.2 "Robot Velocity from Saturation"
% (lines 729-736); DETC2025-167604 (thesis:33) Sec. 3.4 "Robot Velocity
% from Saturation Testing" for the methodology and the corrected
% threshold thesis.tex had rounded down to 0.2 (that is the separate
% training-data-exclusion threshold of Section~\ref{sec:ch2:hsv}, not
% the velocity-linearity threshold).
\label{sec:ch2:sat_velocity}

Because saturation encodes field magnitude, and a robot's commanded
velocity in the navigation primitives of Chapter~\ref{ch:estimation}
scales with sensed magnitude, this testbed also validates that robot
velocity responds linearly to sensed saturation above the stiction
floor (Section~\ref{sec:ch2:dynamics}). To test this, lanes of constant
hue at varying saturation were printed, each robot was placed on each
lane and commanded to follow the hue direction at the velocity indicated
by that lane's saturation, and each robot's position was tracked over a
3-second period per lane \cite{39}. Across all four robots,
velocity scaled linearly with saturation above 0.3; saturation values
below 0.3 did not produce a commanded velocity strong enough to overcome
static friction \cite{39}. This 0.3 velocity-linearity threshold
is separate from the 0.2 saturation threshold below which readings were
excluded from calibration training (Section~\ref{sec:ch2:hsv}).

\section{Robot Dynamics Identification}
\label{sec:ch2:dynamics}
% Ported from: Systems paper Sec. III-A (robot controller layer, lines
% 192-206) for the identified values; thesis.tex Ch.4 Sec.
% "Robot and Cluster Dynamics" (around lines 2015-2058) for the same
% values as previously drafted; Draft_11.tex Sec. III (lines 553-559)
% and Table I (lines 561-590) for the pentagon-code operating point;
% omnibot.py docstring and constructor defaults for the code values.
% The first-order lag MODEL equations belong to Chapter 4
% (\ref{ch:estimation}); this section reports only identification and
% measured values.

Each robot is modeled, in both simulation and the control law derivation
of Chapter~\ref{ch:estimation}, as a single omnidirectional point mass
with a first-order actuator lag response, discretized at control period
$\Delta t$ as
\begin{equation}
\mathbf{v}[k+1] = \alpha_{\text{mom}}\mathbf{v}[k] + (1-\alpha_{\text{mom}})\mathbf{v}_{\text{des}}[k],
\qquad \alpha_{\text{mom}} = e^{-\Delta t/\tau},
\label{eq:ch2:momentum}
\end{equation}
where $\tau$ is the actuator time constant identified on the Decabot
hardware. The full continuous-time model this discretization comes from
is given in Chapter~\ref{ch:estimation}.

Robots are limited to a maximum speed of 0.3~m/s, and require a minimum
commanded speed to overcome static friction (stiction) below which
commanded motion does not produce actual motion.

The robot parameters were identified twice, once for the testbed paper
\cite{39} and again after the sensors and drive were recalibrated for
the critical-points work of Chapter~\ref{ch:first_order}. The current
identification gives a time constant $\tau \approx 0.28$~s, so
$\alpha_{\text{mom}} \approx 0.70$ at the 10~Hz control rate; earlier
drafts used a rounder $\tau = 0.3$~s ($\alpha_{\text{mom}} \approx
0.717$). The simulator uses the fixed value $\alpha_{\text{mom}} = 0.7$.
%% Note: omnibot.py uses the constant 0.7, which does not rescale if dt
%% changes (its docstring audit note, 2026-07-03).

The stiction floor is $v_{\text{stiction}} = 0.025$~m/s after
recalibration, the value used by the simulator; the testbed paper's
characterization, before recalibration, reported a usable velocity range
of 0.05 to 0.3~m/s \cite{39}. Draft
11's pentagon-formation operating point (Table~\ref{tab:ch2:params})
carries the same 0.025~m/s stiction floor and 0.3~m/s speed cap for its
double-gyre benchmark, lowering stiction to 0.002 (in the field's own
non-dimensional length units) and setting $\alpha_{\text{mom}} = 0$ for
its ocean HFR trial, where one control cycle corresponds to ten
real-world minutes and any physically reasonable vessel time constant is
negligible by comparison.

\begin{table}[h]
\centering
\caption{Robot dynamics operating points identified on Decabot hardware
and used across simulation trials. Double-gyre and ocean columns are
the second-order operating points of Chapter~\ref{ch:second_order}, in
non-dimensional length/time units; the Decabot column is the identified
hardware value used by the first-order work.}
\label{tab:ch2:params}
\begin{tabular}{lccc}
\toprule
Parameter & Decabot (identified) & Double gyre & Ocean \\
\midrule
$v_{\max}$ & 0.3~m/s & 0.3 & 0.3 \\
$v_{\text{stiction}}$ & 0.025~m/s & 0.025 & 0.002 \\
$\alpha_{\text{mom}}$ & 0.70 & 0.7 & 0 \\
$\tau$ (s) & 0.28 & -- & -- \\
\bottomrule
\end{tabular}
\end{table}
%% Verified 2026-09-29 against Draft_11.tex Table I (tab:params). Double-gyre
%% and ocean columns are in non-dimensional length/time units.

\section{Field Reconstruction as a Simulation Noise Model}
\label{sec:ch2:noise_model}
% Ported from: Systems paper Sec. V-A (lines 344-345) and Sec. IV
% (lines 239-244) for the reconstructed-field noise source; thesis.tex
% Ch.5 "Noise Model" (around lines 4031-4057) for the two-channel
% measurement/position noise formulation; src/control/pentagon_primitives.py
% docstring (lines 393-411) for the current code-level noise hooks.

The printed-map encoding of Section~\ref{sec:ch2:hsv} has imperfections
in both printing and sensing. Rather than model this error abstractly,
the critical-points work of Chapter~\ref{ch:first_order} reconstructed
two fields (a vortex and a saddle) directly from robot-sensed
measurements collected on the hardware testbed, and used these
reconstructed fields, and the interpolation models fit to them, as a
realistic noise condition for simulation, alongside the six analytical,
noise-free canonical fields. This ties the simulator's "noisy" condition
to a measured quantity, the same sensed reconstruction error shown in
Fig.~\ref{fig:ch2:hsv_field}, rather than to an assumed noise
distribution.

The second-order (pentagon) simulation code generalizes this into two
independent, parametric noise channels rather than one measured
reconstruction. Measurement noise is additive on each robot's sensed
$(u,v)$ reading; position noise perturbs the location at which a robot
samples the field, leaving its true, controlled position untouched. In
the current simulator these are exposed as \texttt{cluster.measurement\_noise\_std}
and \texttt{cluster.position\_noise\_std} on the six-robot cluster
object, both defaulting to zero. The three-robot simulator has no
parametric channel; its noisy condition is the reconstructed field itself
(\texttt{NNField}, \texttt{RBFField}, and \texttt{BlendedField} in
\texttt{src/fields/field\_types.py}). The derivation of how each channel
propagates into the fitted coefficients is given in
Chapter~\ref{ch:second_order}.

%% Verified 2026-09-29: pentagon_primitives.py lines 417-418 read both
%% attributes with default 0.0; AnalyticalField takes no noise argument.

\section{Simulation Environment}
\label{sec:ch2:simulator}
% Ported from: thesis.tex Ch.5 "Simulation Program" (around lines
% 3934-3957); Systems paper Sec. IV (lines 239-244); src/robot/omnibot.py,
% src/simulation/runner.py, src/fields/field_types.py.

Simulation trials throughout this dissertation use a Python framework
under \texttt{trunk/Python\_Simulations/Vector\_Fields/VF\_Robot/}, whose
entry point is \texttt{experiments/main\_omni.py}, configured by editing
module-level constants rather than command-line flags. Each robot is
represented as an \texttt{Omnibot} object implementing the first-order
momentum dynamics of Eq.~\eqref{eq:ch2:momentum}
(Section~\ref{sec:ch2:dynamics}), together with the maximum-speed and
stiction constraints of Table~\ref{tab:ch2:params}. The integrator is
explicit Euler at the robots' native 10~Hz control rate
($\Delta t = 0.1$~s); a simulation loop
(\texttt{src/simulation/runner.py}) advances each robot's commanded
velocity, integrates position, and, for time-varying fields, advances
the field's internal clock once per control step.

Vector fields are represented through a common interface,
\texttt{get\_value(x, y)}, implemented by an \texttt{AnalyticalField}
wrapper around a closed-form field function for the canonical fields of
Appendix~\ref{app:fields}, and by neural-network and reconstructed-field
classes for the hardware-derived fields of
Section~\ref{sec:ch2:noise_model}. The same interface, together with the
noise channels of Section~\ref{sec:ch2:noise_model} and the cluster-space
kinematics of Chapter~\ref{ch:estimation}, is reused unchanged across the
three-robot and six-robot formations studied in this dissertation.

\section{Summary of Prior and New Contributions}
\label{sec:ch2:summary}
% New text, modeled on Hart (2023) Ch. 2.4's pattern of crediting prior
% tools before listing new contributions.

Table~\ref{tab:ch2:contributions} separates what this chapter's tools
inherit from prior Robotic Systems Laboratory work from what this
dissertation adds to them.

\begin{table}[h]
\centering
\caption{Prior versus new contributions among this dissertation's
verification tools.}
\label{tab:ch2:contributions}
\begin{tabular}{ll}
\toprule
Prior work & New in this dissertation \\
\midrule
Decabot rovers \cite{41} & HSV field encoding \\
Low-cost indoor testbed \cite{42} & Neural-network sensor calibration \\
OptiTrack motion capture & Post-maintenance recalibration \\
Three-layer control architecture & Saturation-to-velocity characterization \\
 & Robot dynamics identification ($\tau$, stiction) \\
 & Field-reconstruction simulation noise model \\
 & Simulation environment (Omnibot, noise channels) \\
\bottomrule
\end{tabular}
\end{table}

The platform, the Decabot rovers and the indoor testbed they operate in,
is prior work. This dissertation's contribution to that platform is the
HSV encoding that replaces direct RGB field representation, the
neural-network calibration that corrects for the resulting sensor's
channel interference, and the recalibration, dynamics identification, and
simulation tooling built on top of it. These tools are used without
further comment for the remainder of this dissertation: the estimation
framework of Chapter~\ref{ch:estimation} and the zeroth-, first-, and
second-order results of Chapters~\ref{ch:zeroth}
through~\ref{ch:second_order} all run on the testbed and simulator
described here.

% ===================================================================
% CHAPTER 3: VECTOR FIELD PRELIMINARIES
% ===================================================================
\chapter{Vector Field Preliminaries}
\label{ch:preliminaries}

\section{Planar Vector Fields}
\label{sec:ch3:planar_fields}
% Ported from: Systems paper Sec. II-A (opening) and Draft_11.tex
% Sec. II-A (Local Vector Field Polynomials, opening). New synthesis
% text joins the two and frames the chapter.

A planar vector field assigns a velocity $\mathbf{v}(\mathbf{p}) = (u(\mathbf{p}), v(\mathbf{p})) \in \mathbb{R}^2$ to each position $\mathbf{p} = (x,y)$ in a domain $\mathcal{D} \subset \mathbb{R}^2$. In the applications this dissertation targets, $\mathbf{v}$ is a physical fluid velocity, a surface ocean current, a printed laboratory analogue, or a simulated benchmark flow, sensed by robots that occupy points of $\mathcal{D}$ and record the local flow.

A critical point $\mathbf{p}^* \in \mathcal{D}$ is a location where the field vanishes, $\mathbf{v}(\mathbf{p}^*) = \mathbf{0}$. Critical points organize the surrounding flow. A sink gathers floating material into a convergence zone; a source models an outflow or a leak; a vortex traps and carries a body of fluid, such as the plankton bloom inside a mesoscale ocean eddy, across a basin. Beyond isolated critical points, a flow can carry extended boundaries, the separatrices and coherent structures of Section~\ref{sec:ch3:coherent}, along which floating material collects or separates.

This chapter collects the field-theoretic definitions used in the rest of this dissertation: the velocity gradient tensor (Section~\ref{sec:ch3:jacobian}), the classification of critical points by its eigenvalues (Section~\ref{sec:ch3:critical_points}), the rotation-strain decomposition and the Okubo-Weiss partition it induces (Section~\ref{sec:ch3:ow}), the two families of coherent structure that motivate tracking a boundary rather than only a point (Section~\ref{sec:ch3:coherent}), what it means for a quantity built from the field to be objective (Section~\ref{sec:ch3:objectivity}), and the benchmark fields used throughout (Section~\ref{sec:ch3:benchmarks}).

\section{The Velocity Gradient Tensor}
\label{sec:ch3:jacobian}
% Ported from: thesis.tex Ch.4 "Flow Fields and Local Polynomial
% Models" (lines 2723-2757), Jacobian definition only; the polynomial
% estimator itself belongs to Chapter \ref{ch:estimation}.

The velocity gradient tensor, or field Jacobian, of $\mathbf{v}$ at $\mathbf{p}$ is
\begin{equation}
\mathbf{J}(\mathbf{p}) = \begin{bmatrix} \partial u/\partial x & \partial u/\partial y \\ \partial v/\partial x & \partial v/\partial y \end{bmatrix},
\label{eq:ch3:jacobian}
\end{equation}
abbreviated $u_x, u_y, v_x, v_y$ for the four partial derivatives where the point of evaluation is clear from context. This dissertation reserves bold $\mathbf{J}$ for the field Jacobian throughout. A kinematic Jacobian relating robot velocities to formation shape-parameter velocities is written $\mathbf{J}_c$ instead.

Two scalar invariants of $\mathbf{J}$ recur throughout this dissertation. The trace is the divergence of the field, $\operatorname{tr}(\mathbf{J}) = u_x + v_y = \nabla \cdot \mathbf{v}$, and is zero for an incompressible flow. The determinant,
\begin{equation}
D(\mathbf{p}) = \det(\mathbf{J}(\mathbf{p})) = u_x v_y - u_y v_x,
\label{eq:ch3:D_def}
\end{equation}
locates and classifies critical points (Section~\ref{sec:ch3:critical_points}) and, read as a field over a region rather than a value at a single point, marks the boundary between rotation-dominated and strain-dominated flow (Section~\ref{sec:ch3:ow}).

\section{Critical Points and Their Classification}
\label{sec:ch3:critical_points}
% Ported from: Systems paper Sec. II-A (lines 71-79, 88-104), Table I;
% thesis.tex Ch.3 "Vector Field Environments and Critical Points"
% (lines 1863-1879) for the Hartman-Grobman caveat on centers. New
% text for the trace-determinant plane derivation.

Near a critical point $\mathbf{p}^*$, the field is well approximated by its linearization, and the local topology is determined by $\mathbf{J}$ evaluated at $\mathbf{p}^*$. For a $2\times2$ system with $D(\mathbf{p}^*) \neq 0$, a hyperbolic or elliptic critical point, the eigenvalues of $\mathbf{J}(\mathbf{p}^*)$ fix the type. The characteristic equation is
\begin{equation}
\lambda^2 - \operatorname{tr}(\mathbf{J})\,\lambda + D = 0,
\label{eq:ch3:char_eq}
\end{equation}
with roots
\begin{equation}
\lambda = \frac{\operatorname{tr}(\mathbf{J}) \pm \sqrt{\operatorname{tr}(\mathbf{J})^2 - 4D}}{2}.
\label{eq:ch3:eig_roots}
\end{equation}
The sign of $D$, and where $D>0$ the sign of the discriminant $\operatorname{tr}(\mathbf{J})^2 - 4D$, partition the trace-determinant plane into the six types listed in Table~\ref{tab:ch3:field_summary} and shown in Fig.~\ref{fig:ch3:allfield}. A negative determinant gives real eigenvalues of opposite sign, a saddle. A positive determinant with a nonnegative discriminant gives real eigenvalues of the same sign, a node, stable for $\operatorname{tr}(\mathbf{J})<0$ and unstable for $\operatorname{tr}(\mathbf{J})>0$. A positive determinant with a negative discriminant gives a complex conjugate pair $\lambda = \alpha_{\text{eig}} \pm i\omega$, a spiral for $\alpha_{\text{eig}} \neq 0$ (stable if $\alpha_{\text{eig}}<0$, unstable if $\alpha_{\text{eig}}>0$) and a center for $\alpha_{\text{eig}}=0$.

\begin{table}[!t]
\centering
\caption{Classification of critical points by Jacobian eigenvalue structure ($\lambda$ = eigenvalue; a complex pair $\alpha_{\text{eig}} \pm i\omega$ has real part $\alpha_{\text{eig}}$, imaginary part $\omega$).}
\label{tab:ch3:field_summary}
\begin{tabular}{lll}
\toprule
\textbf{Field Type} & \textbf{Critical Point} & \textbf{Eigenvalues} \\
\midrule
Vortex & Center & $\lambda = \pm i\omega$ \\
Sinking Vortex & Stable Spiral & $\lambda = \alpha_{\text{eig}} \pm i\omega$, $\alpha_{\text{eig}} < 0$ \\
Sink & Stable Node & $\lambda_1, \lambda_2 < 0$ \\
Source & Unstable Node & $\lambda_1, \lambda_2 > 0$ \\
Spewing Vortex & Unstable Spiral & $\lambda = \alpha_{\text{eig}} \pm i\omega$, $\alpha_{\text{eig}} > 0$ \\
Saddle & Saddle Point & $\lambda_1 < 0 < \lambda_2$ \\
\bottomrule
\end{tabular}
\end{table}

\begin{figure}[!t]
\centering
\includegraphics[width=0.75\textwidth]{ch3_six_vector_fields.png}
% FIGURE SOURCE: Paper_Writing/Systems Paper/cwaight_systems_paper/six_vector_fields.png
\caption{Six linear vector field types with critical points.}
\label{fig:ch3:allfield}
\end{figure}

This classification holds rigorously for hyperbolic critical points. The center is the exception: a nonlinear perturbation of a linear center can produce a spiral, since the Hartman-Grobman theorem does not apply at $\alpha_{\text{eig}} = 0$, and under measurement noise an estimated center's real eigenvalue part hovers near zero, making center-versus-weak-spiral discrimination ill-posed in practice. $D(\mathbf{p}^*) = 0$ identifies degenerate cases, non-hyperbolic critical points or regions of near-uniform flow, that fall outside the scope of the estimation and control results in Chapters~\ref{ch:zeroth} through~\ref{ch:second_order}.

\section{Rotation, Strain, and the Okubo-Weiss Partition}
\label{sec:ch3:ow}
% Ported from: Draft_11.tex Sec. II-D (Eulerian Surrogates, lines
% 362-401, strain eigenvalue field and relation between the fields,
% excluding the observer-frame content); thesis.tex Ch.4 "The
% Okubo-Weiss Field and det(J)" (lines 2789-2805) for the Q formula.

The velocity gradient tensor decomposes into a symmetric and an antisymmetric part,
\begin{equation}
\mathbf{J} = \mathbf{S} + \mathbf{W}, \qquad \mathbf{S} = \tfrac{1}{2}(\mathbf{J}+\mathbf{J}^\top), \qquad \mathbf{W} = \tfrac{1}{2}(\mathbf{J}-\mathbf{J}^\top).
\label{eq:ch3:S_W_decomp}
\end{equation}
$\mathbf{S}$ is the rate-of-strain tensor. $\mathbf{W}$ is the spin tensor and carries the vorticity $\omega = v_x - u_y$ as $\mathbf{W} = \begin{bmatrix} 0 & -\omega/2 \\ \omega/2 & 0 \end{bmatrix}$.

Write the mean, deviatoric normal, and shear strain as $\mu = (u_x+v_y)/2$, $s_n = (u_x-v_y)/2$, and $s_s = (u_y+v_x)/2$. The eigenvalues of $\mathbf{S}$ are
\begin{equation}
s_{1,2} = \mu \mp \sqrt{s_n^2+s_s^2}, \qquad s_1 \le s_2,
\label{eq:ch3:s1s2}
\end{equation}
with eigenvectors $\mathbf{e}_1$ (the compression direction) and $\mathbf{e}_2$ (the stretching direction). For an incompressible flow $\mu = 0$, so $s_2 = -s_1$.

The Okubo-Weiss parameter partitions a flow into rotation-dominated and strain-dominated regions using only the local velocity gradient,
\begin{equation}
Q = \tfrac{1}{4}(s_n^2+s_s^2-\omega^2) = \tfrac{1}{4}(\nabla\cdot\mathbf{v})^2 - D,
\label{eq:ch3:Q_def}
\end{equation}
proposed by Okubo and, independently, Weiss \cite{43,44}. Hunt, Wray, and Moin generalized the same second-invariant idea to three dimensions as the Q-criterion for identifying eddies, streams, and convergence zones in turbulent flow \cite{45}. Hua and Klein revisited the two-dimensional criterion for oceanographic use, sharpening the conditions under which it separates elliptic and hyperbolic regions of a nearly two-dimensional turbulent flow \cite{46}, and the criterion has since been applied at scale in mesoscale eddy censuses \cite{7}.

For the incompressible flows considered in this dissertation, $\nabla\cdot\mathbf{v}=0$ and $Q=-D$ exactly, so $D$ carries the same partition with the opposite sign. The convention used throughout this dissertation is $D>0$ marking a rotation-dominated (elliptic) region and $D<0$ a strain-dominated (hyperbolic) region, with the zero level set $D=0$ the Okubo-Weiss boundary between them.

Where the vorticity vanishes, $\mathbf{J} = \mathbf{S}$ and $D = s_1 s_2$, which is $-s_1^2$ in incompressible flow: along such a curve the two fields share the same trench. Elsewhere the spin $\mathbf{W}$, which carries the vorticity, separates them, and it also separates them across observers (Section~\ref{sec:ch3:objectivity}).

Chapter~\ref{ch:second_order} uses trenches of $D$ and $s_1$ inside strain-dominated regions as surrogates for a separatrix or other coherent-structure boundary, estimated from a six-robot quadratic fit rather than computed from a known field; the argument for why those trenches mark the boundary, how the surrogates are estimated, and how a tracker rides one, is developed there.

\section{Coherent Structures}
\label{sec:ch3:coherent}

A coherent structure is a boundary that organizes material transport in a flow. Two frameworks define one, a Lagrangian framework built from particle trajectories and an instantaneous Eulerian framework built from the local velocity gradient.

\subsection{Lagrangian Coherent Structures and FTLE}
\label{sec:ch3:lcs_ftle}
% Ported from: Draft_11.tex Introduction (lines 112-114) and
% thesis.tex Ch.4 "Lagrangian Coherent Structures and the Objectivity
% Question" (lines 2541-2559, 2566-2572), citations converted to
% sep:N.

A separatrix in a steady flow is the union of the stable and unstable manifolds of a saddle point, partitioning the domain into regions whose trajectories do not mix. In a flow whose structure varies in time, the analogous object is a Lagrangian coherent structure (LCS), a distinguished material surface that organizes tracer transport over a finite time horizon \cite{47}. Haller's review frames the resulting body of theory as the search for the Lagrangian skeleton of a flow, and states the property such a skeleton must have to be physically meaningful: objectivity, invariance under the time-dependent Euclidean frame changes that a material response cannot depend on (Section~\ref{sec:ch3:objectivity}) \cite{48}. Shadden, Lekien, and Marsden gave the standard computable definition, an LCS as a ridge of the finite-time Lyapunov exponent (FTLE) field \cite{49}.

Computing an FTLE ridge requires integrating dense grids of particle trajectories over a finite time horizon, which presumes knowledge of the field's future evolution, in advance, over the whole region the tracers travel. A robot team that measures the flow only at its own footprint, at the current instant, cannot evaluate it. Nolan, Serra, and Ross show that FTLE fields converge to Eulerian rate-of-strain quantities as the integration horizon shrinks to zero \cite{50}, and Harms, Brunton, and McKeon show that standard gradient-based LCS diagnostics degrade on the sparse tracer trajectories typical of ocean drifters, proposing a regression-based alternative that linearizes in time rather than in space to recover usable gradients from sparse data \cite{51}.

\subsection{Objective Eulerian Coherent Structures}
\label{sec:ch3:oecs}
% Ported from: Draft_11.tex Sec. II-D (Eulerian Surrogates, lines
% 386-392) and Introduction (line 114); thesis.tex Ch.4 "Eulerian
% Surrogates" (lines 2601-2609) for the Cooper et al. magnitude-only
% comparison.

The second route to a coherent structure is instantaneous: a scalar or tensor field built from the local velocity gradient alone, evaluated at a single point in time, with no trajectory integration. The Okubo-Weiss parameter of Section~\ref{sec:ch3:ow} is one such criterion. Because these criteria are properties of the instantaneous field rather than of material transport, Haller's review excludes them from the Lagrangian framework: they can frame a coherent feature of the velocity field, but on their own they are not objective \cite{48}. Reducing a sensed field to a scalar this way can also lose the information needed to tell structures apart; Cooper et al.'s initial multirobot study of vector field navigation reduced each measurement to its magnitude, which yields a descent direction but cannot distinguish a sink from a saddle from a vortex, since the magnitude vanishes at every critical point regardless of type \cite{38}.

Serra and Haller placed an instantaneous criterion on objective footing by defining objective Eulerian coherent structures (OECS) from the eigenvalues and eigenvectors of the rate-of-strain tensor $\mathbf{S}$ rather than from the raw velocity gradient, and showed that OECS are the limit LCS approach as the integration horizon shrinks to zero \cite{52}. An attracting OECS is a trench of $s_1$ tangent to the stretching eigenvector $\mathbf{e}_2$; floating material has been observed at sea to collect onto such curves \cite{11}. A repelling OECS is a ridge of $s_2$ tangent to the compression eigenvector $\mathbf{e}_1$, along which neighboring material separates. In incompressible flow $s_2=-s_1$, so a repelling structure is also a trench of $s_1$, and a single tracker built on $s_1$ can ride both (Chapter~\ref{ch:second_order}).

\section{Objectivity and Observer Frames}
\label{sec:ch3:objectivity}
% Ported from: Draft_11.tex Sec. II-D (Eulerian Surrogates, lines
% 400-411, general transformation law and which quantities are
% objective) and Sec. VI-C ("Behavior Under a Rotating Observer",
% lines 946-947, frame transformation) for the operational definition.
% The rotating-observer trial results belong to Chapter
% \ref{ch:second_order}, not this chapter.

Objectivity, also called frame-indifference, requires a physical quantity computed from the flow to take the same value for any two observers related by a time-dependent Euclidean transformation, a rotation and a translation, since such a transformation changes nothing about the material motion itself \cite{48}. This matters for a robot cluster whose heading is dead-reckoned rather than directly measured: the cluster is, in effect, a rotating observer of the field it senses.

For an observer rotating at rate $\Omega$ with rotation matrix $\mathbf{Q}$, the observed field is
\begin{equation}
\mathbf{v}'(\mathbf{p}', t) = \mathbf{Q}\,\mathbf{v}(\mathbf{Q}^\top\mathbf{p}') + \Omega\,(-y', x')^\top,
\label{eq:ch3:rotating_frame}
\end{equation}
the field rotated into the moving frame plus the frame's own solid-body rotation.

Not every quantity built from $\mathbf{J}$ is objective under this transformation. The rate-of-strain tensor $\mathbf{S}$ co-rotates with the frame, so its eigenvalues, including $s_1$, are objective. The spin tensor $\mathbf{W}$ and the vorticity are not: the observed vorticity is $\omega' = \omega + 2\Omega$, the true vorticity plus the frame's own rotation rate. Because $D$ mixes $\mathbf{S}$ and $\mathbf{W}$, it is not objective either. Expanding $D$ in the rotating frame,
\begin{equation}
D' = D + \Omega\,\omega + \Omega^2,
\label{eq:ch3:D_not_objective}
\end{equation}
so two observers rotating relative to each other disagree about where the trenches of $D$ lie, while they agree about the trenches of $s_1$ \cite{52}. This is the tradeoff Chapter~\ref{ch:second_order} evaluates between the two surrogates of Section~\ref{sec:ch3:ow}.

\section{Benchmark Fields}
\label{sec:ch3:benchmarks}

\subsection{Canonical Linear Fields}
\label{sec:ch3:linear_fields}
% Ported from: Systems paper Sec. VII "Vector Field Environments"
% (lines 510-526); thesis.tex Appendix A (lines 5066-5107) for the
% $\rho,\theta$ renaming.

Six linear vector field environments, shown in Fig.~\ref{fig:ch3:allfield} and classified in Table~\ref{tab:ch3:field_summary}, are used as simulation benchmarks in this dissertation. Centered at the origin, with $\rho=\sqrt{x^2+y^2}$ and $\theta=\operatorname{atan2}(y,x)$, the base fields are the vortex (pure rotational flow), the sink (radial inward flow), whose negation is the source (radial outward flow), and a saddle rotated $45^\circ$ so its stable and unstable directions lie along the diagonals $y=x$ and $y=-x$. The sinking vortex (spiraling inward) and spewing vortex (spiraling outward) are linear combinations of the vortex and sink fields. Full closed forms are given in Appendix~\ref{app:fields}.

\subsection{The Steady Double Gyre}
\label{sec:ch3:double_gyre}
% Ported from: Draft_11.tex Sec. II-E "The Double-Gyre Example" (lines
% 413-491), near verbatim; internal section references retargeted to
% this chapter's labels.

The double gyre is a benchmark flow on which both surrogates of Section~\ref{sec:ch3:ow} mark the material transport structure exactly, and every quantity of interest has a closed form. The flow is two counter-rotating gyres in a closed basin, separated by a separatrix that runs between two saddle points on the basin walls. This dissertation uses the steady member of the Shadden family \cite{49}, shifted from the canonical domain $[0,2]\times[0,1]$ to $[-1,1]\times[-0.5,0.5]$ to center it on the origin. The shift is $x_f = x+1$, $y_f = y+0.5$, where $(x_f,y_f)$ are field-frame and $(x,y)$ world-frame coordinates. Both are non-dimensional throughout the double-gyre benchmark. The field is
\begin{equation}
u = -\pi A \sin(\pi x_f)\cos(\pi y_f), \quad v = \pi A \cos(\pi x_f)\sin(\pi y_f),
\label{eq:ch3:dg_u}
\end{equation}
with $A=0.1$. The separatrix is the line $x=0$, the saddle points are $\mathbf{p}_1^* = (0,-0.5)$ and $\mathbf{p}_2^* = (0,0.5)$, and the gyre centers are $(\pm0.5, 0)$.

The analytical Jacobian entries are $u_x = -\pi^2 A \cos\pi x_f \cos\pi y_f = -v_y$ and $u_y = \pi^2 A \sin\pi x_f \sin\pi y_f = -v_x$. The flow is therefore incompressible, and its determinant and Hessian are
\begin{equation}
\begin{aligned}
D(\mathbf{p}) &= -\tfrac{1}{2}\pi^4 A^2 \left(\cos 2\pi x_f + \cos 2\pi y_f\right), \\
\mathbf{H}_D &= 2\pi^6 A^2 \begin{bmatrix} \cos 2\pi x_f & 0 \\ 0 & \cos 2\pi y_f \end{bmatrix}.
\end{aligned}
\label{eq:ch3:det_analytic}
\end{equation}
The Hessian is diagonal everywhere, so its eigenvectors lie on the coordinate axes.

The landscape of $D$ is a grid of trenches around rotation-dominated gyres, each inside an Okubo-Weiss diamond $x_f \pm y_f \in \tfrac{1}{2}+\mathbb{Z}$ where $D>0$. The trenches lie on the lines $x \in \{-1,0,1\}$ and $y=\pm0.5$, the separatrix and the four basin walls, which cross at right angles at $\mathbf{p}_1^*$, $\mathbf{p}_2^*$, and the four domain corners.

Along the separatrix $D = -\pi^4 A^2 \sin^2\pi y$. The trench is deepest at the two saddles and rises to a crest, $D=0$, at the origin (Fig.~\ref{fig:ch3:dg_fields}(b)). The crest divides the separatrix into an upper and a lower segment. On $x=0$ the Hessian reduces to $\mathbf{H}_D = 2\pi^6 A^2\,\operatorname{diag}(1,-\cos2\pi y)$. The transverse curvature $2\pi^6 A^2$ is the same along the whole separatrix. The along-trench curvature is negative for $|y|<0.25$ and positive beyond, so each segment is concave along its length on its inner half and convex on its outer half. At a saddle both curvatures are positive, and $D$ has a local minimum where the separatrix crosses the wall.

The vorticity ties the $s_1$ field to the $D$ field on this flow. It is $\omega = -2\pi^2 A \sin\pi x_f \sin\pi y_f$, which vanishes where $x_f$ or $y_f$ is an integer, exactly on the trench grid of $D$. There $\mathbf{J} = \mathbf{S}$ and $D=-s_1^2$ (Section~\ref{sec:ch3:ow}), so the two fields share every trench. The shear strain vanishes identically, so $\mathbf{S} = \operatorname{diag}(u_x,-u_x)$ everywhere. This gives
\begin{equation}
s_1 = -\pi^2 A \left|\cos\pi x_f \cos\pi y_f\right|,
\label{eq:ch3:s1_analytic}
\end{equation}
with minima $s_1 = -\pi^2 A$ where the grid lines cross, at $\mathbf{p}_1^*$, $\mathbf{p}_2^*$, and the four domain corners. The eigenvectors of $\mathbf{S}$ lie on the coordinate axes except on the lines $x=\pm0.5$ and $y=0$ through the gyre centers, where $u_x=0$, $\mathbf{S}$ vanishes, and both eigenvectors are undefined. The line $y=0$ meets the separatrix at the origin, an isotropic point of $\mathbf{S}$ at the $D$ crest.

On the upper segment $u_x<0$, so the stretching eigenvector $\mathbf{e}_2$ is tangent to the trench. On the lower segment $u_x>0$, and the compression eigenvector $\mathbf{e}_1$ is tangent. By the OECS definition of Section~\ref{sec:ch3:oecs}, the upper segment is an attracting OECS and the lower a repelling OECS. The two meet at the isotropic point, where the tangent swaps (Fig.~\ref{fig:ch3:dg_fields}(c)).

\begin{figure}[!tb]
\centering
\includegraphics[width=0.75\textwidth]{ch3_double_gyre_fields.png}
% FIGURE SOURCE: Paper_Writing/Separatrix_and_OW_Paper/figures/double_gyre_fields.png
\caption{Steady double gyre. (a) Streamlines, with the separatrix $x=0$ (dashed), the saddle points (crosses), and the gyre centers (dots). (b) Signed $D$ as an elevation surface. The gyre centers rise as peaks ($D>0$), and the separatrix (bold) is a trench of $D$, deepest at the saddles and rising to $D=0$ at the origin. (c) $s_1$ over the same domain. The separatrix is a trench on both segments, deepest at the saddles. The trench tangent is $\mathbf{e}_2$ on the upper, attracting segment and $\mathbf{e}_1$ on the lower, repelling segment, and the two swap at the isotropic point of $\mathbf{S}$ at the origin, where the $D$ crest of (b) sits.}
\label{fig:ch3:dg_fields}
\end{figure}

% ===================================================================
% CHAPTER 4: COOPERATIVE ESTIMATION FRAMEWORK
% ===================================================================
\chapter{Cooperative Estimation Framework}
\label{ch:estimation}

% New text
Every primitive in this dissertation runs on the same control
architecture and follows the same estimation pattern. The robots sample
the field at one instant, a local polynomial is fit to those samples, and
a control law steers the cluster on what the fit reveals. This chapter
describes the architecture and the estimation pattern once.

\section{Multilayer Control Architecture}
\label{sec:ch4:architecture}

% Ported from: Draft_11_6ec7e51.tex
The control architecture is the established multilayer approach of prior
cluster space adaptive navigation studies \cite{2}, shown in
Fig.~\ref{fig:ch4:control_architecture}. There are three layers: (1) an
adaptive navigation layer that estimates the local field and generates the
cluster velocity command, (2) a cluster controller layer that maintains
the formation and translates between cluster space and robot space
variables, and (3) a robot controller layer that converts each robot's
velocity command into actuation. Field samples enter only at the adaptive
navigation layer. The layers below it operate on velocity commands and
positions alone.

% Ported from: IDETC 2025 (DETC2025-167604), Sec. 2.1, transcribed from the published PDF
This architecture was selected because its modular layered approach
allows for independent development and testing of different components
of the overall architecture. This approach also maximizes system
reusability, enabling systematic comparison of navigation strategies
using standardized components in downstream layers. Additionally, the
architecture enables the same control laws to be applied across different
robot platforms and formations.

% Ported from: cwaight_systems_paper.tex
At each control cycle (10~Hz), each robot reports its position and vector field measurement. The cluster space kinematics are computed, and the cluster shape velocity commands needed for formation control are determined. The adaptive navigation layer computes the critical point estimate and generates a cluster velocity command. The cluster space controller maps all velocity commands to individual robot velocities via the inverse kinematic Jacobian.

% Ported from: cwaight_systems_paper.tex
\begin{figure}[htbp]
    \centering
    \includegraphics[width=0.55\textwidth]{ch4_control_architecture.png}
    \caption{Hierarchical control architecture: robot, cluster space, and adaptive navigation layers.}
    \label{fig:ch4:control_architecture}
\end{figure}

\subsection{Robot Controller Layer}
\label{sec:ch4:robot_layer}

% Ported from: cwaight_systems_paper.tex
Each robot's controller receives velocity commands and translates them into motion. Each robot also samples the local vector field $\mathbf{v}(\mathbf{p}_i) = (u_i, v_i)$ and reports its global position $\mathbf{p}_i$, providing the inputs to the estimation framework of Section~\ref{sec:ch4:estimation}.

% New text
Section~\ref{sec:ch2:dynamics} gives the first-order lag, speed limit, and
stiction floor identified for the Decabots, which the simulations of
later chapters reproduce.

\subsection{Cluster Space Controller Layer}
\label{sec:ch4:cluster_space}

% Ported from: IDETC 2025, Sec. 2.3, transcribed from the published PDF
To maintain a fixed formation while commanding all robots with individual
velocity commands, a cluster space controller was used. Cluster space
control is a technique developed by Dr. Chris Kitts of Santa Clara
University \cite{52}. In this technique, all robot positions are
described geometrically relative to a cluster frame via a set of
kinematic equations.

% New text
It is an operational space control approach in which the multirobot
formation is represented as a virtualized full degree-of-freedom
articulating mechanism.

% Ported from: cwaight_systems_paper.tex
We refer to the group of robots as a cluster: a formation treated as a single virtual rigid body whose centroid, orientation, and internal geometry can be independently specified and controlled \cite{52,53}.

% Ported from: IDETC 2025, Sec. 2.3, transcribed from the published PDF
Cluster space control also provides a closed form method for converting
between cluster-level commands and robot-level commands. This allows for
the user or adaptive navigation layer to specify the behaviour commands
for the entire cluster, while the individual robots are then commanded by
the robot control layer to react accordingly. These conversions happen on
a centralized computer, after the individual robots communicate their
instantaneous position and sensed information.

\subsubsection{Three-Robot SAS Formation}
\label{sec:ch4:sas}

% Ported from: cwaight_systems_paper.tex
The cluster space controller maintains the cluster's formation geometry while executing centroid velocity commands from the adaptive navigation layer. It translates cluster-level velocity commands into individual robot velocity commands using an inverse kinematic Jacobian, and transforms robot position data back into cluster variables for feedback control.

% Ported from: cwaight_systems_paper.tex
Chapters~\ref{ch:zeroth} and~\ref{ch:first_order} consider a cluster of three robots forming a triangular configuration. The shape can be parameterized using Side-Angle-Side (SAS) representation with parameters $(p, \beta, q)$, where $p$ is the distance between robots 1 and 2, $\beta$ is the interior angle at robot 2, and $q$ is the distance between robots 2 and 3, as visualized in Fig.~\ref{fig:ch4:threerobot}.

% Ported from: cwaight_systems_paper.tex
\begin{figure}[htbp]
    \centering
    \includegraphics[width=0.45\textwidth]{ch4_SAS_Robots.png}
    \caption{Three-robot triangular formation parameterized by SAS variables: side lengths $p$, $q$ and interior angle $\beta$.}
    \label{fig:ch4:threerobot}
\end{figure}

% Ported from: cwaight_systems_paper.tex
The forward kinematic equations map the robot positions $(x_1, y_1), (x_2, y_2), (x_3, y_3)$ to the cluster shape parameters, and the inverse kinematic equations map cluster variables back to robot positions. Differentiating them gives analytic forms for the cluster space kinematic Jacobian and its inverse, which translate velocities between robot space and cluster space. The full equations are given in Appendix~\ref{app:kin3}.

\subsubsection{Six-Robot Pentagon-Plus-Center Formation}
\label{sec:ch4:pentagon}

% Ported from: Draft_11_6ec7e51.tex
Cluster space control treats a multirobot formation as a virtual rigid
body that can be moved, scaled, and articulated \cite{52}. For six planar
robots the cluster state has twelve variables. The centroid $(x_c, y_c)$
and heading $\theta_c$ place the formation, and nine shape variables set
its geometry. The robots form three pairs whose midpoints define an SAS
triangle $(p, \beta, q)$ (Fig.~\ref{fig:ch4:pentagon}). The cluster controller layer maps the
cluster-space velocity vector to robot velocities through the analytic
inverse Jacobian $\mathbf{J}_c^{-1}$ and returns measured robot positions to
the layer above through the forward kinematics (Appendix~\ref{app:kin6}).

% Ported from: Draft_11_6ec7e51.tex
\begin{figure}[htbp]
    \centering
    \includegraphics[width=0.45\textwidth]{ch4_pentagon_formation.png}
    \caption{Pentagon-plus-center formation. Robot 1 sits at the centroid
    $\mathbf{p}_c$ and robots 2 to 6 on a ring (dotted). The cluster
    variables shown are the pair separations $L_2$, $L_3$, $L_4$, a pair
    orientation $\theta_3$, the SAS parameters $(p, \beta, q)$ of the
    triangle of pair midpoints (squares), and the heading $\theta_c$.}
    \label{fig:ch4:pentagon}
\end{figure}

\subsection{Adaptive Navigation Layer}
\label{sec:ch4:an_layer}

% Ported from: IDETC 2025, Sec. 2.4, transcribed from the published PDF
The adaptive navigation layer integrates the cluster shape policy,
feature estimator, and control law, determining the desired actions of
the entire cluster in cluster space variables.

% New text
The feature estimator is the part that changes from rung to rung.
Section~\ref{sec:ch4:estimation} describes it in general, and each of
Chapters~\ref{ch:zeroth} to~\ref{ch:second_order} gives the control law
that consumes it.

\section{Local Polynomial Field Estimation}
\label{sec:ch4:estimation}

% New text. Generalizes the first-order fit of the critical-points paper
% and the second-order fit of Draft 11 to order k; checked 2026-09-29.
Each estimator fits a polynomial to the velocity samples the robots take
at one instant. The order of that polynomial decides what the cluster can
learn about the field.

\subsection{Monomial Basis and Formation Matrix}
\label{sec:ch4:basis}

Let $\boldsymbol{\phi}_k(\mathbf{p})$ collect the monomials $x^a y^b$
with $a + b \leq k$, each scaled by $1/(a!\,b!)$, with coordinates
measured from the cluster centroid. The scaling makes each fitted
coefficient a Taylor derivative of the field at the centroid. The basis
has
\begin{equation}
    n_k = \tfrac{1}{2}(k+1)(k+2)
    \label{eq:ch4:nk}
\end{equation}
entries. Robot $i$ at $\mathbf{p}_i$ contributes the row
$\boldsymbol{\phi}_k(\mathbf{p}_i)^{\top}$ to the formation matrix
$\boldsymbol{\Phi}_k$, and each velocity component is fit separately,
\begin{equation}
    \boldsymbol{\Phi}_k\,\hat{\boldsymbol{\theta}}_u = \mathbf{m}_u, \qquad
    \boldsymbol{\Phi}_k\,\hat{\boldsymbol{\theta}}_v = \mathbf{m}_v ,
    \label{eq:ch4:fit}
\end{equation}
where $\mathbf{m}_u$ and $\mathbf{m}_v$ stack the robots' readings. The
first-order fit of Chapter~\ref{ch:first_order} is the case $k = 1$, and
the second-order fit of Chapter~\ref{ch:second_order} is $k = 2$. The two
chapters order the columns differently, which changes nothing in the fit.

\subsection{Minimality and Degeneracy}
\label{sec:ch4:minimality}

Each robot adds one equation per component, so a unique fit needs at
least $n_k$ robots. The first order needs three and the second order
needs six. With exactly $n_k$ robots $\boldsymbol{\Phi}_k$ is square. It
is singular exactly when a nonzero coefficient vector $\mathbf{c}$
satisfies $\boldsymbol{\Phi}_k\mathbf{c} = \mathbf{0}$, which says that
the polynomial $\mathbf{c}^{\top}\boldsymbol{\phi}_k(\mathbf{p})$
vanishes at every robot. The fit therefore fails exactly when the robots
lie on a common algebraic curve of degree $k$. For $k = 1$ that curve is
a line, and for $k = 2$ it is a conic. With more than $n_k$ robots the
system is overdetermined and is solved by least squares.

\subsection{Formation Size}
\label{sec:ch4:size}

Formation size enters through the columns of $\boldsymbol{\Phi}_k$. If a
formation keeps its shape and is scaled to radius $\rho$, a monomial of
degree $q$ scales as $\rho^{\,q}$, so the noise reaching a coefficient of
order $q$ scales as $\rho^{-q}$ (Section~\ref{sec:ch8:noise_order}).
Truncation error runs the other way. The neglected terms are of degree
$k + 1$, so the bias they leave grows with $\rho$. A larger formation
therefore trades truncation bias for noise suppression.

\section{The Fit-Order Ladder}
\label{sec:ch4:ladder}

Table~\ref{tab:ch4:ladder} organizes this dissertation by fit order. Each
rung fixes a failure of the one below it. The zeroth-order primitives of
Chapter~\ref{ch:zeroth} use the mean sensed vector, or the gradient of
the sensed magnitude, and never assemble the velocity gradient. They
supply a heading, which drifts on an orbit. The first-order fit of
Chapter~\ref{ch:first_order} locates and classifies a critical point, but
it gives the velocity gradient at one point only, so it cannot follow a
curve. The second-order fit of Chapter~\ref{ch:second_order} gives the
velocity gradient as a field over the formation, which is enough to
acquire and ride a separatrix. A third-order fit from ten robots would
recover the third derivatives the second-order trackers lack, and
Chapter~\ref{ch:limitations} leaves it as future work.

\begin{table}[htbp]
\centering
\caption{The fit-order ladder. The minimum robot count is $n_k$ of
(\ref{eq:ch4:nk}), and the degeneracy is the curve on which that many
robots cannot support the fit.}
\label{tab:ch4:ladder}
\begin{tabular}{>{\raggedright\arraybackslash}p{1.6cm} c >{\raggedright\arraybackslash}p{2.4cm} >{\raggedright\arraybackslash}p{5.2cm} c}
\hline
Fit order & Robots & Degenerate on & Features reachable & Chapter \\
\hline
Zeroth & 1 & none & heading only, which drifts on orbits & \ref{ch:zeroth} \\
First & 3 & common line & critical points, with location, type, and orbit & \ref{ch:first_order} \\
Second & 6 & common conic & separatrices, OECS, flow saddles as minima & \ref{ch:second_order} \\
Third & 10 & common cubic & third derivatives, the true $\mathbf{H}_D$ and transverse curvature & \ref{ch:limitations} \\
\hline
\end{tabular}
\end{table}

% ===================================================================
% CHAPTER 5: ZEROTH-ORDER BASELINES
% ===================================================================
\chapter{Zeroth-Order Baselines}
\label{ch:zeroth}

% Ported from DETC2025-167604, "A Functional Indoor Testbed for Multirobot
% Adaptive Navigation in Vector Field Environments" (ref. 39), text at
% PhD Thesis/sources/idetc2025_testbed.txt. Its HSV encoding and NN
% calibration content is Chapter~\ref{ch:tools}'s. The unpublished 2024
% DETC2024-142746 manuscript is deliberately not used (author decision,
% 2026-09-29).

\section{Direction-Only Navigation}
\label{sec:ch5:direction_only}
% Ported from: idetc2025_testbed.txt Secs. 1, 2.5; opening paragraph new.

This chapter presents the zeroth-order control primitives of this
dissertation: primitives that command a cluster's heading from an
instantaneous, unweighted read of the team's vector-field measurements,
without fitting any spatial derivative of the field. Each robot returns a
single instantaneous vector reading of the local field, and the adaptive
navigation layer combines these readings into a commanded cluster
velocity, using the multilayer control architecture and cluster-space
formation control introduced in Chapters~\ref{ch:tools}
and~\ref{ch:estimation}. Robots are treated as omnidirectional; robot
orientation is not used, and the cluster's own orientation pose is not
commanded by any primitive in this chapter.

Two zeroth-order primitives from prior work \cite{38} are examined here:
the vector-sum technique, which uses the direction of the summed
readings, and the vector-to-scalar technique, which discards direction
and uses only the magnitude of each reading. Both were first validated
in simulation, with known limitations documented \cite{38}, and were
the first control primitives run on the Decabot testbed of
Chapter~\ref{ch:tools} \cite{39}
(Section~\ref{sec:ch5:hardware_trials}). Both are purely reactive: each
commands a velocity from the current instantaneous readings alone, with
no state carried between control cycles.

\section{Vector-Sum Primitive}
\label{sec:ch5:vector_sum}
% Ported from: idetc2025_testbed.txt Sec. 2.5.

The vector-sum technique combines the vectors sensed by each robot into
a resultant vector. The normalized direction of this resultant then
specifies the commanded cluster velocity $(\dot{x}_c,\dot{y}_c)$
\cite{39}. Because the readings are not normalized before summing, a
reading with larger magnitude carries proportionally more weight in the
resultant direction.

\section{Vector-to-Scalar Primitive}
\label{sec:ch5:vector_to_scalar}
% Ported from: idetc2025_testbed.txt Sec. 2.5, Eqs. 2-7.

In the vector-to-scalar technique, only the magnitude of each robot's
vector reading is used, reducing navigation to a scalar-field gradient
problem. Treating each robot's reading magnitude as a $z$-coordinate over
its $(x,y)$ position, the differences between robots' readings determine
the normal to the plane through the three points, which is then projected
to give the direction of steepest change. For robots $1$, $2$, $3$ at
positions $(x_i,y_i)$ with reading magnitudes $z_i$,
\begin{equation}
\vec{R}_{12} = \begin{pmatrix} x_2-x_1 \\ y_2-y_1 \\ z_2-z_1 \end{pmatrix}, \qquad
\vec{R}_{13} = \begin{pmatrix} x_3-x_1 \\ y_3-y_1 \\ z_3-z_1 \end{pmatrix},
\label{eq:ch5:R_edges}
\end{equation}
\begin{equation}
\vec{N} = -\vec{R}_{13} \times \vec{R}_{12},
\label{eq:ch5:normal}
\end{equation}
\begin{equation}
b_{\text{grad}} = \operatorname{ATAN2}(N_y, N_x) + \frac{\pi}{2},
\label{eq:ch5:bgrad}
\end{equation}
and the angle is converted to a commanded cluster velocity \cite{39},
\begin{equation}
\dot{x}_c = \cos(b_{\text{grad}}), \qquad
\dot{y}_c = \sin(b_{\text{grad}}).
\label{eq:ch5:vel_cmd}
\end{equation}
%% Corrected from source Eq. 7 (cos -> sin): as printed in
%% idetc2025_testbed.txt / DETC2025-167604, the pair reads
%% $\dot{x}_c = \cos(b_{\text{grad}})$ (Eq. 6, correct) and
%% $\dot{y}_c = \cos(b_{\text{grad}})$ (Eq. 7, printed with cos where sin
%% is needed for $(\dot{x}_c,\dot{y}_c)$ to be a unit vector at angle
%% $b_{\text{grad}}$). The affected equation is source Eq. 7, immediately
%% following the correctly-printed Eq. 6.

The gradient becomes ill-defined if the three robots are collinear.
Equation~\eqref{eq:ch5:bgrad} gives the direction of steepest ascent; the
technique navigates to a local extremum by moving in this direction
(source-seeking, or maximum-seeking) or its negation (sink-seeking, or
minimum-seeking, used for the hardware trials of
Section~\ref{sec:ch5:hardware_trials}), and orbits an extremum by moving
perpendicular to it.

\section{Hardware Trials on Fixed and Sinking Vortices}
\label{sec:ch5:hardware_trials}
% Ported from: idetc2025_testbed.txt Secs. 2.4-2.5 (trial cluster and
% primitive setup), 3.2 (Eqs. 8-13, field definitions), 3.5 (trial
% methodology), 4.2 (results); DETC2025-167604 \cite{39}. The
% Decabot testbed itself (HSV encoding, NN calibration, architecture) is
% Chapter~\ref{ch:tools}'s.

The vector-sum and vector-to-scalar primitives of
Sections~\ref{sec:ch5:vector_sum} and~\ref{sec:ch5:vector_to_scalar},
previously demonstrated only in simulation, were demonstrated on the
Decabot hardware testbed of Chapter~\ref{ch:tools} in \cite{39}, on
a fixed vortex and a sinking vortex field. The commanded formation was a
three-robot cluster with a constant shape policy, $p = q = 0.35$~m and
$\beta = 1.05$~rad, forming an equilateral triangle.

The first-order trials of Chapter~\ref{ch:first_order} used a smaller,
$0.33$~m equilateral formation.

Both fields are built from a fixed vortex centered at
$(\mathrm{center}_x,\mathrm{center}_y)$: with
$r = \sqrt{(x-\mathrm{center}_x)^2+(y-\mathrm{center}_y)^2}$ and
$\theta = \operatorname{atan2}(y-\mathrm{center}_y,\, x-\mathrm{center}_x)$,
\begin{equation}
u = r\sin\theta, \qquad v = -r\cos\theta.
\label{eq:ch5:fixed_vortex}
\end{equation}
The sinking vortex adds a radially inward term to this fixed vortex,
\begin{equation}
u \mathrel{+}= -\frac{c_s(x-\mathrm{center}_x)}{r^2+10^{-6}}, \qquad
v \mathrel{+}= -\frac{c_s(y-\mathrm{center}_y)}{r^2+10^{-6}},
\label{eq:ch5:sinking_vortex}
\end{equation}
a point sink of strength $c_s = 0.2$ whose inward speed decays as $1/r$, with
the $10^{-6}$ term added only to avoid division by zero at the center. After construction, both fields are rescaled so that vector
magnitude falls in $[0.3,1]$ for printing; magnitudes below 0.2 were
excluded from calibration, and below 0.3 did not produce a robot velocity
sufficient to overcome static friction.
%% Sign: DETC2025-167604 Eqs. 12-13 print the sink term with a plus sign
%% (outward), a typo; the hardware plotting script
%% trunk/robots_3/saddle_tests/plotter_avg_run_Sinking_Vortex.m uses the
%% inward form above, and the trials spiral inward.
%% c_s = 0.2 follows the testbed paper's own plotting package
%% (trunk/robots_3/Results_for TestBed_Paper/plotting_package/
%% plotter_avg_run_Sinking_Vortex.m); the published equations print 0.1.
%% Author decision 2026-09-29: use the code value.
%% This 1/r sink differs from the linear sinking vortex of
%% Chapter~\ref{ch:first_order} (Appendix~\ref{app:fields}).

Testing used a fixed 10~Hz OptiTrack sampling rate and 10 trial runs per
environment-primitive combination. The fixed vortex was tested with both
the vector-sum and vector-to-scalar (minimum-seeking) primitives; the
sinking vortex was tested only with vector-sum. Both fixed and randomized
starting positions were used for the vector-to-scalar trials on the fixed
vortex. Performance was evaluated by trajectory shape relative to the
expected path, steady-state error between final and target cluster
position, and formation-shape maintenance (Section~\ref{sec:ch5:formation_stats}).

On the fixed vortex, the vector-sum primitive produced a consistent
outward drift across all 10 runs, with the drift rate stable at all
angles, a behavior not predicted in simulation. Section~\ref{sec:ch5:why_not_location}
returns to this result. On the sinking vortex, the same primitive settled
the cluster to within approximately 0.1~m of the true center, with the
target location falling within the triangle connecting the three robots.
%% NOTE: the 0.2 m value comes from the text of DETC2025-167604 (third test
%% set, vector-to-scalar, fixed vortex, its Fig. 18), but that paper's Fig. 18
%% caption says "sinking vortex". Text and caption disagree in the published
%% paper; the text is followed here. Check against the Fig. 18 data.
The vector-to-scalar (minimum-seeking) primitive on the fixed vortex
settled the cluster within approximately 0.2~m of center, following an
approximately straight-line path with some deviation attributed to
printing noise and to variance in the saturation-to-magnitude estimate;
performance was essentially unchanged when starting positions were
randomized rather than fixed, since the velocity command is computed from
the local gradient of saturation rather than from saturation itself, and
near the center this gradient commands speeds below 0.05~m/s, too low
to overcome static friction as the robots were characterized at the
time (the floor was later remeasured at 0.025~m/s after recalibration,
Chapter~\ref{ch:tools}), regardless of where the cluster started. Neither orbiting nor saddle-point navigation, the subjects of
Chapter~\ref{ch:first_order}, was demonstrated on hardware at this stage;
both require the Jacobian-based estimation developed there.

\begin{figure}[h]
  \centering
  \begin{subfigure}[t]{0.32\textwidth}
    \centering
    \includegraphics[width=\textwidth]{ch5_fixed_vortex_vector_sum.png}
    \caption{Vector-sum, fixed vortex.}
  \end{subfigure}\hfill
  \begin{subfigure}[t]{0.32\textwidth}
    \centering
    \includegraphics[width=\textwidth]{ch5_sinking_vortex_vector_sum.png}
    \caption{Vector-sum, sinking vortex.}
  \end{subfigure}\hfill
  \begin{subfigure}[t]{0.32\textwidth}
    \centering
    \includegraphics[width=\textwidth]{ch5_fixed_vortex_vector_to_scalar.png}
    \caption{Vector-to-scalar, fixed vortex.}
  \end{subfigure}
  \caption{Average robot and cluster-center trajectories over 10 hardware
  runs per case, shaded to one standard deviation \cite{39}.}
  \label{fig:ch5:hardware_trial}
\end{figure}
% FIGURE SOURCE: trunk/robots_3/Results_for TestBed_Paper/ (Figs. 14, 15,
% 17 of DETC2025-167604): Fixed_vortex_Vector_Sum/average_trajectories.png,
% Sinking_Vortex_Vector_Sum/average_trajectories.png,
% Fixed_vortex_Vector_Sum/Fixed_Vortex_mag_Descent_consistent_Start/average_trajectories.png

\begin{figure}[h]
  \centering
  \begin{subfigure}[t]{0.48\textwidth}
    \centering
    \includegraphics[width=\textwidth]{ch5_sinking_vortex_distance.png}
    \caption{Vector-sum, sinking vortex.}
  \end{subfigure}\hfill
  \begin{subfigure}[t]{0.48\textwidth}
    \centering
    \includegraphics[width=\textwidth]{ch5_vector_to_scalar_distance.png}
    \caption{Vector-to-scalar, fixed vortex.}
  \end{subfigure}
  \caption{Average distance of the cluster center from the field center
  over time, 10 hardware runs per case \cite{39}.}
  \label{fig:ch5:hardware_distance}
\end{figure}
% FIGURE SOURCE: Figs. 16 and 18 of DETC2025-167604; average_distance_vs_time.png
% in the same two folders.

\section{Formation Maintenance Statistics}
\label{sec:ch5:formation_stats}
% Ported from: idetc2025_testbed.txt Sec. 4.3, Table 2 \cite{39}.

Formation shape was held close to its commanded values throughout the
hardware trials of Section~\ref{sec:ch5:hardware_trials}
(Table~\ref{tab:ch5:formation}) \cite{39}, with deviations on the
order of 1~cm in $p$ and $q$, consistent with the resolution limits of
the OptiTrack motion capture system itself (Chapter~\ref{ch:tools}).

\begin{table}[h]
  \centering
  \caption{Cluster formation-holding accuracy across control primitive
  trials, commanded $p=q=0.35$~m, $\beta=1.05$~rad, 10 trial runs per
  column \cite{39}.}
  \label{tab:ch5:formation}
  \begin{tabular}{lccc}
    \toprule
     & $p$ (m) & $q$ (m) & $\beta$ (rad) \\
    \midrule
    Fixed vortex, vector-sum & $0.36 \pm 0.028$ & $0.36 \pm 0.032$ & $1.03 \pm 0.122$ \\
    Fixed vortex, vector-to-scalar & $0.35 \pm 0.015$ & $0.36 \pm 0.009$ & $1.02 \pm 0.045$ \\
    Sinking vortex, vector-sum & $0.35 \pm 0.019$ & $0.35 \pm 0.015$ & $1.02 \pm 0.056$ \\
    \bottomrule
  \end{tabular}
\end{table}

\section{Why Direction Is Not Location}
\label{sec:ch5:why_not_location}
% New text.

Every primitive in this chapter commands a heading. None of them
estimates a position: nothing in Sections~\ref{sec:ch5:vector_sum}
or~\ref{sec:ch5:vector_to_scalar} produces a coordinate for the feature
being sought, only a direction to move in on the current control cycle.
On a purely rotational field component, this is a structural limitation
rather than a tuning problem. A vortex's field vectors are everywhere
tangential to circles about its center, so a primitive that follows the
locally sensed direction, as vector-sum does, tracks that tangential direction with no term correcting
for radial distance from the center. The fixed-vortex hardware result of
Section~\ref{sec:ch5:hardware_trials}, a consistent outward drift for the
vector-sum primitive across all 10 runs and stable at all angles, is this
failure mode observed directly, and on hardware it was worse than
simulation predicted: \cite{39} attributes the drift to the robots
always moving tangent to the circle while momentum effects, not captured
in the simulation model, carry them further outward on each cycle.
Orbiting at a fixed radius is reported there as an open problem the
testbed is positioned to address, not one these primitives solve.

The same qualitative failure appears in the critical-points paper's own
comparison of primitives (Figure~\ref{fig:ch6:primitive_comparison}): both the
vector-sum and vector-to-scalar primitives exhibit outward spiral drift
in a vortex field when compared against a Jacobian-based orbital
controller, because neither has a mechanism to detect or correct radial
error with respect to a specific point. In a saddle field the comparison
is more decisive still: the field vectors do not all point toward the
critical point, so the vector-sum primitive cannot produce an orbital
trajectory at all.


This is why direction-only control cannot certify arrival at a point:
arrival is a statement about position, and a heading command carries no
position information to certify it with. The vector-to-scalar primitive
comes closest, since it converts the problem to scalar-field gradient
ascent or descent and can locate an extremum of the magnitude field, but
it still has no notion of the field's vector structure away from that
extremum, and so cannot classify what kind of critical point it has
found or navigate to a saddle. Recovering that structure requires
estimating the field's local Jacobian, which is what makes a critical
point's location, rather than merely its uphill direction, computable
(Chapter~\ref{ch:first_order}). The same three-robot cluster suffices
for that estimate; the change is that the readings are fitted to an
affine model instead of summed.

\section{Summary}
\label{sec:ch5:summary}
% New text.

This chapter presented two zeroth-order control primitives: vector-sum,
built from the direction of the summed readings, and vector-to-scalar,
built from their magnitudes alone. On 10 hardware trial runs per
environment-primitive combination \cite{39}, vector-sum settled near the
center of a sinking vortex and vector-to-scalar settled near the center
of a fixed vortex, with commanded formation shape held to within about
1~cm. What they cannot do is certify that the cluster has arrived
anywhere: with no position estimate, a rotational field component drives
a steady outward drift that these primitives have no mechanism to
correct, worse on hardware than simulation predicted.
Chapter~\ref{ch:first_order} takes the next step in the ladder of local
polynomial fit order, estimating the field's Jacobian from the same
three robots to recover a critical point's location and type rather than
only its local uphill direction.

% ===================================================================
% CHAPTER 6: FIRST-ORDER PRIMITIVES -- CRITICAL POINTS
% ===================================================================
\chapter{First-Order Primitives: Critical Points}
\label{ch:first_order}

\section{Introduction and Related Work}
\label{sec:ch6:intro}

% New text: chapter framing against Chapter 5.
This chapter is the first rung above the zeroth-order primitives of
Chapter~\ref{ch:zeroth}. At zeroth order, a robot team reduces its
measurements to a single scalar, a mean heading or a flow magnitude, and
steers by that reduction alone. That policy reports a direction, not a
location. It carries no representation of where the field vanishes, so it
cannot certify arrival at a feature, and when applied to station-keeping
around a rotational feature it drifts outward under uncorrected centrifugal
error (Section~\ref{sec:ch6:comparison}). This chapter fits an affine model
to the positions and field measurements of three robots, one Jacobian per
time step, and recovers both the location of the nearby critical point and,
from the same fit, its topological type. The resulting estimate
$\mathbf{p}^*$ supports two control primitives: a proportional law that
drives the cluster onto $\mathbf{p}^*$, and an orbital law that holds a
commanded standoff radius around it.

% Ported from: cwaight_systems_paper.tex Sec. I (contributions list, lines
% 65-67), adapted to chapter framing.
This chapter's contributions are four. It presents a distributed algorithm
that estimates critical point location from three robots' instantaneous
measurements, with no prior map of the field. It proves that three robots
are the minimum number for which this estimation is well posed. It verifies
in simulation that the resulting estimate is sufficient to drive attraction
and orbital navigation across six canonical vector field types. It
validates both primitives in 157 hardware attraction trials and 12 hardware
orbital trials on the testbed of Chapter~\ref{ch:tools}.

% Ported from: cwaight_systems_paper.tex Sec. I, paragraphs 3-6 (lines 53,
% 57-61), trimmed to critical-point-specific material; shared adaptive
% navigation and artificial-vector-field literature moved to Chapter 1
% (\ref{ch:intro}).
By contrast, only a few works address navigation of physically sensed
vector fields \cite{55,39,34}. Work that senses the physical
field falls into two camps. One camp presumes prior knowledge of the
environment, unavailable in the unmapped settings that motivate this
chapter \cite{32,33}. The other operates online in two behaviors:
minimally actuated drifters that ride and station-keep along ambient
streamlines, and PIM-triple-inspired teams that straddle a saddle and track
a single stable or unstable manifold from local velocity measurements
\cite{34,1,37}. Both supply a heading along the flow or track
the manifold of a single assumed saddle; neither recovers the coordinate of
a critical point together with its identity across the full taxonomy of
feature types. Estimating the locations of critical points in vector
fields remains an open problem for multirobot systems.

Recovering those coordinates is harder than following a heading, for
reasons of both information and implementation. The information problem is
that the location and classification of a critical point are absent from
any single velocity sample; they are carried by the spatial derivatives of
the field \cite{56}. A single platform must move, sample again, and
difference its readings, which conflates the spatial gradient with the
field's change in time and can return features that are not physically
present. A distributed team removes this temporal aliasing by sampling at
one instant \cite{39,34,42}, but recovering the derivatives is
still not recovering the feature. Magnitude-based methods \cite{38}
reduce the flow to a scalar and yield a descent direction, but the
magnitude vanishes at every critical point regardless of type, and
estimating the two velocity components as independent scalar fields still
does not locate the point. Both the coordinate and the identity are
properties of the assembled Jacobian: inverting it recovers the location
where the field vanishes, and its eigenvalues determine the type
\cite{2}.

The implementation problem is that the local field must be reconstructed
on physical hardware from noisy distributed measurements. The conditioning
of this reconstruction degrades as the formation approaches collinearity,
so the sampling shape must be held by a formation controller under real
actuator dynamics and communication latency (Chapter~\ref{ch:estimation}).
Both the estimator and the mechatronics of the cluster limit the accuracy
of the recovered features.

\section{Critical Point Estimation from Three Robots}
\label{sec:ch6:estimation}

\subsection{Formation Model and Closed-Form Estimate}
\label{sec:ch6:formation_model}

% Ported from: cwaight_systems_paper.tex Sec. II-B, eqs. 2-9 (lines
% 107-142); thesis.tex "Distributed Critical Point Estimation"
% (lines 1909-1951).
Chapter~\ref{ch:estimation} states the local-polynomial field estimator,
its formation matrix, and its minimality condition for a general fit
order. This section specializes that machinery to fit order one, the
affine case, fit from three robots.

Given three robots sampling a planar vector field with positions
$\mathbf{p}_i = (x_i, y_i)$ and field measurements $(u_i, v_i)$ for
$i \in \{1,2,3\}$, the two field components are approximated over the
cluster's spatial extent by linear planar models
\begin{align}
U(x,y) &= ax + by + c, \label{eq:ch6:affine_u}\\
V(x,y) &= dx + ey + f. \label{eq:ch6:affine_v}
\end{align}
The coefficients solve
\begin{equation}
\mathbf{A}\boldsymbol{\theta}_u = \mathbf{u}, \qquad
\mathbf{A}\boldsymbol{\theta}_v = \mathbf{v},
\label{eq:ch6:coeff_system}
\end{equation}
where
\begin{equation}
\mathbf{A} = \begin{bmatrix} x_1 & y_1 & 1 \\ x_2 & y_2 & 1 \\ x_3 & y_3 & 1 \end{bmatrix},
\quad
\boldsymbol{\theta}_u = \begin{bmatrix} a \\ b \\ c \end{bmatrix},
\quad
\boldsymbol{\theta}_v = \begin{bmatrix} d \\ e \\ f \end{bmatrix},
\label{eq:ch6:formation_matrix}
\end{equation}
and $\mathbf{u} = [u_1,u_2,u_3]^{\mathsf T}$,
$\mathbf{v} = [v_1,v_2,v_3]^{\mathsf T}$ are the measurement vectors.

The critical point $\mathbf{p}^* = (x^*,y^*)$ satisfies $U(x^*,y^*) = 0$
and $V(x^*,y^*)=0$, which in matrix form reads
$\mathbf{J}\mathbf{p}^* = -\mathbf{h}$ with
\begin{equation}
\mathbf{J} = \begin{bmatrix} a & b \\ d & e \end{bmatrix}, \qquad
\mathbf{h} = \begin{bmatrix} c \\ f \end{bmatrix},
\label{eq:ch6:jacobian_offset}
\end{equation}
so that, provided $\det(\mathbf{J}) \neq 0$,
\begin{equation}
\mathbf{p}^* = -\mathbf{J}^{-1}\mathbf{h}.
\label{eq:ch6:pstar}
\end{equation}
The eigenvalues of $\mathbf{J}$ then identify the type of the nearby
critical point (Chapter~\ref{ch:preliminaries}).

If the true field is affine over the cluster's spatial extent,
\eqref{eq:ch6:affine_u}--\eqref{eq:ch6:affine_v} reproduce it without
error and the recovery \eqref{eq:ch6:pstar} is exact at any formation
size. On a smooth nonlinear field, the fit carries a first-order
linearization error that scales with the local curvature of the field;
this error decreases as the formation contracts toward the critical
point.

\subsection{Three Robots Are Minimal and Sufficient}
\label{sec:ch6:minimality}

% Ported from: cwaight_systems_paper.tex Sec. II-C (lines 145-153);
% thesis.tex "Minimum Robot Requirements" (lines 1952-1969).
This is the fit-order-one instance of the general minimality statement of
Chapter~\ref{ch:estimation}. Equations~\eqref{eq:ch6:affine_u}--%
\eqref{eq:ch6:affine_v} each carry three unknown coefficients, and each
robot contributes one equation to each system: robot $i$ supplies
$u_i = ax_i+by_i+c$ for the $U$-component and $v_i = dx_i+ey_i+f$ for the
$V$-component. Determining three unknowns uniquely requires at least
three equations, so three robots are the minimum for critical point
estimation. With fewer than three robots, $\mathbf{A}$ in
\eqref{eq:ch6:coeff_system} has fewer rows than columns, admitting
infinitely many coefficient solutions and preventing unique
identification of $\mathbf{p}^*$. With more than three robots the system
is overdetermined and solved by least squares.

Three non-collinear robots are also sufficient: $\det(\mathbf{A})$ equals
twice the signed area of the formation triangle, so $\mathbf{A}$ is
invertible exactly when the robots are not collinear. This is the
fit-order-one case of the common-curve degeneracy of
Chapter~\ref{ch:estimation}: a formation degenerates when its robots lie
on a common curve of the degree the fit requires, a line at first order.

\section{Critical Point Classification}
\label{sec:ch6:classification}

% Ported from: cwaight_systems_paper.tex Sec. II-A (lines 77-79);
% thesis.tex "Vector Field Environments and Critical Points"
% (lines 1863-1879, Hartman-Grobman caveat). New paragraph: planned
% hardware classification-accuracy analysis.
Chapter~\ref{ch:preliminaries} classifies a critical point by the
eigenvalue structure of the field Jacobian $\mathbf{J}$ evaluated there.
The three-robot estimator of Section~\ref{sec:ch6:formation_model}
recovers $\mathbf{J}$ directly as part of the same fit that locates
$\mathbf{p}^*$, so a type classification is available at the same
measurement cost as the location estimate, with no additional sampling.

This classification holds rigorously for hyperbolic critical points. The
center is the exception: a nonlinear perturbation of a linear center can
produce a spiral, since the Hartman-Grobman theorem does not apply at
$\alpha_{\text{eig}}=0$, and under measurement noise the estimated real
part of a center's eigenvalue hovers near zero, making
center-versus-weak-spiral discrimination ill-posed from noisy data.
Eigenvalue analysis also identifies the degenerate case
$\det(\mathbf{J})=0$, corresponding to a non-hyperbolic critical point or
a region of near-uniform flow; this case is detectable through the same
estimate but is not addressed further in this chapter.

% New analysis, 2026-09-29: experiments/classify_from_hardware_logs.py,
% results in experiments/outputs/classify_hw/summary.json (VF_Robot).
% Numbers below are from that run and await the author's review.
\subsection{Classification on Hardware}
\label{sec:ch6:classification_hw}

The hardware trials of Section~\ref{sec:ch6:attraction_hw} were designed
to test location, but every control cycle also produced a full Jacobian
estimate. Recomputing that estimate offline from the logged robot
positions and sensed hue and saturation reproduces the controller's
logged critical point estimates to within $5\times10^{-7}$~m, so the
Jacobians analyzed here are the ones the controller computed. The logs
hold 78 saddle trials and 80 vortex trials of 101 cycles each (10~s at
10~Hz).
%% TODO(reconcile): the logs contain 80 vortex runs; the critical-points
%% paper reports 79 (157 total), and its bias and precision statistics,
%% quoted elsewhere in this chapter, use that count.

Each estimate was classified by the rules of Chapter~\ref{ch:preliminaries}:
a saddle when $\det\hat{\mathbf{J}} < 0$, a node when the eigenvalues are
real and of one sign, and a rotational type when they form a complex pair
$\alpha_{\text{eig}} \pm i\omega$. Because a measured center never has
$\alpha_{\text{eig}}$ exactly zero, a complex pair was labeled a center
when $|\alpha_{\text{eig}}| < 0.2\,\omega$ and a spiral otherwise; the
threshold is a choice, and results are given for $0.1$ and $0.3$ as well.

\begin{table}[h]
\centering
\caption{Offline classification of the hardware Jacobian estimates. The
terminal window is the final 2~s of each trial; the trial vote is the
majority label over that window.}
\label{tab:ch6:classification_hw}
\begin{tabular}{lcc}
\toprule
 & Saddle map & Vortex map \\
\midrule
Trials & 78 & 80 \\
Class correct (saddle, rotational, node), all cycles & 99.9\% & 99.9\% \\
Class correct, terminal window & 100\% & 100\% \\
Class correct, trial vote & 78/78 & 80/80 \\
Exact type, terminal window, threshold 0.1 & 100\% & 71\% \\
Exact type, terminal window, threshold 0.2 & 100\% & 89\% \\
Exact type, terminal window, threshold 0.3 & 100\% & 99\% \\
\bottomrule
\end{tabular}
\end{table}

The estimate names the class of both fields correctly in effectively
every cycle, including the approach from more than 0.8~m away, so the
affine model's truncation error does not reach the classification in
these fields. Every vortex error is a stable spiral with a small negative
real part; none is a saddle or a node. This is the center-versus-spiral
ambiguity described above, and the terminal window resolves it for most
trials once the threshold exceeds the residual damping in the estimate.

The result carries one condition. Classification, unlike location,
depends on how each printed map's hue directions are registered to the
motion-capture axes: a reflection of the sensed vectors changes the sign
of $\det\hat{\mathbf{J}}$ and exchanges saddles with nodes, while leaving
the estimated critical point where it was. The two maps were printed and
placed under different frame conventions (reflections, rotations, and the
camera's mirroring among them). The results above apply to each map the
registration used in the author's own field reconstructions of
Chapter~\ref{ch:tools}, which negates the sensed $v$ component for the
saddle map only. The controller applied no such registration, since it
used the estimate only for location, and its unregistered Jacobian read
the saddle map as a stable node or spiral throughout the terminal window.
A deployment that reports feature type therefore needs the sensor frame
registered to the navigation frame, a requirement that location alone
does not impose.

\section{Attraction Primitive}
\label{sec:ch6:attraction}

\subsection{Definition}
\label{sec:ch6:attraction_def}

% Ported from: cwaight_systems_paper.tex Sec. II-D.1, eq. 10; thesis.tex
% "Critical Point Attraction" (lines 1979-1993).
The three-robot cluster centroid is
$\mathbf{p}_c = \tfrac{1}{3}\sum_{i=1}^{3}\mathbf{p}_i$. To navigate
$\mathbf{p}_c$ onto the estimated critical point, a proportional control
law is applied:
\begin{equation}
\dot{\mathbf{p}}_c = k(\mathbf{p}^* - \mathbf{p}_c), \qquad k>0.
\label{eq:ch6:attraction}
\end{equation}
The mapping from this cluster-level velocity command to individual robot
velocities is handled by the three-robot SAS cluster-space controller of
Chapter~\ref{ch:estimation} and Appendix~\ref{app:kin3}.

\subsection{Stability}
\label{sec:ch6:attraction_stability}

% Ported from: cwaight_systems_paper.tex Sec. II-D.1 (lines 160-164);
% thesis.tex "Critical Point Attraction" (lines 1987-1992).
\begin{proposition}
\label{thm:ch6:attraction_stability}
Under \eqref{eq:ch6:attraction} and a stationary estimate
($\dot{\mathbf{p}}^*=\mathbf{0}$), the centroid converges globally and
exponentially to $\mathbf{p}^*$ with time constant $\tau = 1/k$.
\end{proposition}
\begin{proof}
Define the error $\mathbf{e} = \mathbf{p}_c - \mathbf{p}^*$. Since
$\dot{\mathbf{p}}^*=\mathbf{0}$, $\dot{\mathbf{e}} = \dot{\mathbf{p}}_c =
-k\mathbf{e}$, so $\mathbf{e}(t) = \mathbf{e}(0)e^{-kt}$, independent of
$\mathbf{e}(0)$.
\end{proof}

On an exactly affine field the estimate $\mathbf{p}^*$ is exact at every
measurement, hence automatically stationary in the noise-free case
(Section~\ref{sec:ch6:formation_model}). Sensor noise and formation
conditioning perturb $\mathbf{p}^*$ between control cycles and are
treated separately in Section~\ref{sec:ch6:error_analysis}; actuator
dynamics and command limits, not part of this kinematic model, are also
addressed there.

\subsection{Simulation Verification}
\label{sec:ch6:attraction_sim}

% Ported from: cwaight_systems_paper.tex Sec. IV, IV-A (lines 239-256),
% Table II (lines 259-280); thesis.tex "Convergence to Critical Points"
% (lines 2074-2119).
Simulation trials used the robot dynamics and simulation environment of
Chapter~\ref{ch:tools}, with a cluster of three robots in an equilateral
triangular formation of side length 0.33~m and attraction gain
$k=1$~s$^{-1}$ in \eqref{eq:ch6:attraction}. Six noise-free analytical
fields, one for each canonical critical point type (Chapter~%
\ref{ch:preliminaries}; field formulas in Appendix~\ref{app:fields}),
were used for trials. After verifying performance in noise-free
environments, trials were also run on two fields reconstructed from
robot-sensed measurements collected on the hardware testbed (vortex and
saddle), incorporating realistic sensor noise.

For each of the eight fields, the cluster was initialized at a random
position drawn uniformly from a 1~m~$\times$~1~m box centered on the true
critical point and run for 100 time steps at 10~Hz. This was repeated
1000 times per field; the final centroid position of each trial gave
systematic bias (mean error) and precision (standard deviation). A trial
was counted as convergent if the trajectory remained bounded and
terminated within the stiction-limited neighborhood of the critical
point.

The cluster converged in 100\% of the 1000 trials for every field tested.
On the six analytical fields, systematic bias remained below 0.001~m and
precision ranged from 0.014 to 0.016~m. On the two reconstructed fields,
systematic bias was 0.028~m for the vortex and 0.005~m for the saddle,
with precision near 0.016~m in both cases (Table~\ref{tab:ch6:sim_attraction}).

\begin{figure}[h]
  \centering
  \includegraphics[width=0.75\textwidth]{ch6_figure_5.png}
  \caption{Successful navigation to critical points from random starting
  positions across six canonical vector field types.}
  \label{fig:ch6:attraction_sim}
\end{figure}
% FIGURE SOURCE: Paper_Writing/Systems Paper/cwaight_systems_paper/figure_5.png

\begin{table}[h]
\footnotesize
\centering
\caption{Simulation results across eight field types (six analytical
plus two reconstructed). Every field achieved 100\% convergence over
1000 trials.}
\label{tab:ch6:sim_attraction}
\setlength{\tabcolsep}{4pt}
\begin{tabular}{lccc}
\toprule
\textbf{Field} & \textbf{Avg.\ Rest Point (x, y) m} & \textbf{Bias (m)} & \textbf{$\sigma$ (m)} \\
\midrule
Sink & (0.0006, $-$0.0001) & 0.0006 & 0.0156 \\
Source & ($-$0.0005, $-$0.0001) & 0.0006 & 0.0140 \\
Vortex & ($-$0.0002, $-$0.0000) & 0.0002 & 0.0163 \\
Sinking Vortex & ($-$0.0002, $-$0.0004) & 0.0004 & 0.0145 \\
Spewing Vortex & (0.0001, 0.0006) & 0.0006 & 0.0157 \\
Saddle & (0.0004, $-$0.0001) & 0.0004 & 0.0157 \\
Vortex (reconstructed) & (0.0206, $-$0.0187) & 0.0278 & 0.0162 \\
Saddle (reconstructed) & ($-$0.0048, $-$0.0007) & 0.0048 & 0.0157 \\
\bottomrule
\end{tabular}
\end{table}

\subsection{Hardware Verification}
\label{sec:ch6:attraction_hw}

% Ported from: cwaight_systems_paper.tex Sec. V-B, V-C (lines 352-403);
% thesis.tex "Attraction Controller Results" (lines 2216-2243).
Hardware validation used the Decabot testbed of Chapter~\ref{ch:tools},
focused on two field types spanning the key challenges of vector field
navigation: rotational dynamics (vortex) and hyperbolic instability
(saddle); the saddle is the only canonical field type here that contains
separatrices. For each field, eight starting locations were selected
around the perimeter of the workspace and the attraction controller
\eqref{eq:ch6:attraction} was run 10 times from each location using the
three-robot equilateral triangle formation with a 0.33~m side length,
yielding 80 trials per field. Two saddle trials and one vortex trial were
lost to file management errors, giving 157 usable trials (79 vortex, 78
saddle).

The cluster converged to the critical point in 100\% of the 157 hardware
trials. In the vortex field, systematic bias was 0.012~m with precision
of 0.009~m; in the saddle field, systematic bias was 0.005~m with
precision of 0.014~m (Table~\ref{tab:ch6:hw_attraction}). The narrow
standard deviation bounds across the 10 repetitions per starting position
confirm that the residual scatter reflects sensor noise rather than
uncontrolled experimental variation.

\begin{figure}[h]
  \centering
  \includegraphics[width=0.75\textwidth]{ch6_allrunsadjusted.png}
  \caption{Convergence trajectories from eight starting positions to the
  center of the vortex (a) and saddle (b) using three robots.}
  \label{fig:ch6:attraction_hw}
\end{figure}
% FIGURE SOURCE: Paper_Writing/Systems Paper/cwaight_systems_paper/allrunsadjusted.png

\begin{figure}[h]
  \centering
  \includegraphics[width=\textwidth]{ch6_fig9.png}
  \caption{Distance to center over time, with mean trajectory and
  standard deviation bounds.}
  \label{fig:ch6:attraction_hw_distance}
\end{figure}
% FIGURE SOURCE: Paper_Writing/Systems Paper/cwaight_systems_paper/fig9.png

\begin{table}[h]
\footnotesize
\centering
\caption{Hardware experimental results. Saddle: 78 runs. Vortex: 79 runs.
Bias measures mean offset from the true critical point; precision
($\sigma$) measures repeatability.}
\label{tab:ch6:hw_attraction}
\setlength{\tabcolsep}{4pt}
\begin{tabular}{lccc}
\toprule
\textbf{Field} & \textbf{Avg.\ Rest Point (x, y) m} & \textbf{Bias (m)} & \textbf{$\sigma$ (m)} \\
\midrule
Saddle (printed map) & (0.004, $-$0.001) & 0.005 & 0.014 \\
Vortex (printed map) & ($-$0.004, 0.012) & 0.012 & 0.009 \\
\bottomrule
\end{tabular}
\end{table}

\section{Orbital Primitive}
\label{sec:ch6:orbital}

\subsection{Definition}
\label{sec:ch6:orbital_def}

% Ported from: cwaight_systems_paper.tex Sec. II-D.2, eq. 11; thesis.tex
% "Orbital Navigation" (lines 1994-2014).
To maintain a circular orbit of radius $r_d$ around $\mathbf{p}^*$,
define the radial vector $\mathbf{r} = \mathbf{p}^* - \mathbf{p}_c$ with
magnitude $r = \lVert\mathbf{r}\rVert$. The unit radial vector is
$\hat{\mathbf{r}} = \mathbf{r}/r$, and the tangent vector
$\hat{\mathbf{r}}_\perp$ is obtained by a $-90^\circ$ rotation. The
orbital velocity command is
\begin{equation}
\dot{\mathbf{p}}_c = k_t\hat{\mathbf{r}}_\perp + k_r(r-r_d)\hat{\mathbf{r}},
\label{eq:ch6:orbital}
\end{equation}
where $k_t$ is the commanded orbital (tangential) speed, whose sign fixes
direction, and $k_r>0$ drives radial convergence.

\subsection{Stability}
\label{sec:ch6:orbital_stability}

% Ported from: cwaight_systems_paper.tex Sec. II-D.2 (lines 173-174);
% thesis.tex "Orbital Navigation" (lines 2007-2013).
\begin{proposition}
\label{thm:ch6:orbital_stability}
Under \eqref{eq:ch6:orbital} and a stationary estimate, the radial error
$e_r = r-r_d$ decays exponentially to zero, and as $r\to r_d$ the angular
velocity stabilizes to $\omega = k_t/r_d$, producing steady circular
motion.
\end{proposition}
\begin{proof}
Equation~\eqref{eq:ch6:orbital} decouples radial and angular dynamics:
the radial error $e_r=r-r_d$ obeys $\dot e_r = -k_r e_r$, so
$e_r(t)=e_r(0)e^{-k_r t}$, ensuring exponential convergence to $r_d$. As
$r\to r_d$, the angular velocity stabilizes to $\omega=k_t/r_d$,
producing steady circular motion.
\end{proof}

\subsection{Simulation Verification}
\label{sec:ch6:orbital_sim}

% Ported from: cwaight_systems_paper.tex Sec. IV-B (lines 283-311),
% Table III; thesis.tex "Orbital Control" (lines 2120-2163).
Using the simulation setup of Section~\ref{sec:ch6:attraction_sim},
orbital control was evaluated on the same eight fields at six commanded
radii with gains $k_t=0.2$~m/s and $k_r=0.5$~s$^{-1}$ in
\eqref{eq:ch6:orbital}. For each radius, the cluster was positioned at
the desired radius to initialize the run, then the orbital controller
was used for 600 time steps; mean radial error and standard deviation
were computed after discarding the first 100 time steps (10~s) as a
settling transient.

On the six noise-free fields, all six produced identical performance at
each radius, with standard deviations below $2\times10^{-5}$~m
(Table~\ref{tab:ch6:sim_orbital}). This is expected: on an exactly
affine noise-free field the three-robot fit recovers $\mathbf{p}^*$
exactly at every step, so the controller input, and hence the
trajectory, is independent of the field type; the residual standard
deviation is the steady-state radius ripple rather than field variation,
hence the single noise-free column in the table. Mean radial errors had
a slight positive bias that decreased with increasing radius, from
0.146~m at a commanded radius of 0.01~m to 0.042~m at 0.50~m. This
pattern held on the reconstructed vortex, with mean errors from 0.150~m
to 0.027~m, and on the reconstructed saddle, from 0.146~m to
$-0.003$~m. In both reconstructed fields, error standard deviations
increased with commanded radius, from 0.012~m to 0.061~m.

No figure in the source paper covers this simulated-orbital experiment;
it is reported there only as Table~\ref{tab:ch6:sim_orbital}.


\begin{table}[h]
\centering
\caption{Simulated orbital control across eight field types. Negative
error indicates the cluster orbited inside the commanded radius.}
\label{tab:ch6:sim_orbital}
\footnotesize
\setlength{\tabcolsep}{3pt}
\begin{tabular}{cccc}
\toprule
\textbf{Radius (m)} & \textbf{Noise-free (6 fields)} & \textbf{Vortex (reconstructed)} & \textbf{Saddle (reconstructed)} \\
 & \textbf{Error (m) $\pm$ std} & \textbf{Error (m) $\pm$ std} & \textbf{Error (m) $\pm$ std} \\
\midrule
0.01 & 0.146 $\pm$ $3{\times}10^{-6}$ & 0.150 $\pm$ 0.015 & 0.146 $\pm$ 0.012 \\
0.10 & 0.109 $\pm$ $7{\times}10^{-6}$ & 0.109 $\pm$ 0.013 & 0.098 $\pm$ 0.015 \\
0.20 & 0.081 $\pm$ $1{\times}10^{-5}$ & 0.104 $\pm$ 0.014 & 0.072 $\pm$ 0.024 \\
0.30 & 0.063 $\pm$ $1{\times}10^{-5}$ & 0.094 $\pm$ 0.018 & 0.059 $\pm$ 0.042 \\
0.40 & 0.050 $\pm$ $2{\times}10^{-5}$ & 0.067 $\pm$ 0.026 & 0.035 $\pm$ 0.054 \\
0.50 & 0.042 $\pm$ $2{\times}10^{-5}$ & 0.027 $\pm$ 0.038 & $-$0.003 $\pm$ 0.061 \\
\bottomrule
\end{tabular}
\end{table}

\subsection{Hardware Verification}
\label{sec:ch6:orbital_hw}

% Ported from: cwaight_systems_paper.tex Sec. V-B, V-D (lines 359-361,
% 407-418), Table V; thesis.tex "Orbital Control Results"
% (lines 2244-2280).
Using the same testbed and field pair as
Section~\ref{sec:ch6:attraction_hw}, the orbiting controller
\eqref{eq:ch6:orbital} was tested at six commanded radii matching the
simulation distances, each run for 3000 cycles at 10~Hz, more than three
complete orbits. Mean radial error and standard deviation were recorded
over the full run, including the initial transient (unlike the
simulation statistics of Table~\ref{tab:ch6:sim_orbital}, which discard a
10~s settling transient).

In the vortex field, mean radial errors decreased from 0.246~m at a
commanded radius of 0.01~m to $-0.035$~m at 0.50~m, with standard
deviations ranging from 0.012 to 0.053~m. In the saddle field, mean
radial errors ranged from 0.114~m to 0.060~m over the same commanded
radii, with standard deviations increasing from 0.011 to 0.095~m
(Table~\ref{tab:ch6:hw_orbital}).

\begin{figure}[h]
  \centering
  \includegraphics[width=\textwidth]{ch6_orbit040_analytical_vs_real.png}
  \caption{Orbital trajectory comparison at a commanded radius of
  0.40~m: analytical prediction vs. hardware performance in the saddle
  field.}
  \label{fig:ch6:orbit_hw}
\end{figure}
% FIGURE SOURCE: Paper_Writing/Systems Paper/cwaight_systems_paper/orbit040_analytical_vs_real.png

\begin{table}[h]
\centering
\caption{Hardware orbital control results. Negative error indicates the
cluster orbited inside the commanded radius.}
\label{tab:ch6:hw_orbital}
\begin{tabular}{ccc}
\toprule
\textbf{Radius (m)} & \textbf{Vortex (printed map)} & \textbf{Saddle (printed map)} \\
 & \textbf{Error (m) $\pm$ std} & \textbf{Error (m) $\pm$ std} \\
\midrule
0.01 & 0.246 $\pm$ 0.053 & 0.114 $\pm$ 0.011 \\
0.10 & 0.190 $\pm$ 0.012 & 0.072 $\pm$ 0.018 \\
0.20 & 0.143 $\pm$ 0.019 & 0.058 $\pm$ 0.037 \\
0.30 & 0.087 $\pm$ 0.015 & 0.089 $\pm$ 0.059 \\
0.40 & 0.032 $\pm$ 0.020 & 0.090 $\pm$ 0.079 \\
0.50 & $-$0.035 $\pm$ 0.023 & 0.060 $\pm$ 0.095 \\
\bottomrule
\end{tabular}
\end{table}

\section{Comparison with Zeroth-Order Primitives}
\label{sec:ch6:comparison}

% Ported from: cwaight_systems_paper.tex Sec. VI-B "Comparison with
% Other Primitives" (lines 462-476); thesis.tex "Comparison with Other
% Primitives" (lines 2328-2345), reframed against Chapter 5.
The modular control architecture of Chapter~\ref{ch:estimation} allows
navigation controllers to be interchanged while the cluster-space and
robot controller layers remain fixed, so performance differences reflect
algorithmic rather than implementation variation. This isolates the
contribution of location estimation: the orbital primitive of
Section~\ref{sec:ch6:orbital} is compared against the vector-sum and
vector-to-scalar primitives of Chapter~\ref{ch:zeroth} in the vortex
field. The vector-sum primitive commands motion in the direction of the
average perceived flow; the vector-to-scalar primitive \cite{38},
modified here to move tangentially to the direction of steepest descent,
reduces the flow to its magnitude.

Both zeroth-order alternatives exhibit outward spiral drift in the
vortex field (Figure~\ref{fig:ch6:primitive_comparison}). Both derive
their velocity commands from direction-only guidance; neither has a
mechanism to detect or correct radial error with respect to a specific
point, so centrifugal drift goes uncorrected. The vector-sum drift has
been confirmed in hardware \cite{39}. In the saddle field the
comparison is more decisive: the field vectors do not all point toward
the critical point, so the vector-sum controller cannot produce an
orbital trajectory at all, and the vector-to-scalar technique would
still drift. The orbital primitive of this chapter avoids both failure
modes because it holds an explicit estimate of $\mathbf{p}^*$ and closes
the loop on the radial error $r-r_d$ directly, rather than acting on
direction alone.

\begin{figure}[h]
  \centering
  \includegraphics[width=\textwidth]{ch6_orbit_demonstration.png}
  \caption{Orbital trajectories: vector-sum primitive versus the
  proposed critical point controller.}
  \label{fig:ch6:primitive_comparison}
\end{figure}
% FIGURE SOURCE: Paper_Writing/Systems Paper/cwaight_systems_paper/orbit_demonstration.png

\section{Error Analysis}
\label{sec:ch6:error_analysis}

% Ported from: thesis.tex "Error Analysis" (lines 2283-2327; primary
% source for the robot-dynamics, equilateral-triangle kappa(A)=7.42, and
% Jacobian-conditioning material); cwaight_systems_paper.tex Discussion
% Sec. VI-A (sigma_p ~ ||J^{-1}|| sigma_v / sqrt(3) noise-propagation
% paragraph, line 454, not carried into thesis.tex's condensed version).
Estimation and control errors arise from four sources: robot dynamics,
sensor noise, cluster geometry, and field conditioning.

Perfect convergence does not occur in simulation. However, repeating
experiments with no stiction, no maximum velocity, and
$\alpha_{\text{mom}}=0$ (instantaneous velocity response,
Chapter~\ref{ch:tools}) confirms convergence to the controller's
$10^{-6}$~m command deadband, with the critical point estimate itself
exact to machine precision. This confirms the residual variance is due
to robot dynamics, not the estimation method.

Positive radial bias at small orbital radii is also a result of robot
dynamics. At small radii, the commanded angular rate $k_t/r_d$ is high
relative to the radial correction bandwidth $k_r$, and the centripetal
demand ($k_t^2/r$) pushes the cluster outward faster than the radial
controller can correct; as the commanded radius increases, this effect
diminishes. More aggressive radial gains would reduce this error, at the
cost of increased sensitivity to noise in the critical point estimate.

Comparing the noise-free and reconstructed fields in simulation isolates
the effect of noisy sensor data: sensor noise is the driving factor of
the simulation-reality gap, with sensing errors propagating directly
through the critical point estimation equations
\eqref{eq:ch6:jacobian_offset}--\eqref{eq:ch6:pstar}. To first order, the
centered equilateral fit maps per-component sensor noise $\sigma_v$ into
estimate scatter
\begin{equation}
\sigma_p \approx \frac{\lVert\mathbf{J}^{-1}\rVert\,\sigma_v}{\sqrt{3}};
\label{eq:ch6:sigma_p}
\end{equation}
the calibration errors of Chapter~\ref{ch:tools} give
$\sigma_v \approx 0.02$~m/s, predicting roughly 0.03~m of scatter on the
printed vortex. This matches the observed 0.009--0.016~m precision after
closed-loop time-averaging, while the 0.012--0.028~m bias reflects
systematic distortion that averaging cannot remove.

To avoid errors from robots becoming collinear (the condition number of
$\mathbf{A}$ approaching infinity), an equilateral triangle was chosen
as the cluster shape: among 200{,}000 random formations of the same size
centered on the sampled region, none beat the equilateral triangle's
$\kappa(\mathbf{A})=7.42$. This optimality is not translation-invariant
(the same formation centered elsewhere has a different condition
number), so it is stated for a formation centered on the sampled region,
as here.

Further errors arise from poor conditioning of the estimated field
Jacobian $\kappa(\mathbf{J})$, directly affecting the critical point
estimate \eqref{eq:ch6:pstar}. The condition number grows when the
Jacobian eigenvalues have very different magnitudes, as occurs near
highly asymmetric nodes or saddle points with disparate stable and
unstable manifold strengths, and diverges as $\det(\mathbf{J})\to0$,
corresponding to non-hyperbolic critical points or near-uniform flow.
The negative bias at the largest commanded radii traces to degradation
of the reconstructed field near the edge of the mapped region: there the
estimated critical point distance is systematically inflated, so the
radial controller pulls the cluster inside the commanded radius, the
same range-dependent error amplification captured by $\kappa(\mathbf{J})$,
since far from the critical point the sampled patch approaches uniform
flow and small field errors produce large estimate shifts. Unlike
$\kappa(\mathbf{A})$, this factor is a property of the field, not the
formation, and is not under the designer's control.

Across all four error sources, estimation errors shift the equilibrium
location of both controllers but preserve the stability established in
Sections~\ref{sec:ch6:attraction_stability} and~\ref{sec:ch6:orbital_stability}.

\section{Summary}
\label{sec:ch6:summary}

% Ported from: thesis.tex "Conclusion" (lines 2388-2397), trimmed to
% this chapter's own results; forward pointers to Chapter 7 and to
% future-work material move to Chapters 9 and 10.
This chapter presented a distributed estimation framework that uses
instantaneous measurements from three robots to recover both the
location and the type of a nearby critical point, and two control
primitives, attraction and orbit, built on that estimate. Simulation
across six canonical field types confirmed convergence of both
primitives (Sections~\ref{sec:ch6:attraction_sim},
\ref{sec:ch6:orbital_sim}). In 157 hardware attraction trials, the
system achieved a 100\% success rate with position errors of 0.012~m or
less (Section~\ref{sec:ch6:attraction_hw}). Orbital control experiments
maintained commanded radii within 0.25~m error at small radii and 0.1~m
at larger radii (Section~\ref{sec:ch6:orbital_hw}). Against the
zeroth-order primitives of Chapter~\ref{ch:zeroth}, the explicit
critical point estimate removes the outward drift both alternatives
exhibit on a rotational feature (Section~\ref{sec:ch6:comparison}).

% ===================================================================
% CHAPTER 7: SECOND-ORDER PRIMITIVES -- COHERENT STRUCTURES
% ===================================================================
\chapter{Second-Order Primitives: Coherent Structures}
\label{ch:second_order}

% Ported from: Separatrix_and_OW_Paper/Draft_11.tex (the Separatrix/
% Okubo-Weiss paper), which supersedes the older port of this material
% in thesis.tex Chapter 5. Draft_11 is treated as the primary source
% throughout this chapter; thesis.tex is used only in the few places
% noted below where it carries description Draft_11 dropped.

\section{Introduction and Related Work}
% New text (opening framing) followed by material ported from
% Draft_11.tex Sec. I (Introduction), trimmed to what is specific to
% coherent structures. Shared adaptive-navigation and scalar-field
% literature is in Chapter 1.

This chapter is the second rung of the fit-order ladder introduced in
Chapter~\ref{ch:intro}. The first rung (Chapter~\ref{ch:first_order}) fits
a plane to each velocity component from three robots, which gives the
velocity gradient at one point. A linear fit gives the velocity gradient
at one point; a second-order fit gives it as a field over the formation,
so any quantity built from it comes with its own gradient in closed form.
A single-point gradient can locate and classify an isolated critical
point, but it carries no information about how the field varies away from
that point, so it cannot find or follow a curve. This chapter fits a full
quadratic to each velocity component from six robots, recovering the
gradient and the Hessian from one synchronized sample, and uses the
result to build two adaptive navigation primitives that acquire and ride
a coherent structure from instantaneous measurements alone.

Formation-based Hessian recovery is established for scalar fields: a ring
of five robots plus one at the center recovers a scalar field's gradient
and Hessian in closed form from one synchronized sample \cite{15}. This
chapter extends that formation to a two-component field, recovering the
Jacobian and its second derivatives together, which is what makes the
surrogate fields of Section~\ref{sec:ch7:surrogates} available in closed
form.

Beyond critical points, vector fields have curves that are useful for
navigation. Dynamic ocean flows contain naturally occurring boundaries
known as coherent structures that influence where floating materials like
debris and larvae will accumulate or separate \cite{10,11}.
Mapping these exact curves allows teams to strategically deploy resources
in highly targeted search areas \cite{11,12,13,14}.
Without adaptive navigation, a large scale team must sample a much larger
basin to determine where these structures lie.

Coherent structures are typically identified through two primary methods.
The first relies on Lagrangian coherent structures (LCS) \cite{47},
typically computed as ridges in the finite-time Lyapunov exponent field
\cite{49}, with ongoing research addressing the Lagrangian/Eulerian
tradeoff and sparse-trajectory gradient recovery
\cite{48,52,50,51}. The second method uses instantaneous
Eulerian criteria, such as the Okubo-Weiss parameter, to partition flows
based solely on the local velocity gradient \cite{43,44}. This
instantaneous approach has been generalized to three dimensions
\cite{45}, refined to second order \cite{46}, and applied at scale
in mesoscale eddy censuses \cite{7}. However, both of these methods
characterize these structures offline using a reconstructed field, which
is not well suited to real-time adaptive navigation of a field that
evolves in time.

New work is being done on using adaptive navigation techniques with
multirobot systems in these scenarios. Michini et al.\ bracket a presumed
stable or unstable manifold with three robots using only local velocity
measurements \cite{34}, an approach inspired by the PIM triple
procedure, extended to $N$ robots \cite{35} and validated
experimentally in gyre-like flows \cite{36}. The robots begin on the
manifold and are advected by the flow, and the manifold's identity is
supplied in advance. The saddle point is treated by neither the
controller nor its analysis, and that work reports veering away from the
structure near a saddle in flow-tank experiments. Kularatne and Hsieh
confirm the boundary by computing the local FTLE field on board, at the
cost of a nonzero integration horizon and a commitment to attracting
structures \cite{1}. In both cases, the robotic formation
pre-straddles the structure, and does not acquire a structure it does not
already occupy.

This chapter shows that a second-order fit of the local flow is enough to
acquire and track such a structure from instantaneous measurements. The
contributions are
\begin{itemize}
\item a cooperative estimator that recovers the local quadratic model of a
two-component field from one synchronized set of measurements across six
robots, the minimum for the unconstrained model, and that fails only when
the robots lie on a common conic;
\item two surrogates read in closed form from that fit, the
velocity-gradient determinant $D$ and the smaller strain eigenvalue $s_1$
\cite{52};
\item two adaptive navigation primitives, the $D$ tracker and the $s_1$
tracker, that navigate to the structure and ride along it without
pre-straddling;
\item a tradeoff between them, in which the $D$ tracker tolerates more
measurement noise while the $s_1$ tracker, being objective, holds its path
under a rotating observer, which matters for a cluster whose heading is
dead-reckoned rather than measured.
\end{itemize}
Both primitives are demonstrated on a simulated double gyre and on
measured Santa Barbara Channel radar currents.

\section{Six-Robot Quadratic Estimation}
\label{sec:ch7:estimation}
% Ported from: Draft_11.tex Sec. II-A (Local Vector Field Polynomials)
% and Sec. II-D (Determinant Field, Strain Eigenvalue Field), specialized
% to the six-robot quadratic case. The general monomial-basis estimator,
% minimality proof, and common-conic degeneracy result are in Chapter 4
% (Sec. 4.3); only the second-order-specific closed forms are repeated
%% NOTE(new-math, checked 2026-09-29): eq:ch7:grad_D is not written in Draft_11.
%% Checked against eq:ch7:Jhat (u_xx=a5, u_xy=a4, u_yy=a6): both rows match.
% here. The gradient of $D$ (eq. \ref{eq:ch7:grad_D}) is a new derivative
% of Draft_11's eq. \ref{eq:ch7:Jhat}, written out because Draft_11 never
% states it explicitly (it is used only implicitly through the
% orthogonal-gradient argument of Sec.~\ref{sec:ch7:d_definition}); the
% same closed form appears, in $u_x, u_y, \ldots$ notation, in
% thesis.tex's "Closed-Form Recovery of $D$, $\nabla D$, $\mathbf{H}_D$."

Chapter~\ref{ch:estimation} states the general local-polynomial
estimator, proves that six robots are necessary and (generically)
sufficient for the unconstrained planar quadratic model, and
characterizes the resulting degeneracy as six points lying on a common
conic. This section specializes that estimator to give the closed-form
quantities the two trackers of this chapter consume.

A planar vector field assigns a velocity $\mathbf{v}(\mathbf{p}) =
(u(\mathbf{p}), v(\mathbf{p}))$ to each position $\mathbf{p} = (x, y)$.
Coordinates are measured from the formation centroid $\mathbf{p}_c$
throughout this chapter, so $\mathbf{p}_c = \mathbf{0}$. The second-order
Taylor expansion about the centroid is
\begin{align}
    u(\mathbf{p}) &= a_1 + a_2 x + a_3 y
        \nonumber \\
        &\quad + a_4 xy + \tfrac{1}{2} a_5 x^2 + \tfrac{1}{2} a_6 y^2
        + \mathcal{O}(\|\mathbf{p}\|^3),
        \label{eq:ch7:quad_model_u} \\
    v(\mathbf{p}) &= b_1 + b_2 x + b_3 y
        \nonumber \\
        &\quad + b_4 xy + \tfrac{1}{2} b_5 x^2 + \tfrac{1}{2} b_6 y^2
        + \mathcal{O}(\|\mathbf{p}\|^3),
        \label{eq:ch7:quad_model_v}
\end{align}
where $\mathcal{O}(\|\mathbf{p}\|^3)$ is the truncation error. The twelve
coefficients are constants, each a derivative of the field at the
centroid: $a_1 = u(\mathbf{p}_c)$, $a_2 = \partial u/\partial
x|_{\mathbf{p}_c}$, $a_3 = \partial u/\partial y|_{\mathbf{p}_c}$, $a_4 =
\partial^2 u/\partial x \partial y|_{\mathbf{p}_c}$, $a_5 = \partial^2
u/\partial x^2|_{\mathbf{p}_c}$, $a_6 = \partial^2 u/\partial
y^2|_{\mathbf{p}_c}$, and likewise for $b_1, \ldots, b_6$ with $v$ in
place of $u$. The one-half factors cancel the factor of 2 that
differentiation introduces. The zeroth-order coefficients give the fitted
flow at the centroid, $\hat{\mathbf{v}}_0 = (\hat{a}_1, \hat{b}_1)$, used
by both trackers below as the measured local velocity.

With robot positions $\mathbf{p}_i = (x_i, y_i)$, define the quadratic
basis vector
\begin{equation}
    \boldsymbol{\phi}(\mathbf{p}) =
    \bigl[\,1, \; x, \; y, \; xy, \;
    \tfrac{1}{2}x^2, \; \tfrac{1}{2}y^2\,\bigr]^\top,
\end{equation}
whose entries are the terms of
(\ref{eq:ch7:quad_model_u})--(\ref{eq:ch7:quad_model_v}) in order.
Stacking the six basis vectors row-wise gives the $6 \times 6$ formation
matrix $\boldsymbol{\Phi}$, and the coefficient estimates
$\hat{\boldsymbol{\theta}}_u = [\hat{a}_1, \ldots, \hat{a}_6]^\top$ and
$\hat{\boldsymbol{\theta}}_v = [\hat{b}_1, \ldots, \hat{b}_6]^\top$ solve
\begin{equation}
    \boldsymbol{\Phi}\,\hat{\boldsymbol{\theta}}_u = \mathbf{m}_u, \qquad
    \boldsymbol{\Phi}\,\hat{\boldsymbol{\theta}}_v = \mathbf{m}_v,
    \label{eq:ch7:solve}
\end{equation}
where $\mathbf{m}_u = [u_1, \ldots, u_6]^\top$ and $\mathbf{m}_v = [v_1,
\ldots, v_6]^\top$ collect the velocity components sampled by robot $i$.
Six robots suffice provided $\boldsymbol{\Phi}$ is nonsingular, which
fails exactly when the robots lie on a common conic (Chapter
\ref{ch:estimation}); a standard six-robot ring is degenerate at every
radius, since the ring itself is the conic.

For this study we chose the pentagon-plus-center formation and use it
throughout. It is the smallest uniform ring plus center on which
\cite{15} recovers a scalar Hessian exactly for quadratic signals,
independent of the formation's orientation. The same symmetry makes the
noise gains of the fit here independent of formation heading
(Chapter~\ref{ch:estimation}).

\paragraph{Closed-form recovery of $D$, $\nabla D$, and $\mathbf{H}_D$.}
The fitted coefficients define the velocity gradient at every point of
the footprint. Differentiating the quadratic model gives a Jacobian whose
entries are affine in position,
\begin{equation}
    \hat{\mathbf{J}}(\mathbf{p}) =
    \begin{bmatrix}
        \hat{a}_2 + \hat{a}_5 x + \hat{a}_4 y &
        \hat{a}_3 + \hat{a}_4 x + \hat{a}_6 y \\
        \hat{b}_2 + \hat{b}_5 x + \hat{b}_4 y &
        \hat{b}_3 + \hat{b}_4 x + \hat{b}_6 y
    \end{bmatrix}.
    \label{eq:ch7:Jhat}
\end{equation}
The determinant $\hat{D}(\mathbf{p}) = \det\hat{\mathbf{J}}(\mathbf{p})$
is therefore exactly quadratic over the footprint, and its value,
gradient, and Hessian at the centroid are polynomials in the
coefficients. The value is $\hat{D}_0 = \hat{a}_2\hat{b}_3 -
\hat{a}_3\hat{b}_2$. Differentiating (\ref{eq:ch7:Jhat}) once more gives
the gradient,
\begin{equation}
    \hat{\mathbf{g}}_{D,0} = \nabla\hat{D}|_0 =
    \begin{bmatrix}
        \hat{a}_5\hat{b}_3 + \hat{a}_2\hat{b}_4
            - \hat{a}_4\hat{b}_2 - \hat{a}_3\hat{b}_5 \\
        \hat{a}_4\hat{b}_3 + \hat{a}_2\hat{b}_6
            - \hat{a}_6\hat{b}_2 - \hat{a}_3\hat{b}_4
    \end{bmatrix},
    \label{eq:ch7:grad_D}
\end{equation}
and the Hessian is
\begin{equation}
    \hat{\mathbf{H}}_{D,0} =
    \begin{bmatrix}
        2(\hat{a}_5\hat{b}_4 - \hat{a}_4\hat{b}_5) &
        \hat{a}_5\hat{b}_6 - \hat{a}_6\hat{b}_5 \\
        \hat{a}_5\hat{b}_6 - \hat{a}_6\hat{b}_5 &
        2(\hat{a}_4\hat{b}_6 - \hat{a}_6\hat{b}_4)
    \end{bmatrix},
    \label{eq:ch7:hess_det}
\end{equation}
which is constant over the footprint, since it is the Hessian of the
fitted $\hat{D}$ rather than of the true $D$. The Hessian of the true $D$
also carries products of $\mathbf{J}$ with third derivatives of the
field, which a quadratic fit cannot recover at any radius.

\paragraph{Closed-form recovery of the strain quantities.} The second
surrogate comes from the rate-of-strain tensor $\hat{\mathbf{S}} =
\tfrac{1}{2}(\hat{\mathbf{J}} + \hat{\mathbf{J}}^\top)$, the symmetric
part of the same fitted Jacobian. At the centroid its mean, deviatoric
normal, and shear strain are $\hat{\mu} = (\hat{a}_2 + \hat{b}_3)/2$,
$\hat{s}_n = (\hat{a}_2 - \hat{b}_3)/2$, and $\hat{s}_s = (\hat{a}_3 +
\hat{b}_2)/2$, and each is affine in position. The eigenvalues of
$\hat{\mathbf{S}}$ are $\hat{s}_{1,2} = \hat{\mu} \mp \hat{r}$ with
$\hat{r} = \sqrt{\hat{s}_n^2 + \hat{s}_s^2}$, and their eigenvectors
$\mathbf{e}_1$ and $\mathbf{e}_2$ are the compression and stretching
directions. The gradient of the smaller eigenvalue follows by the chain
rule,
\begin{equation}
    \nabla\hat{s}_1 = \nabla\hat{\mu}
        - \frac{\hat{s}_n\nabla\hat{s}_n + \hat{s}_s\nabla\hat{s}_s}{\hat{r}},
    \label{eq:ch7:grad_s1}
\end{equation}
with $\nabla\hat{\mu}$, $\nabla\hat{s}_n$, and $\nabla\hat{s}_s$ read from
the second-order coefficients. Unlike $\hat{D}$, $\hat{s}_1$ has no usable
Hessian: it is an affine function minus the norm of an affine map, so it
is concave over the footprint for every field and formation.

\section{Eulerian Surrogates}
\label{sec:ch7:surrogates}
% Ported from: Draft_11.tex Sec. II-D (Eulerian Surrogates), the
% argument for why D and s1 mark the separatrix. The closed-form
% equations these subsections reference are given in
% Sec.~\ref{sec:ch7:estimation} above rather than repeated here.

Lagrangian methods locate coherent structures by integrating particle
trajectories over a finite time window \cite{47,49}. This chapter
instead reads them from the instantaneous geometry of the fitted velocity
field, as Eulerian coherent structures \cite{52}. Two fields computed
from the six-robot fit mark the separatrix: the determinant of the
velocity gradient, and the smaller eigenvalue of its symmetric part.

\subsection{Determinant Field}

For incompressible flow the determinant $\hat{D}$ of
(\ref{eq:ch7:Jhat})--(\ref{eq:ch7:hess_det}) is a negative multiple of the
Okubo-Weiss parameter \cite{43,44}, so the two carry the same
partition of the flow. Where $D < 0$ the flow is strain-dominated, and
where $D > 0$ it is rotation-dominated. In the flows studied here,
separatrices run along the trenches of $D$ inside the strain-dominated
regions.

\subsection{Strain Eigenvalue Field}

Serra and Haller define objective Eulerian coherent structures (OECS)
from the eigenvalues and eigenvectors of $\hat{\mathbf{S}}$
\cite{52}. An attracting OECS is a trench of $s_1$ tangent to
$\mathbf{e}_2$, and floating material has been observed at sea to collect
onto such curves \cite{11}. A repelling OECS is a ridge of $s_2$
tangent to $\mathbf{e}_1$, along which neighboring material separates. In
incompressible flow $s_2 = -s_1$, so repelling structures are also
trenches of $s_1$, and one tracker rides both.

\subsection{Relation Between the Surrogates}

Both fields are read from the same fitted Jacobian, so moving between
them changes nothing in the estimator. Where the vorticity $\omega =
\partial v/\partial x - \partial u/\partial y$ vanishes, $\mathbf{J}$
equals $\mathbf{S}$ and $D = s_1 s_2$, which is $-s_1^2$ in incompressible
flow. Along such a curve the two fields share their trenches. Elsewhere
the spin $\mathbf{W} = \mathbf{J} - \mathbf{S}$, the antisymmetric part of
$\mathbf{J}$ that carries the vorticity, separates them. It also
separates them across observers. A frame rotating clockwise at rate
$\Omega$, with rotation $\mathbf{Q}$, adds a uniform spin that raises
$\omega$ to $\omega + 2\Omega$. The strain tensor only co-rotates, so
$s_1$ is objective \cite{52}, while the determinant keeps the spin,
\begin{equation}
    D' = D + \Omega\,\omega + \Omega^2 .
    \label{eq:ch7:D_not_objective}
\end{equation}
Relatively rotating observers therefore disagree about where the trenches
of $D$ lie. The double-gyre example of Chapter~\ref{ch:preliminaries}
gives both fields in closed form and shows where their trenches coincide
and where they part.

\section{$D$ Tracker}
\label{sec:ch7:d_tracker}

\subsection{Definition}
\label{sec:ch7:d_definition}
% Ported from: Draft_11.tex Sec. III-C (The D Tracker). The Table I
% operating points (Draft_11 Sec. III-A) are reproduced here since both
% trackers draw their parameters from it; the surrounding three-layer
% control architecture description is in Chapter 4.

Both trackers of this chapter share the operating point of
Table~\ref{tab:ch7:params}, fixed across all experiments in this chapter
unless noted.

\begin{table}[!t]
\caption{Operating points. The double-gyre column is non-dimensional,
with one time unit equal to one second of the Decabot's identified
dynamics. Section~\ref{sec:ch7:sb_setup} converts the ocean column to
physical units.}
\label{tab:ch7:params}
\centering
\begin{tabular}{llll}
\hline
Parameter & Double gyre & Ocean & Units \\
\hline
$c_{\max}$              & 0.04  & 0.033  & length/time \\
$k$                     & 3.0   & 1.94   & none \\
$\rho$                  & 0.075 & 0.098  & length (5.0 km) \\
$v_{\text{max,robot}}$  & 0.3   & 0.3    & length/time \\
$v_{\text{stiction}}$   & 0.025 & 0.002  & length/time \\
$\alpha_{\text{mom}}$   & 0.7   & 0      & none \\
$\varepsilon_{\text{raw}}$ & $10^{-3}$ & $10^{-3}$ & $1/\text{time}^2$ \\
$\varepsilon_{\dim}$    & 0.025 & 0.025  & length$^2$ \\
$g_\perp$               & 1.0   & 1.0    & length$^2$ \\
$s_{\mathrm{trim}}$     & 0.05  & 0.3    & 1/time \\
$g_{\text{capture}}$    & 0.15  & 0.15   & 1/(time$\cdot$length) \\
$r_{\text{band}}$       & 0.05  & 0.05   & 1/time \\
\hline
\end{tabular}
\end{table}

On the double-gyre benchmark of Chapter~\ref{ch:preliminaries}, the $D$
landscape around the separatrix is a mountain pass. Two minima of $D$ sit
at the flow saddles, and the separatrix is the path of steepest descent
to each from a saddle point of $D$ at the origin, the crest. Around the
crest the landscape is a hyperbolic paraboloid, curving up across the
separatrix and down along it. The $D$ tracker drives the centroid onto
this trench and along it, over the crest, and down toward the next
minimum.

A trench is taken here in the height-ridge sense \cite{57}, a curve on
which the gradient is orthogonal to a Hessian eigenvector with positive
eigenvalue. The tracker reads this frame from $\hat{\mathbf{H}}_{D,0}
\mathbf{w}_j = \lambda_j \mathbf{w}_j$, with $\|\mathbf{w}_j\| = 1$,
$\lambda_1 \leq \lambda_2$, $\mathbf{w}_2$ across the trench, and
$\mathbf{w}_1$ along it, signed by $\hat{\mathbf{v}}_0^\top \mathbf{w}_1
\geq 0$ to point downstream. The identification is exact on the
double-gyre benchmark, where $\mathbf{H}_D$ is diagonal, and approximate
in general. The primitive assumes the true surface has the shape of the
crest, convex across and concave along; where that assumption fails
(Chapter~\ref{ch:preliminaries}) the fit does not register the change.

A full Newton step $-\hat{\mathbf{H}}_{D,0}^{-1}\hat{\mathbf{g}}_{D,0}$
leads to the stationary point of the fitted hyperbolic paraboloid, which
would stop the cluster at the crest. The tracker keeps only its
across-trench component and sets the along-trench speed separately, from
the along-trench gradient away from the crest and from the measured flow
near it, where that gradient vanishes. The flow takes over inside a
shallow-determinant band $\mathcal{B}$, the set of points where
\begin{equation}
    |\hat{D}_0| < \varepsilon_{\mathrm{raw}}
    \quad\text{or}\quad
    \frac{|\hat{D}_0|}{\|\hat{\mathbf{H}}_{D,0}\|_F}
        < \varepsilon_{\dim}.
    \label{eq:ch7:band}
\end{equation}
The second test measures depth against curvature, in units of length
squared, so the band scales with the landscape. The band surrounds the
zero set of $D$, which the separatrix meets at the crest and the cluster
may also cross during acquisition, where a trigger changes only the
along-trench speed. The along-trench speed is
\begin{equation}
    v_\parallel =
    \begin{cases}
        \hat{\mathbf{v}}_0^\top \mathbf{w}_1,
            & \mathbf{p}_c \in \mathcal{B}, \\[6pt]
        \dfrac{|\mathbf{w}_1^\top \hat{\mathbf{g}}_{D,0}|}{|\lambda_1|},
            & \text{otherwise},
    \end{cases}
    \label{eq:ch7:vpar}
\end{equation}
and the command, scaled by a navigation gain $k$, with each channel
saturated by $\mathrm{sat}(c) = c_{\max}\tanh(c/c_{\max})$, is
\begin{equation}
    \dot{\mathbf{p}}_c =
    \underbrace{\mathrm{sat}(v_\parallel)\,\mathbf{w}_1}_{\text{along the trench}}
    + \underbrace{\mathrm{sat}\!\left(-\frac{\mathbf{w}_2^\top
      \hat{\mathbf{g}}_{D,0}}{\lambda_2}\right)\mathbf{w}_2}_{\text{across the trench}}.
    \label{eq:ch7:d_law}
\end{equation}
The saturation caps the centroid speed at $\sqrt{2}\,k\,c_{\max}$. Both
branches of (\ref{eq:ch7:vpar}) are nonnegative, so the along-trench
motion never opposes the flow.

When $\hat{\mathbf{H}}_{D,0}$ is definite the trench frame is undefined.
The tracker then reverts to the signed Newton step
$\sum_j \mathrm{sat}(-\mathbf{w}_j^\top\hat{\mathbf{g}}_{D,0}/
\lambda_j)\mathbf{w}_j$ off the band and to the saturated flow on it. On a
positive-definite fit that step settles at the fitted minimum, so the
same sign test is the capture condition,
\begin{equation}
    \lambda_1 \lambda_2 \geq 0,
    \label{eq:ch7:d_capture}
\end{equation}
which the true $D$ meets at a flow saddle, where it has a local minimum.
A negative-definite fit carries no stability claim. On the benchmark the
test fires only near the flow saddles and domain corners, where second
derivatives vanish and truncation error sets the sign of the fitted
eigenvalues.

\subsection{Transverse Stability}
\label{sec:ch7:d_stability}
% Ported from: Draft_11.tex Sec. III-C (closing paragraph) and the
% Appendix ("Stability of the Trackers"), statement only. The proof is
% ported to Appendix~\ref{app:stability} by another pass.

The transverse channel of (\ref{eq:ch7:d_law}) reduces to $\dot{n} =
-k\,\mathrm{sat}(a_\perp n)$, where $n$ is the distance across the trench
and $a_\perp = \kappa_\perp/\hat{\kappa}_\perp$ is the ratio of true to
fitted transverse curvature. For $a_\perp > 0$ the cluster converges to
the trench, exponentially once $|n| < c_{\max}/a_\perp$, and estimation
error bounded by $\delta$ leaves it practically stable at $\limsup |n|
\leq \delta/a_\perp$. Since $v_\parallel \geq 0$, traversal never
reverses.

\begin{proposition}
\label{prop:ch7:transverse}
Let $\dot{n} = -k\,\mathrm{sat}(a_\perp n + e)$ with $a_\perp > 0$ and
$|e| \leq \delta$. Then $\limsup_{t\to\infty}|n| \leq \delta/a_\perp$. If
$e = 0$, $|n|$ falls below $c_{\max}/a_\perp$ in finite time, then decays
at rate at least $k\,a_\perp\tanh 1$.
\end{proposition}

For the $D$ tracker, writing the fitted transverse gradient as
$\hat\kappa_\perp n$ and the true one as $\kappa_\perp n + \epsilon$ with
$|\epsilon| \leq \delta_g$ gives $a_\perp = \kappa_\perp/\hat\kappa_\perp$
and $e = \epsilon/\hat\kappa_\perp$ in the proposition above, so the $D$
tracker additionally requires $\hat\kappa_\perp > 0$, that is, the fitted
Hessian must agree in sign with the true one across the trench. The full
derivation and proof are in Appendix~\ref{app:stability}.

\subsection{Simulation Verification}
\label{sec:ch7:d_verification}
% Ported from: Draft_11.tex Sec. IV-B (Clean Runs), the D-tracker-
% specific terminal-step result. The shared setup and acquisition/
% traversal narrative is in Section~\ref{sec:ch7:dg_bench}.

\begin{figure}[h]
  \centering
  \includegraphics[width=0.75\textwidth]{ch7_traverse_vs_logic_c.png}
  % FIGURE SOURCE: Paper_Writing/Separatrix_and_OW_Paper/figures/traverse_vs_logic_c.png
  \caption{$s_1$ tracker versus $D$ tracker from six matched starts
  ($\bullet$), noise-free. Both traverse the same network through the
  isotropic point at the origin. The $s_1$ tracker captures at the far
  saddle ($\times$), while the $D$ tracker continues onto the crossing
  wall trench past each saddle it reaches, ending at the square markers.}
  \label{fig:ch7:d_tracker_sim}
\end{figure}

The $D$ tracker passes within 0.02 of $\mathbf{p}^*_1 = (0, -0.5)$ in all
six runs of Fig.~\ref{fig:ch7:d_tracker_sim} and continues onto the
crossing wall trench. Near the saddle the fitted curvature is set by
truncation, so (\ref{eq:ch7:d_capture}) fires only intermittently
(Section~\ref{sec:ch7:d_definition}).

\section{$s_1$ Tracker}
\label{sec:ch7:s1_tracker}

\subsection{Definition}
\label{sec:ch7:s1_definition}
% Ported from: Draft_11.tex Sec. III-D (The s_1 Tracker).

The $s_1$ landscape has the same two minima as the $D$ landscape, and
between them the separatrix is again a trench. Where the $D$ trench
crosses a mountain pass, the $s_1$ trench crosses a sharp fold. In
incompressible flow $s_1 = -r \leq 0$, so the trench rises at most to
$s_1 = 0$, where the strain vanishes. Since $D = -s_1^2$ along the
separatrix (Section~\ref{sec:ch7:surrogates}), squaring rounds the fold
into the smooth crest of $D$.

The fitted $\hat{s}_1$ is concave for every field and formation
(Section~\ref{sec:ch7:estimation}), so unlike $\hat{D}$ it never shows the
upward curvature a trench has across it, and no Newton step is available.
Both channels are built instead from $\nabla\hat{s}_1$ and the strain
eigenframe, with $\hat{r}$ measuring how well resolved that eigenframe
is.

The trench tangent is a strain eigenvector, $\mathbf{e}_2$ on an
attracting segment and $\mathbf{e}_1$ on a repelling one, so no fixed
choice works across the isotropic point of $\mathbf{S}$. On the trench
the transverse derivative of $s_1$ vanishes, so away from a stationary
point $\nabla\hat{s}_1$ lies almost along the trench, and the tangent is
the eigenvector it is most nearly parallel to, signed against a
reference,
\begin{equation}
    \mathbf{t}_{\mathrm{raw}} =
    \arg\max_{\mathbf{e} \in \{\mathbf{e}_1, \mathbf{e}_2\}}
    \bigl\lvert \nabla\hat{s}_1^\top \mathbf{e} \bigr\rvert,
    \qquad
    \mathbf{t} = \mathrm{sgn}\bigl(\mathbf{t}_{\mathrm{ref}}^\top
    \mathbf{t}_{\mathrm{raw}}\bigr)\,\mathbf{t}_{\mathrm{raw}},
    \label{eq:ch7:tangent_select}
\end{equation}
where $\mathbf{t}_{\mathrm{ref}}$ is the previous step's tangent, or
$\hat{\mathbf{v}}_0$ on the first ride step, since the measured flow is
not objective (Section~\ref{sec:ch7:rotating}). The rule swaps from
$\mathbf{e}_2$ to $\mathbf{e}_1$ at the isotropic point with no dedicated
test.

\begin{figure}[h]
  \centering
  \includegraphics[width=\textwidth]{ch7_s1_channels.png}
  % FIGURE SOURCE: Paper_Writing/Separatrix_and_OW_Paper/figures/s1_channels.png
  \caption{The two channels of the $s_1$ tracker, on the benchmark trench.
  (a) Tangent selection, (\ref{eq:ch7:tangent_select}). At each sample the
  strain eigenframe is drawn, the selected eigenvector heavy and the
  rejected one light, with $\nabla\hat{s}_1$ in orange. The argmax takes
  $\mathbf{e}_2$ on the attracting segment and $\mathbf{e}_1$ on the
  repelling segment, swapping through the isotropic point. (b) The
  command. Travel along $\mathbf{t}$ is open loop at cruise speed, while
  $\mathbf{P}$ removes the along-trench component of the gradient so
  descent acts only across the trench.}
  \label{fig:ch7:s1_channels}
\end{figure}

Until $\hat{s}_1$ first falls below $-s_{\mathrm{trim}}$ the command is
gradient descent, $-\mathrm{sat}(g_\perp\nabla\hat{s}_1)$, with transverse
gain $g_\perp$. From then on the tracker rides,
\begin{equation}
    \dot{\mathbf{p}}_c =
    c_{\max}\tanh(1)\,\mathbf{t}
    - \mathrm{sat}\!\bigl(g_\perp\,\mathbf{P}\,\nabla\hat{s}_1\bigr),
    \qquad
    \mathbf{P} = \mathbf{I} - \mathbf{t}\mathbf{t}^\top,
    \label{eq:ch7:s1_law}
\end{equation}
scaled by $k$ as in (\ref{eq:ch7:d_law}), where $\mathrm{sat}$ saturates a
vector's magnitude and preserves its direction (Fig.
\ref{fig:ch7:s1_channels}). The projector removes the along-trench
component of the gradient, so the cluster descends only across the trench
while riding along it at speed $k\,c_{\max}\tanh 1$. The latch depth
$s_{\mathrm{trim}}$ is a small fraction of the benchmark's well depth
$-\pi^2 A \approx -0.987$ (Chapter~\ref{ch:preliminaries}), so from most
starts the ride begins at the first step and projected descent carries
the cluster onto the trench. Below a strain floor $r_{\mathrm{band}}$ the
eigenframe is unresolved. If $\hat{r} < r_{\mathrm{band}}$ before any
tangent exists, (\ref{eq:ch7:tangent_select}) has no direction, and the
command rides $\hat{\mathbf{v}}_0/\|\hat{\mathbf{v}}_0\|$ with $\mathbf{P}
= \mathbf{I}$. Only the first ride step can reach this case.

The tracker also stops at a well. Its capture test is built on the
gradient, since $\hat{\mathbf{H}}_{s_1}$ is negative semidefinite for
every field and formation (Section~\ref{sec:ch7:estimation}) and cannot
certify a minimum. The fitted gradient tracks the true gradient even
where the fitted curvature carries the wrong sign, and a flow saddle is a
true two-dimensional minimum of $s_1$, so a vanishing gradient suffices.
When $\hat{r} > r_{\mathrm{band}}$, $\|\nabla\hat{s}_1\| <
g_{\mathrm{capture}}$ for a capture threshold $g_{\mathrm{capture}}$
(Table~\ref{tab:ch7:params}), and $\hat{s}_1 < -4s_{\mathrm{trim}}$, the
tracker drops the ride term and descends $\hat{s}_1$, holding the cluster
at the minimum until $\hat{s}_1$ rises back above $-s_{\mathrm{trim}}$.

\subsection{Transverse Stability}
\label{sec:ch7:s1_stability}
% Ported from: Draft_11.tex Sec. III-D (closing paragraph) and the
% Appendix ("Stability of the Trackers"), statement only. The proof is
% ported to Appendix~\ref{app:stability} by another pass.

The transverse channel of (\ref{eq:ch7:s1_law}) reduces to the same form
as Proposition~\ref{prop:ch7:transverse}, with $a_\perp =
g_\perp\kappa_\perp$ set by the field rather than normalized by a fitted
curvature, and $e = g_\perp\epsilon$. Both bounds equal
$\delta_g/\kappa_\perp$, independent of gain and fitted curvature, but
only the $D$ tracker needs $\hat\kappa_\perp > 0$. Where $\kappa_\perp$
vanishes over a length, as at the fold, the ride carries the cluster
across (Appendix~\ref{app:stability}). A frame rotating at $\Omega$
perturbs $\dot n$ by at most $\Omega\lVert\mathbf{p}\rVert$, leaving $n$
bounded while $\Omega\lVert\mathbf{p}\rVert < k\,c_{\max}$
(Section~\ref{sec:ch7:rotating}). The ride direction is seeded once from
the measured flow and is guaranteed only under exact estimates.

\subsection{Simulation Verification}
\label{sec:ch7:s1_verification}
% Ported from: Draft_11.tex Sec. IV-B (Clean Runs), the s1-tracker-
% specific terminal-step result. The shared setup and acquisition/
% traversal narrative is in Section~\ref{sec:ch7:dg_bench}.

\begin{figure}[h]
  \centering
  \includegraphics[width=0.75\textwidth]{ch7_separatrix_formation.png}
  % FIGURE SOURCE: Paper_Writing/Separatrix_and_OW_Paper/figures/separatrix_formation.png
  \caption{The six robots from start S1 (star) under both trackers,
  zoomed to the separatrix $x = 0$ (dashed). Thin lines are robot paths
  and the heavy line is the centroid. Outlines mark the
  pentagon-plus-center formation at the start, during the ride, and at
  the far saddle ($\times$).}
  \label{fig:ch7:s1_tracker_sim}
\end{figure}

The $s_1$ tracker settles within 0.011 to 0.022 of $\mathbf{p}^*_1 = (0,
-0.5)$ in all six runs, capturing at the far saddle rather than
continuing onto the crossing wall trench, as shown for start S1 in
Fig.~\ref{fig:ch7:s1_tracker_sim}.

\section{Behavior on the Double-Gyre Benchmark}
\label{sec:ch7:dg_bench}
% Ported from: Draft_11.tex Sec. IV-A (Setup) and IV-B (Clean Runs).
% Renamed from the skeleton's "Estimator Accuracy on the Double Gyre";
% the content Draft_11 places here is acquisition and traversal
% behavior under clean measurements, not an open-loop estimator-error
% sweep (estimator_accuracy_vs_noise.png discarded by the author).

\subsection{Setup}

The double-gyre experiments use the steady field of
Chapter~\ref{ch:preliminaries} ($A = 0.1$) on $[-1, 1] \times [-0.5,
0.5]$. All trials use the nominal formation with ring radius $\rho =
0.075$; trials vary only the initial heading, drawn uniformly from $[0,
2\pi)$ in the Monte Carlo sweeps of Sections~\ref{sec:ch7:noise} and
\ref{sec:ch7:rotating}, and zero in the noise-free runs below.

\subsection{Clean Runs}

These runs test acquisition, traversal, and the terminal step for both
primitives. On the double gyre the two fields share every trench
(Chapter~\ref{ch:preliminaries}). Clean trials run 400 to 600 steps and
stop at domain exit or, for the $D$ tracker, 150 steps after first saddle
contact.

Both trackers ran noise-free from six matched starts
(Fig.~\ref{fig:ch7:d_tracker_sim}), spanning both gyre halves, the
origin, and points up to $2.7\rho$ off the structure, acquiring the
network from all six in 0 to 13 steps for the $D$ tracker and 0 to 14 for
the $s_1$ tracker. They then ride the same segments in the same direction
and cross the isotropic point at the origin, with the formation
straddling the separatrix throughout (Fig.~\ref{fig:ch7:s1_tracker_sim}).
A run acquiring above the origin therefore rides an attracting OECS to
that point and continues onto a repelling one, on the same primitive and
the same scalar trench throughout. Neither needs a dedicated rule at the
crossing, and the $s_1$ tracker rides through the origin on ordinary
tangent selection. The primitives part only at the terminal step
(Sections~\ref{sec:ch7:d_verification} and~\ref{sec:ch7:s1_verification}).

\section{Behavior Under Noise}
\label{sec:ch7:noise}
% Ported from: Draft_11.tex Sec. IV-C (Behavior Under Noise).

\begin{figure}[h]
  \centering
  \includegraphics[width=0.75\textwidth]{ch7_flip_resolution.png}
  % FIGURE SOURCE: Paper_Writing/Separatrix_and_OW_Paper/figures/flip_resolution.png
  \caption{Far-saddle success of both trackers from the straddling start
  $(0, 0.35)$, $10^4$ trials per point, with each tracker's straddle
  retention (counted on successful trials) dashed. The dotted line marks
  50\% and the ticks each tracker's crossing. (a) Against $\sigma_{uv}$,
  $\sigma_p = 0$. (b) Against $\sigma_p$, $\sigma_{uv} = 0$.}
  \label{fig:ch7:flip_resolution}
\end{figure}

Both primitives sweep measurement and position noise separately from one
straddling start, $(0, 0.35)$, at 10\,000 trials per noise level from
independent headings. Noise-sweep trials run up to 500 steps and stop at
task completion or domain exit. Success is contact within 0.06 of the far
saddle $\mathbf{p}^*_1 = (0, -0.5)$, which a noise-free run from this
start reaches under either primitive, without formation collapse (an RMS
error above $4\rho$ in the fifteen inter-robot distances); each trial also
records straddle retention. A trial straddles when the robots fall on
both sides of the separatrix, and retention requires it never stop once
started, the failure \cite{34} reports. Holding one target fixed makes
a settle at the near saddle a directional failure rather than a loss of
the structure. Both succeed in every noise-free trial. Noise favors the
$D$ tracker, whose 50\% crossing is about four times the $s_1$ tracker's.

The $D$ tracker crosses 50\% at $\sigma_{uv} \approx 0.0077$, where
$\nabla\hat{D}$ stops being informative along the separatrix
($\sigma_{uv} \approx 0.0047$ to $0.0085$), and the $s_1$ tracker at
$\sigma_{uv} \approx 0.0019$ (Fig.~\ref{fig:ch7:flip_resolution}(a)).
Under position noise the ordering holds, $\sigma_p \approx 0.0081$
against 0.0023 (Fig.~\ref{fig:ch7:flip_resolution}(b)). The $s_1$
tracker's failure is a sign inversion at tangent selection
(Section~\ref{sec:ch7:s1_definition}). Supplying a clean gradient and
eigenframe to (\ref{eq:ch7:tangent_select}) at $\sigma_{uv} = 0.002$
restores the correct-saddle rate to 100\%. This fragility is attributed
to how long a sign decision persists before a new measurement checks it,
rather than to per-cycle estimation accuracy. The $D$ tracker re-signs
$\mathbf{w}_1$ against $\hat{\mathbf{v}}_0$ every cycle, while the $s_1$
tracker carries the sign in state through $\mathbf{t}_{\mathrm{ref}}$.
That reliance follows from objectivity and traversal together.
Objectivity rules out the fitted curvature
(Section~\ref{sec:ch7:surrogates}), the measured flow
(Section~\ref{sec:ch7:rotating}), and a world axis as sign references, and
traversal rules out the along-trench component of $\nabla\hat{s}_1$,
which reverses at the along-trench maximum the ride must cross, leaving
the primitive's own prior output. Straddle retention stays within five
points of success for the $D$ tracker but falls further below it for the
$s_1$ tracker, 47.6\% against 71.5\% at $\sigma_{uv} = 0.0015$. Most
$s_1$ trials that lose the straddle do so within 0.1 of the origin, where
the $s_1$ trench loses its transverse curvature
(Section~\ref{sec:ch7:s1_definition}).

\section{Behavior Under a Rotating Observer}
\label{sec:ch7:rotating}
% Ported from: Draft_11.tex Sec. IV-D (Behavior Under a Rotating
% Observer). This is the objectivity result and is reused in
% Chapter 8.

\begin{figure}[h]
  \centering
  \includegraphics[width=\textwidth]{ch7_objectivity_traverser.png}
  % FIGURE SOURCE: Paper_Writing/Separatrix_and_OW_Paper/figures/objectivity_traverser.png
  \caption{Objectivity trial, both primitives started from $(0.05, 0.40)$
  and run in the inertial frame (solid) and in a frame rotating at
  $\Omega = 0.23$ (dashed, pulled back to inertial coordinates). (a) The
  $s_1$ tracker's two runs coincide and capture at the same saddle. (b)
  The $D$ tracker's rotating-frame run loses the straddle across the
  crest (inset), the artifact predicted by (\ref{eq:ch7:D_not_objective}),
  and rejoins near the far saddle.}
  \label{fig:ch7:oecs_objectivity}
\end{figure}

Both primitives were run from $(0.05, 0.40)$, just inside the upper
saddle, on the same physical double gyre in two observer frames, the
inertial one and a frame rotating clockwise at $\Omega = 0.23$, a heading
drift of about $8^\circ$/h at the ocean time unit of
Section~\ref{sec:ch7:sb_setup}, where the field is
$\mathbf{v}'(\mathbf{p}', t) = \mathbf{Q}\mathbf{v}(\mathbf{Q}^\top
\mathbf{p}') + \Omega\,(-y', x')^\top$. The rate keeps the structure's
apparent speed along the tracked path, $\Omega\lVert\mathbf{p}\rVert \leq
0.11$, inside the command authority $k\,c_{\max} = 0.12$, the condition
Appendix~\ref{app:stability} requires. Rotating-frame trajectories are
pulled back to inertial coordinates and compared with the inertial run of
the same primitive.

Between acquisition and the far saddle, the $s_1$ tracker stays within
0.016 of the separatrix in the rotating frame, against 0.010 in the
inertial frame, keeps the straddle throughout, and both runs capture at
the same bottom saddle (Fig.~\ref{fig:ch7:oecs_objectivity}). The $D$
tracker's rotating-frame run is displaced by up to 0.084, against 0.010,
and loses the straddle for 36 steps across the crest, with all six robots
on one side of the separatrix, the failure \cite{34} reports.
Sweeping the rate, the $D$ tracker first loses the straddle at $\Omega
\approx 0.21$, where $\Omega\lVert\mathbf{p}\rVert$ at the crest is a
sixth of the command authority.

The frame's solid-body vorticity shifts the $D'$ landscape
(\ref{eq:ch7:D_not_objective}) and adds a transport term to the measured
flow, so both terms of the $D$ tracker are steered by a frame artifact.
On the separatrix $\omega = 0$ but $\partial_x\omega$ is nonzero
(Chapter~\ref{ch:preliminaries}), so $\Omega\omega$ tilts $D'$ across the
trench, moving the trench most at the origin and least at the saddles, so
the run rejoins the separatrix near $\mathbf{p}_1^*$. The $s_1$ tracker's
tests are scalar inequalities in quantities a rotation leaves invariant,
so its command co-rotates with the frame, and its residual is the
frame-transport term bounded in Appendix~\ref{app:stability}. Its one
non-objective input is the tangent seed taken from the measured flow,
which fixes the ride direction, and the frames agree while
$\Omega\lVert\mathbf{p}_{\text{seed}}\rVert <
\lvert\hat{\mathbf{v}}_0^\top\mathbf{t}_{\text{raw}}\rvert$.

\section{Santa Barbara Channel Demonstration}
\label{sec:ch7:sb_channel}
% Ported from: Draft_11.tex Sec. V (Santa Barbara Channel Trial), full.
% One sentence of data-provenance detail (CORDC, THREDDS-Ocean, retrieval
% date) is added from thesis.tex's "Ocean HFR Simulation" (lines
% ~3958-3968), which describes it more fully than Draft_11 does; every
% numeric operating-point and result value below is Draft_11's, not
% thesis.tex's, since thesis.tex's ocean trial used a different
% formation scale and gain and is not comparable.

\subsection{Setup}
\label{sec:ch7:sb_setup}

The ocean experiment uses hourly surface-current frames for the Santa
Barbara Channel from the HFRNet U.S.\ West Coast Radial-to-Total Vector
product \cite{58}, on the 2~km grid, which is public. The product is
produced by the Scripps Coastal Ocean Observing Research and Development
Center (CORDC) HFRNet network and archived at the NOAA National Centers
for Environmental Information, retrieved from the NOAA THREDDS-Ocean
server. The record is the one \cite{34} uses, the same network over
the same 28-h window, 29 frames from May 16, 2012 08:00 GMT to May 17,
2012 12:00 GMT, so these results are interpretable against prior work on
it. Kularatne and Hsieh track a time-reversed attracting boundary in the
same channel's May 2012 currents \cite{1}. Frames are interpolated
linearly in time. World coordinates on $[-0.65, 0.65]^2$ map affinely onto
a region centered on $(34.2^{\circ}\mathrm{N}, 120.4^{\circ}\mathrm{W})$.

Physical numbers follow from one unit block, length $L_w = 51.4$~km on
both axes and time $6000$~s, with a control period $\Delta t = 0.1$ time
units, so one step advances the field by $600$~s and the $168$-step trial
covers the record. The commanded-speed cap $\sqrt{2}\,k\,c_{\max}$ is then
$0.78$~m/s, within reach of a small autonomous surface vessel but above a
buoyancy glider. With longitude scaled by $\cos(34.2^{\circ})$ the map is
a uniform dilation, so the zero sets, trenches, and eigenframe are the
physical ones. The FTLE reference uses the same constants, over a 24-h
field for Fig.~\ref{fig:ch7:sb_channel} and 12~h for a persistence check.

The pentagon is enlarged to $\rho = 0.098$, $1.3\times$ the nominal
radius and $5.0$~km at this site, so the footprint spans about five cells
of the 2~km grid rather than under four, and starts at heading
$335^\circ$. At one control cycle per ten minutes any vessel time
constant is negligible, so $\alpha_{\text{mom}} = 0$. The gain and radius
set when the cluster reaches the bifurcation of the evolving ridge, and
so which branch it meets. The latch depth $s_{\mathrm{trim}} = 0.3$ is
set from this field's own $s_1$ scale. No noise level is estimated, since
the exactly determined fit leaves no residual, so the noise model of
Section~\ref{sec:ch7:noise} is validated on the double gyre only.

\subsection{Results}
\label{sec:ch7:sb_results}

\begin{figure}[h]
  \centering
  \includegraphics[width=\textwidth]{ch7_ocean_progression_2km.png}
  % FIGURE SOURCE: Paper_Writing/Separatrix_and_OW_Paper/figures/ocean_progression_2km.png
  \caption{$D$ and $s_1$ tracker paths from a single shared start at (a)
  7, (b) 14, (c) 21, and (d) 28~h, each drawn up to that time, with dots
  marking the current positions. The background is the 24-h forward FTLE
  field computed offline from the same 2~km data and anchored at the
  record start. Both follow the dominant ridge from the channel entrance
  to the bifurcation, where the $D$ tracker turns onto the southeastern
  branch toward the middle island. The $D$ run stops at 24~h, when its
  first robot reaches the island.}
  \label{fig:ch7:sb_channel}
\end{figure}

The field is divergent and not vorticity-free, so $s_2 = -s_1$ does not
hold and the two surrogates need not share trenches. Both were run with
no changes to the primitives, at the ocean operating point of
Table~\ref{tab:ch7:params}, from a shared start north of the channel
entrance, $(34.38^{\circ}\mathrm{N}, 120.41^{\circ}\mathrm{W})$. Over the
28-h trial both descend through the channel entrance and track the same
ridge to the bifurcation. There the $D$ tracker turns onto the
southeastern branch, and its run stops after 24~h, when its first robot
reaches the coast of the middle island. The $s_1$ tracker is still
descending when the record ends.

Fig.~\ref{fig:ch7:sb_channel} places both trials against the 24-h
forward-time FTLE field, computed offline in the style of \cite{49}.
Both follow the dominant ridge from the channel entrance to the
bifurcation closely. Mean distance from the tracked path to the nearest
95th-percentile ridge point is 1.5~km for the $D$ tracker and 2.1~km for
the $s_1$ tracker. The grid is 2~km and the formation radius 5.0~km. That
ridge was available to neither, and FTLE recomputed at four anchor times
shows it persisting while the finer structure downstream evolves.

The trial exercises tangent selection where neither the eigenframe nor
the trench network is axis-aligned, which the double gyre cannot do. No
two-dimensional minimum of $s_1$ lies within reach, so the capture test
never engages.

\section{Summary}
\label{sec:ch7:summary}
% New text, summarizing this chapter's own results. Draft_11's
% Conclusion also carries the cross-cutting objectivity synthesis
% (Chapter 8's territory) and future-work items (Chapter 9's
% territory, per the settled decision that Draft_11 Sec. VI stays out
% of this chapter); neither is repeated here.

A six-robot pentagon-plus-center formation recovers the local quadratic
model of a two-component field from one synchronized sample, giving the
Jacobian, its determinant $D$, and the smaller rate-of-strain eigenvalue
$s_1$ in closed form, together with their gradients and, for $D$, its
Hessian (Section~\ref{sec:ch7:estimation}). Both surrogates mark the
separatrix as a trench (Section~\ref{sec:ch7:surrogates}), and the $D$
tracker and $s_1$ tracker built on them acquire and ride the structure
from instantaneous measurements alone, with no pre-straddling and no
integration horizon.

On the double-gyre benchmark both trackers acquire, traverse, and cross
the isotropic point identically under clean measurements, parting only at
the terminal saddle. Under noise the $D$ tracker withstands about four
times the measurement noise the $s_1$ tracker does, with a matching
margin under position noise. Under a rotating observer the ranking
reverses: above a heading drift of about $7^\circ$/h the $D$ tracker's
formation no longer straddles the structure, while the $s_1$ tracker,
being objective, keeps the straddle and returns the same material curve
in both frames. A cluster with an absolute heading reference should carry
the $D$ tracker; one whose heading is dead-reckoned should carry the
$s_1$ tracker. On Santa Barbara Channel currents both primitives followed
the same transport corridor to its bifurcation, at mean distances of 1.5
and 2.1~km from an FTLE ridge neither could see, on a field neither
trials nor the double-gyre benchmark can fully stand in for.

% If hardware verification for second-order primitives happens before
% defense (per the March-June 2027 window discussed with the advisor),
% add a Hardware Verification subsection to each tracker above and a
% corresponding \section{Hardware Verification} here before this
% Summary.

% ===================================================================
% CHAPTER 8: CROSS-CUTTING FINDINGS
% ===================================================================
\chapter{Cross-Cutting Findings}
\label{ch:crosscutting}

% NEW SYNTHESIS CHAPTER. Every section here must state something
% neither source paper states alone; if a section only restates
% Chapter 6 or 7, fold it into Chapter 10 instead and drop to nine
% chapters total.

The zeroth-, first-, and second-order estimators of Chapters~\ref{ch:zeroth}, \ref{ch:first_order}, and \ref{ch:second_order} were built and validated one at a time, against different fields, formations, and failure modes. This chapter puts their stated results next to each other. Some sections below add nothing beyond that juxtaposition and say so; the remainder point out where the papers, read together, support a claim that neither states on its own.

\section{Noise Amplification with Fit Order}
\label{sec:ch8:noise_order}

% SOURCE: combine Ch.6 Error Analysis (sigma_p ~ ||J^{-1}|| sigma_v /
% sqrt(3)) with Ch.7's noise cliffs (Draft_11.tex Sec. VI-B,
% lines 873-924). The comparison across orders is new.

The first-order estimator (Chapter~\ref{ch:first_order}) reports one noise-propagation formula, mapping per-component sensor noise $\sigma_v$ into position scatter $\sigma_p$ through the inverse Jacobian of the fitted plane,
\begin{equation}
    \sigma_p \approx \frac{\lVert\mathbf{J}^{-1}\rVert\,\sigma_v}{\sqrt{3}} .
    \label{eq:ch8:sigma_p}
\end{equation}
% EVIDENCE: Systems Paper/cwaight_systems_paper.tex, Sec. VI-A (Error Analysis), the sigma_p formula and the sqrt(3) attributed to three-robot averaging.
With the calibration error of the printed-map sensor, $\sigma_v \approx 0.02$~m/s, this formula predicts roughly 0.03~m of scatter on the printed vortex, matching the observed 0.009--0.016~m precision after closed-loop time-averaging reported for the two hardware fields. Equation~(\ref{eq:ch8:sigma_p}) carries no explicit dependence on formation radius; it is evaluated at the single 0.33~m equilateral triangle used throughout that paper's hardware and simulation trials.
% EVIDENCE: Systems Paper/cwaight_systems_paper.tex, Sec. VI-A (Error Analysis).

The second-order estimator (Chapter~\ref{ch:second_order}) states its noise-propagation model in more granular form. Combining measurement noise $\sigma_{uv}$ on each robot's reading with position noise $\sigma_p$ at measurement time gives one effective noise level,
\begin{equation}
    \sigma_{\mathrm{eff}}^2 = \sigma_{uv}^2 + \tfrac{1}{2}\lVert\mathbf{J}\rVert_F^2\,\sigma_p^2 ,
    \label{eq:ch8:sigma_eff}
\end{equation}
% EVIDENCE: Separatrix_and_OW_Paper/Draft_11.tex, Sec. II-C (Estimator Sensitivity, Position Noise subsection), eq. sigma_eff.
and states directly that formation size scales the noise reaching each fitted coefficient by $\rho^{-q}$, where $q$ is the polynomial order of that coefficient (1 for the gradient terms, 2 for the curvature terms) and $\rho$ is the formation radius, so a larger formation suppresses noise most in the second-order terms.
% EVIDENCE: Separatrix_and_OW_Paper/Draft_11.tex, Sec. II-C (Estimator Sensitivity, Formation subsection).
Both surrogates built on this fit, $\hat{D}$ and $\hat{s}_1$, are read from the same coefficients but at different orders: the value $\hat{D}_0$ is a product of first-order ($q=1$) coefficients, while the Hessian $\hat{\mathbf{H}}_{D,0}$ and the gradient $\nabla\hat{s}_1$ draw on the second-order ($q=2$) coefficients.
% EVIDENCE: Separatrix_and_OW_Paper/Draft_11.tex, Sec. II-D (Eulerian Surrogates), eq. hess_det and eq. grad_s1.
Under Monte Carlo noise sweeps from a straddling start, the $D$ tracker's far-saddle success crosses 50\% at $\sigma_{uv} \approx 0.0077$ and the $s_1$ tracker's at $\sigma_{uv} \approx 0.0019$, roughly a factor of four apart; under position noise the crossings are $\sigma_p \approx 0.0081$ against $0.0023$.
% EVIDENCE: Separatrix_and_OW_Paper/Draft_11.tex, Sec. IV-B (Behavior Under Noise).
Both crossing values are small fractions of the commanded speed scale $c_{\max} = 0.04$ used in the same trials (Table~II of Chapter~\ref{ch:second_order}), so both second-order trackers fail at noise levels that are modest relative to the robots' own commanded speed.
% EVIDENCE: Separatrix_and_OW_Paper/Draft_11.tex, Table params (c_max, double-gyre column) and Sec. IV-B.

%% NOTE(new-math, checked 2026-09-29): the derivation below is not in
%% either paper. It generalizes Draft_11's rho^{-q} statement to every fit
%% order and recovers the critical-points paper's sigma_p formula.
Both results follow from one scaling argument that holds at every fit
order. Let the formation have a fixed shape $\boldsymbol{\xi}_i$ scaled to
radius $\rho$, so $\tilde{\mathbf{p}}_i = \rho\,\boldsymbol{\xi}_i$. Each
entry of $\boldsymbol{\phi}_k$ is a monomial of degree $|\alpha|$, so
\begin{equation}
    \boldsymbol{\Phi}_\rho = \boldsymbol{\Phi}_1\,\mathbf{S}_\rho,
    \qquad
    \mathbf{S}_\rho = \operatorname{diag}\bigl(\rho^{|\alpha|}\bigr),
    \label{eq:ch8:phi_scaling}
\end{equation}
and with independent measurement noise of standard deviation $\sigma$ on
each reading, the fitted coefficient of degree $q$ has standard deviation
\begin{equation}
    \sigma_{c_\alpha} = \rho^{-q}\,\sigma\,
        \bigl\lVert \mathbf{e}_\alpha^{\mathsf T}\boldsymbol{\Phi}_1^{-1}\bigr\rVert_2,
    \qquad q = |\alpha|,
    \label{eq:ch8:coef_noise}
\end{equation}
where the row norm depends on the formation's shape only. This is the
$\rho^{-q}$ law of the second-order paper, stated for any $k$.

A feature is located by driving some fitted quantity $f$ to zero, and a
perturbation $\delta f$ moves the estimated zero by
$\delta n \approx \delta f / |\partial f/\partial n|$ along the direction
$n$ in which $f$ changes. If $f$ is built from field derivatives of order
up to $m$, Equation~\eqref{eq:ch8:coef_noise} makes $\delta f$ scale as
$\rho^{-m}\sigma$ to leading order, so
\begin{equation}
    \delta n \;\sim\; \frac{\rho^{-m}\,\sigma}{|\partial f/\partial n|}.
    \label{eq:ch8:terminal_noise}
\end{equation}
For the critical point, $f = \mathbf{v}$ itself ($m = 0$) and its slope is
$\mathbf{J}$. From $\hat{\mathbf{p}}^* = \mathbf{p}_c -
\hat{\mathbf{J}}^{-1}\hat{\mathbf{v}}_c$,
\begin{equation}
    \delta\hat{\mathbf{p}}^* = -\mathbf{J}^{-1}\,\delta\hat{\mathbf{v}}_c
        + \mathbf{J}^{-1}\,\delta\hat{\mathbf{J}}\,\mathbf{J}^{-1}\mathbf{v}_c ,
    \label{eq:ch8:pstar_pert}
\end{equation}
and the second term vanishes as the centroid reaches the critical point,
where $\mathbf{v}_c \to \mathbf{0}$. For any nondegenerate triangle the
affine interpolant's value at the centroid is the mean of the three
readings, so $\sigma(\hat{\mathbf{v}}_c) = \sigma_v/\sqrt{3}$ per component
and Equation~\eqref{eq:ch8:sigma_p} follows, with no dependence on $\rho$.
For both second-order trackers the zero sought is that of a transverse
gradient, $\mathbf{w}_2^{\mathsf T}\nabla\hat{D}$ or
$\mathbf{P}\nabla\hat{s}_1$ (Chapter~\ref{ch:second_order}), which is
built from second derivatives ($m = 2$) and whose slope is the transverse
curvature of $D$ or $s_1$. The terminal error of the first-order law is
therefore independent of formation radius, while that of either
second-order tracker grows as $\rho^{-2}$. Equation~\eqref{eq:ch8:terminal_noise}
gives the scaling only; it does not by itself account for the factor of
four between the two second-order trackers, which share $m = 2$, and
neither paper runs the sweep over fit order at fixed radius that would
test it directly.

\section{Formation Size: Truncation Versus Noise}\section{Formation Size: Truncation Versus Noise}
\label{sec:ch8:formation_size}

% SOURCE: Ch.4 Sec. "Truncation Error and Formation Radius" applied
% side by side to the first- and second-order results.

The first-order paper states one half of the size trade-off explicitly: on a smooth nonlinear field the plane fit carries a linearization error that scales with the local curvature of the field, and this error decreases as the formation contracts toward the critical point.
% EVIDENCE: Systems Paper/cwaight_systems_paper.tex, Sec. II-B (Distributed Critical Point Estimation), paragraph after eq. 9.
It does not state the complementary half, that a smaller formation also raises the noise sensitivity of the same fit; the noise formula (\ref{eq:ch8:sigma_p}) is evaluated at one fixed formation size rather than swept, so the paper's own results do not show noise growing as the triangle shrinks. Equation~\eqref{eq:ch8:pstar_pert} explains why: at the critical point
only the order-zero coefficient enters the location error, so shrinking
the triangle does not raise the terminal scatter. The radius enters
through $\delta\hat{\mathbf{J}} \propto \rho^{-1}$ only while the
centroid is away from the critical point, during approach.

The second-order paper states both halves in closed form on the same page. Truncation error in the quadratic Taylor model is $\mathcal{O}(\lVert\mathbf{p}\rVert^3)$, which grows with the formation radius $\rho$ by construction, while the noise reaching each fitted coefficient scales as $\rho^{-q}$ (Section~\ref{sec:ch8:noise_order}), so a larger $\rho$ suppresses noise at the cost of truncation accuracy, and a smaller $\rho$ does the reverse.
% EVIDENCE: Separatrix_and_OW_Paper/Draft_11.tex, Sec. II-A (eq. quad_model_u/v, the O(||p||^3) truncation term) and Sec. II-C (Estimator Sensitivity, Formation subsection).
The paper reports a formation radius of $\rho = 0.075$ on the double-gyre benchmark and $\rho = 0.098$ (about $1.3\times$ the nominal radius, 5.0~km on the Santa Barbara Channel grid) on the ocean trial, chosen so the enlarged footprint spans about five cells of the 2~km radar grid rather than under four, but does not report a formal sweep of $\rho$ against either error term.
% EVIDENCE: Separatrix_and_OW_Paper/Draft_11.tex, Sec. III-A (Double-Gyre Benchmark, Setup) and Sec. V-A (Santa Barbara Channel Trial, Setup).

The two papers therefore differ in how completely each states the size trade-off inherent to its own fit order, not in whether the trade-off exists. The first-order paper's Error Analysis and Limitations sections describe the truncation side and leave the noise side implicit; the second-order paper's Estimator Sensitivity section states both sides as closed-form functions of $\rho$. Neither paper runs a formation-radius sweep against the estimation error it reports (Section~\ref{sec:ch8:noise_order}), so this remains a statement about what each source proves, not a joint prediction of an optimal radius.
% EVIDENCE: Systems Paper/cwaight_systems_paper.tex, Sec. VI-C (Limitations), fourth item (affine-model assumption over the cluster's spatial extent).

\section{Formation Conditioning Versus Field Conditioning}
\label{sec:ch8:conditioning}

Both papers separate two independent sources of error amplification, one set by the geometry of the robots and one set by the field being sampled, and both name each with a condition number, though on different matrices.

On the formation side, the first-order paper reports the condition number $\kappa(\mathbf{A})$ of the position matrix in its plane fit, minimized among formations of a given size by the equilateral triangle used throughout, and diverging as the three robots approach collinearity.
% EVIDENCE: Systems Paper/cwaight_systems_paper.tex, Sec. VI-A (Error Analysis), paragraph on the equilateral triangle and kappa(A).
The second-order paper reports the condition number $\kappa(\boldsymbol{\Phi})$ of its $6\times 6$ formation matrix, equal to 8.26 at every radius for the pentagon-plus-center formation used throughout, and rising to 564 when the center robot is displaced to $0.99\rho$, within 0.6\% of the ring.
% EVIDENCE: Separatrix_and_OW_Paper/Draft_11.tex, Sec. II-C (Estimator Sensitivity, Formation subsection), the kappa(Phi) = 8.26 and 564 values.
Both $\kappa(\mathbf{A})$ and $\kappa(\boldsymbol{\Phi})$ are properties of robot placement alone, computable before any field measurement is taken, and both are stated to be under the formation controller's control.

On the field side, the first-order paper names a second, independent condition number, $\kappa(\mathbf{J})$, that measures how the local field structure amplifies coefficient errors into critical-point location errors, growing without bound as $\det(\mathbf{J}) \to 0$, and states explicitly that unlike $\kappa(\mathbf{A})$ this factor is an intrinsic property of the field and not under the designer's control.
% EVIDENCE: Systems Paper/cwaight_systems_paper.tex, Sec. VI-A (Error Analysis), paragraph beginning "Further errors could arise from poor conditioning of the estimated Jacobian."
The second-order paper does not name a field-side condition number.
Equation~\eqref{eq:ch8:terminal_noise} identifies what plays that role:
the field-side amplifier is the inverse of the slope of whatever quantity
the controller drives to zero. For the critical point that is
$\lVert\mathbf{J}^{-1}\rVert$, which diverges as $\det\mathbf{J} \to 0$.
For the $D$ tracker it is $1/|\lambda_2|$, the inverse transverse
curvature of the fitted $D$ that already divides the across-trench
command in Equation~(\ref{eq:ch7:d_law}), which diverges as the trench
flattens; for the $s_1$ tracker it is the inverse transverse curvature of
$s_1$. The ratio $a_\perp = \kappa_\perp/\hat{\kappa}_\perp$ in the
second-order stability proposition (Appendix~\ref{app:stability}) is a
different quantity, the ratio of true to fitted curvature, which measures
how far the estimated amplifier is from the true one.
%% NOTE(new-math, checked 2026-09-29): identification of the field-side
%% amplifier as the inverse slope follows from eq:ch8:terminal_noise.

The chapter keeps $\kappa(\mathbf{A})$, $\kappa(\boldsymbol{\Phi})$, $\kappa(\mathbf{J})$, and $a_\perp$ notationally distinct throughout, since the sources never equate them and use them in different fits.

\section{Instantaneous Estimation and Time-Varying Fields}
\label{sec:ch8:time_varying}

Both estimators are memoryless. The first-order paper states that its estimation framework uses only instantaneous measurements and carries no state between control cycles, and argues from this that the estimator does not assume field stationarity, so it should in principle apply to a field whose shape or strength changes with time, distinct from a critical point that merely translates while the field's structure around it stays fixed. It states plainly that validating the estimator on a genuinely time-varying field remains future work, and that none of its own trials, on six analytical fields or two static reconstructed fields, exercises this case.
% EVIDENCE: Systems Paper/cwaight_systems_paper.tex, Sec. VI-E (Future Work), second and third paragraphs.
It separately flags that its robot dynamics respond only to commanded velocity, so a field that physically advects the platform, as an ocean current would, adds a disturbance term the controller does not compensate for, and states that extending the architecture to this case is a future-work priority.
% EVIDENCE: Systems Paper/cwaight_systems_paper.tex, Sec. III-A (Robot Controller Layer) and Sec. VI-E (Future Work), fourth paragraph.

The second-order paper's estimator is instantaneous in the same sense, each cycle's fit seeing only that cycle's measurements, and it runs a trial the first-order paper does not: 29 hourly frames of measured Santa Barbara Channel surface currents over a 28-hour window, with both the $D$ and $s_1$ trackers tracking the same evolving ridge to its bifurcation.
% EVIDENCE: Separatrix_and_OW_Paper/Draft_11.tex, Sec. II-C (Estimator Sensitivity, Measurement Noise subsection) and Sec. V (Santa Barbara Channel Trial).
It is explicit, however, that this is a finding on one time-varying record rather than a stability guarantee: neither tracker has an analytic guarantee for a time-varying flow, and the field is sampled from the gridded product with no added noise, so the trial does not test the estimator's noise behavior (Section~\ref{sec:ch8:noise_order}) under time variation.
% EVIDENCE: Separatrix_and_OW_Paper/Draft_11.tex, Sec. VI (Limitations and Future Work), third and fourth paragraphs.
Both papers share the same no-advection premise: the sensed field enters the estimator only as a reading, never as a force or drift term on the robot itself, and both flag this as an idealization rather than a demonstrated property of the physical world. The first-order paper states this is exact on its own printed-map testbed, where no physical flow acts on the robots, and the second-order paper states the same premise is exact for the planned Decabot follow-on of its own method, but is an untested extrapolation on the ocean field it actually ran.
% EVIDENCE: Systems Paper/cwaight_systems_paper.tex, Sec. III-A. Separatrix_and_OW_Paper/Draft_11.tex, Sec. VI (Limitations and Future Work), fourth paragraph.

Read together, the second-order paper carries out the experiment the first-order paper's Future Work section calls for, an instantaneous estimator run against a field that changes in time, but under a weaker field model (real radar currents rather than a controlled sweep of field shape or strength) and without closing the advection gap either paper flags.

\section{Objectivity and Sign References}
\label{sec:ch8:objectivity}

% SOURCE: Ch.7 "Behavior Under a Rotating Observer," generalized as a
% cross-cutting statement about what each fit order can and cannot
% certify under an observer transformation.

The second-order paper decomposes the fitted Jacobian $\hat{\mathbf{J}} = \hat{\mathbf{S}} + \hat{\mathbf{W}}$ into a symmetric rate-of-strain part $\hat{\mathbf{S}}$ and an antisymmetric spin part $\hat{\mathbf{W}}$ carrying the vorticity, and shows that a frame rotating at rate $\Omega$ adds a uniform spin, raising the vorticity by $2\Omega$ and shifting the determinant surrogate by $D' = D + \Omega\omega + \Omega^2$, while the strain eigenvalue $s_1$ is unaffected because $\hat{\mathbf{S}}$ only co-rotates. It runs a rotating-observer trial at $\Omega = 0.23$: the $D$ tracker's rotating-frame run is displaced by up to 0.084 from its inertial-frame run (against 0.010 for the $s_1$ tracker) and loses the straddle for 36 steps across the crest, while the $s_1$ tracker's two runs stay within 0.016 of each other and both capture at the same saddle.
% EVIDENCE: Separatrix_and_OW_Paper/Draft_11.tex, Sec. II-D (Eulerian Surrogates, Relation Between the Fields) and Sec. IV-C (Behavior Under a Rotating Observer).
It also states which sign references remain valid under this transformation: objectivity rules out the fitted curvature, the measured flow, and a world axis as sign references for the $s_1$ tracker's tangent selection, leaving only the primitive's own prior output, carried in state as $\mathbf{t}_{\mathrm{ref}}$.
% EVIDENCE: Separatrix_and_OW_Paper/Draft_11.tex, Sec. IV-B (Behavior Under Noise), paragraph on the D tracker re-signing w1 against v0-hat versus the s1 tracker carrying sign through t_ref.

%% NOTE(new-math, checked 2026-09-29): extension to the first-order
%% classification; the first-order paper runs no rotating-observer trial.
The same transformation acts on the first-order classification, which is
read off the eigenvalues of the full Jacobian (Table~I of
Chapter~\ref{ch:first_order}). Writing $\mathbf{J} = \mathbf{S} +
\tfrac{\omega}{2}\mathbf{E}$ with $\mathbf{E} = \left[\begin{smallmatrix}
0 & -1\\ 1 & 0\end{smallmatrix}\right]$, an observer rotating at rate
$\Omega$, signed as in Chapter~\ref{ch:second_order} so that the measured
vorticity rises by $2\Omega$, replaces $\mathbf{J}$ by $\mathbf{J}' =
\mathbf{J} + \Omega\mathbf{E}$ in the co-rotating basis. Since $\operatorname{tr}\mathbf{E} = 0$,
\begin{equation}
    \operatorname{tr}\mathbf{J}' = \operatorname{tr}\mathbf{J},
    \qquad
    \det\mathbf{J}' = \det\mathbf{S} + \bigl(\tfrac{\omega}{2}+\Omega\bigr)^2
        = \det\mathbf{J} + \Omega\,\omega + \Omega^2 ,
    \label{eq:ch8:rotating_J}
\end{equation}
the second identity being the shift $D' = D + \Omega\omega + \Omega^2$ of
the second-order paper. The real part of the eigenvalues,
$\alpha_{\text{eig}} = \tfrac{1}{2}\operatorname{tr}\mathbf{J}$, is
unchanged, so an observer's rotation cannot turn a sink into a source or
either into a center. The determinant and the discriminant
$(\operatorname{tr}\mathbf{J})^2 - 4\det\mathbf{J}$ do change, so a
sufficiently fast rotation can relabel a saddle ($\det\mathbf{J} < 0$) as
a node or spiral, or a node as a spiral. The location of the critical
point is not objective either: the rotating observer's velocity differs
from $\mathbf{v}$ by a rigid rotation about the rotation center, which
is nonzero away from that center, so the zeros of the measured field
move. The first-order paper
tests none of this, so these are consequences of the shared
decomposition, not demonstrated failures. The orbital law's tangent,
obtained by a fixed $-90^\circ$ rotation of the radial vector, is also a
world-frame sign convention of the kind the second-order paper rules out
as a reference for a rotating or dead-reckoned cluster.
% EVIDENCE: Systems Paper/cwaight_systems_paper.tex, Table I (field_summary) and Sec. II-D (Adaptive Navigation Control Laws, Orbital Navigation, the -90 degree rotation).
A reflection of the sensor frame is harsher than a rotation: it reverses
the sign of $\det\hat{\mathbf{J}}$, so it exchanges saddles with nodes
while leaving the located critical point in place. The hardware logs
show this directly (Chapter~\ref{ch:first_order}): with the saddle map's
frame unregistered, the estimate that located the saddle correctly read
it as a stable node or spiral in every terminal cycle. Location is insensitive to frame
conventions that type classification depends on.
% EVIDENCE: experiments/outputs/classify_hw/summary.json (raw_unregistered_terminal_confusion).

\section{What a Fit of Each Order Can Certify at the Terminal Step}
\label{sec:ch8:terminal_certify}

% NEW. Per the plan's own note: this is the section most likely to be
% genuinely new thinking. Draft_11.tex states the quadratic fit can
% only certify a vanishing gradient; the first-order controller's
% terminal behavior is stiction-limited. State these side by side.

The zeroth-order primitives (Chapter~\ref{ch:zeroth}) certify nothing about location. The vector-sum and vector-to-scalar primitives compress the sensed field to a heading or a scalar gradient direction and command motion along it, with no mechanism to detect or correct radial error with respect to any point. The first-order paper's own simulated comparison reports both drifting outward in a vortex field and states that in a saddle field the vector-sum controller cannot produce an orbital trajectory at all.
% EVIDENCE: Systems Paper/cwaight_systems_paper.tex, Sec. VI-B (Comparison with Other Primitives).
A separate hardware study of these two zeroth-order primitives on a 3-robot Decabot cluster (fixed formation $p=q=0.35$~m, $\beta=1.05$~rad), 10 trial runs per environment-primitive combination, found the outward-drift failure mode is not simulation-only: orbiting the fixed vortex with vector-sum produced a stable outward drift at all angles that the source states was not predicted by its own simulation studies and attributes to momentum effects the simulation did not capture, and it states that orbiting a center at a fixed radius with these primitives remains an open problem the testbed is positioned to address. The same study found vector-sum converges when the target field itself supplies inward flow: on a sinking vortex, vector-sum drove the cluster to settle around 0.1~m from center, matching its own simulation prediction. The vector-to-scalar primitive, tested only on the fixed vortex (10 runs, plus a second set of 10 with randomized starts), settled within 0.2~m of center, with a similar steady-state error in the randomized-start set; it was not tested on the sinking vortex.
% EVIDENCE: PhD Thesis/sources/idetc2025_testbed.txt \cite{39}, Sec. 3.5 (Control Primitive Testing Methodology) and Sec. 4.2 (Control Primitive Results on Testbed).

The first-order attraction and orbital laws (Chapter~\ref{ch:first_order}) certify convergence to the estimated critical point exactly, in closed form, given a stationary estimate: the proportional law's linear error dynamic guarantees global exponential convergence with time constant $\tau = 1/k$. In the idealized limit, no stiction, no velocity cap, and instantaneous actuator response, the reported trials confirm convergence to the controller's $10^{-6}$~m command deadband, with the critical-point estimate itself exact to machine precision on the six analytical fields. What limits the terminal step in every reported physical or physically modeled trial is not the estimator or the control law but actuator stiction: a trial is counted as convergent once the trajectory remains bounded and terminates within the stiction-limited neighborhood of the critical point, and the reported final-position statistics (Section~\ref{sec:ch8:noise_order}) quantify that neighborhood, not a limit on what the linear fit can certify.
% EVIDENCE: Systems Paper/cwaight_systems_paper.tex, Sec. II-D (eq. 10, the exponential convergence claim), Sec. IV-A (the stiction-limited-neighborhood convergence criterion), and Sec. VI-A (the no-stiction, alpha=0 control experiment).

The second-order trackers (Chapter~\ref{ch:second_order}) certify less than the first-order law does, and say so explicitly. The $s_1$ tracker's capture test is built on the gradient rather than the curvature, because its fitted Hessian $\hat{\mathbf{H}}_{s_1}$ is negative semidefinite for every field and formation by construction, so it cannot itself certify a minimum; the paper states that a vanishing gradient suffices only because the physical target, a flow saddle, happens to be a true two-dimensional minimum of $s_1$, a fact about the structure being tracked rather than a fact the fit can verify on its own. The $D$ tracker's capture test does examine the fitted curvature, firing when the fitted Hessian eigenvalues satisfy $\lambda_1\lambda_2 \geq 0$, a condition the true $D$ meets at a flow saddle; but the paper states directly that a negative-definite fit carries no stability claim, and that near the flow saddles on its own benchmark, second derivatives vanish and truncation error sets the sign of the fitted eigenvalues, so even this sign test is not a certificate away from the one benchmark case it was checked against.
% EVIDENCE: Separatrix_and_OW_Paper/Draft_11.tex, Sec. III-B (The D Tracker, eq. d_capture and surrounding paragraph) and Sec. III-C (The s1 Tracker, capture-test paragraph).

%% NOTE(new-math, checked 2026-09-29): pattern across the three fit orders; no source paper compares itself to the other two. See eq:ch8:terminal_noise.
Across the three orders, what the terminal step certifies gets weaker, not stronger, as the fit order rises. The zeroth-order primitives certify nothing about location, only a direction. The first-order law certifies exact convergence to its own estimate in closed form, and the reported terminal error traces entirely to actuator stiction, a limit external to the estimator. The second-order trackers, with a richer fit and two more surrogate quantities available, explicitly cannot certify a minimum from their own curvature estimate at all; their capture rules work only because the physical target being sought is independently known to be the kind of point ($D$: a local minimum of the fitted quantity; $s_1$: a vanishing-gradient point) the rule checks for, a fact supplied from outside the fit rather than recovered by it.

\section{Primitive Library}
\label{sec:ch8:primitive_library}

% SOURCE: modeled on Hart's Appendix B primitive table. One row per
% primitive: fit order, minimum robots, feature, surrogate, stability
% result, verification level (sim only / sim + hardware).

Table~\ref{tab:ch8:primitive_library} collects the primitives discussed in Chapters~\ref{ch:zeroth}, \ref{ch:first_order}, and \ref{ch:second_order} in one place. Stability and verification entries are copied from what each source states; "none stated" marks a primitive whose source paper reports no formal stability result. The zeroth-order rows' hardware entries are drawn from the testbed paper \cite{39}.
% EVIDENCE: zeroth-order rows from idetc2025_testbed.txt \cite{39} (hardware) and the critical-points paper's primitive comparison (simulation); first- and second-order rows from the sections cited in Sections \ref{sec:ch8:noise_order}--\ref{sec:ch8:terminal_certify} above.

\begin{table}[t]
  \centering
  \footnotesize
  \caption{Primitive library across the three fit orders. Primitives built on fewer robots than the fit's theoretical minimum (vector-sum, vector-to-scalar) instead report the cluster size used in their own trials.}
  \label{tab:ch8:primitive_library}
  \begin{center}
  \setlength{\tabcolsep}{3pt}
  \begin{tabular}{p{1.5cm}cc p{1.9cm} p{2.2cm} p{2.9cm} p{2.5cm}}
    \toprule
    \textbf{Primitive} & \textbf{Order} & \textbf{Min.\ robots} & \textbf{Feature} & \textbf{Surrogate / estimate} & \textbf{Stability result} & \textbf{Verification} \\
    \midrule
    Vector-sum & 0 & 3 (as tested) & heading toward mean flow & mean sensed velocity vector & none stated; outward drift orbiting a fixed vortex (sim.\ and hw.), converges on a sinking vortex (hw.), no orbit possible in a saddle (sim.) & simulation \cite{38}; hardware, fixed and sinking vortex, 10 runs each \cite{39} \\
    Vector-to-scalar & 0 & 3 (as tested) & steepest-ascent / descent direction on flow magnitude & flow magnitude and its finite-difference normal & none stated; converges on fixed vortex (hw.); also drifts when adapted for orbital use in a saddle (sim.) & simulation \cite{38}; hardware, fixed vortex only, 10 runs + 10 randomized-start \cite{39} \\
    Attraction & 1 & 3 & critical point & $\mathbf{p}^* = -\mathbf{J}^{-1}\mathbf{h}$ from the fitted plane & global exponential convergence to the estimated point, $\tau = 1/k$, given a stationary estimate & hardware, 157 convergence trials, 100\% success \\
    Orbital & 1 & 3 & fixed-radius standoff about a critical point & $\mathbf{p}^*$ as above, plus radial/tangential decomposition & radial error obeys $\dot{e}_r = -k_r e_r$, exponential convergence to $r_d$; decoupled from angular dynamics & hardware, 12 trials (6 radii $\times$ 2 fields) \\
    $D$ tracker & 2 & 6 & Okubo--Weiss / determinant separatrix (mountain-pass trench) & $\hat{D}_0 = \det\hat{\mathbf{J}}(\mathbf{p}_c)$ and Hessian $\hat{\mathbf{H}}_{D,0}$ & practically stable transverse to trench, $\limsup|n| \leq \delta/a_\perp$; capture test carries no stability claim on a negative-definite fit & simulation only: double-gyre clean and noise trials ($10^4$/pt) plus 28-h Santa Barbara Channel trial \\
    $s_1$ tracker & 2 & 6 & objective (rate-of-strain) separatrix & smaller strain eigenvalue $\hat{s}_1$ and $\nabla\hat{s}_1$ & same practical-stability bound as $D$ tracker; fitted Hessian negative semidefinite by construction, cannot certify a minimum & simulation only: double-gyre clean and noise trials ($10^4$/pt) plus 28-h Santa Barbara Channel trial \\
    \bottomrule
  \end{tabular}
  \end{center}
\end{table}

% ===================================================================
% CHAPTER 9: LIMITATIONS AND FUTURE WORK
% ===================================================================
\chapter{Limitations and Future Work}
\label{ch:limitations}

\section{Limitations}

\subsection{Planar Fields}
% Ported from: cwaight_systems_paper.tex Sec. VI-C (Limitations);
% thesis.tex Ch.4 Limitations (lines 2348-2350); idetc2025_testbed.txt
% Sec. 5 (System Limitations, 2D constraint).

The zeroth-order primitives of Chapter~\ref{ch:zeroth}, hardware-tested
on the HSV-encoded testbed of Chapter~\ref{ch:tools}, and the
first-order attraction and orbital primitives of
Chapter~\ref{ch:first_order} and the second-order trackers of
Chapter~\ref{ch:second_order} are all restricted to planar vector
fields. Extending the first-order estimator to three dimensions would
require at minimum four non-coplanar robots, since an affine model in
three dimensions carries four unknowns per velocity component rather
than three. Extending the testbed's own HSV floor-map encoding to a
third dimension would separately require additional sensing capability
beyond the color sensor used throughout \cite{39}. No
three-dimensional extension of either the estimator or the testbed
representation is developed in this dissertation.

\subsection{No Advection}
\label{sec:ch9:no_advection}
% Ported from: cwaight_systems_paper.tex Sec. III-A (lines 213-215);
% Draft_11.tex Sec. VII (lines 1088-1094).

In both control architectures, the sensed field enters only the
estimation layer. It does not appear as a force or drift term in the
robot dynamics, so each robot's velocity responds to its commanded
value alone, and no physical flow acts on the platform. Every result in
Chapters~\ref{ch:first_order} and~\ref{ch:second_order}, including the
second-order tracker's run on real Santa Barbara Channel current data,
assumes the field is sensed rather than felt. That premise is exact for
a Decabot deployment, where the field is encoded and read rather than
physically experienced \cite{39}, but on a genuine ocean deployment
it is an extrapolation, since the current would advect the vehicle
directly, and it is untested against real advection in either chapter.

\subsection{Centralized Control and Team Size}
% Ported from: Draft_11.tex Sec. VII (lines 1061-1065);
% cwaight_systems_paper.tex Sec. VI-C (lines 484-485).

Cluster space control is centralized, requires global state, becomes
intractable for suitably large clusters, and is impractical without
reliable high-bandwidth communication among the robots \cite{53}.
Every result in this dissertation, three robots for the zeroth-order
and first-order primitives of Chapters~\ref{ch:zeroth}
and~\ref{ch:first_order}, and six for the
second-order trackers of Chapter~\ref{ch:second_order}, is a small
cluster in a limited workspace with ample bandwidth. Team size was not
varied within a fit order to probe this boundary; the estimation
framework of Chapter~\ref{ch:estimation} supports larger, overdetermined
teams at each order, but larger formations raise communication
bandwidth, computation, and formation-control-stability costs this
dissertation does not quantify. All three chapters also share a
centralized architecture; extending any of the three primitive families
to a decentralized, swarm-style architecture remains open, and the
testbed of Chapter~\ref{ch:tools} was not built with the controller
substitution such a comparison would need \cite{39}.

\subsection{Static and Simulated Fields}
% Ported from: cwaight_systems_paper.tex Sec. VI-C (lines 481-483);
% Draft_11.tex Sec. VII (line 1088); idetc2025_testbed.txt Sec. 5 (System
% Limitations, static-field restriction).

The testbed of Chapter~\ref{ch:tools}, and with it the zeroth-order
primitives of Chapter~\ref{ch:zeroth} and the first-order primitives of
Chapter~\ref{ch:first_order}, is restricted to static printed vector
fields; there is no clear path from a printed colormap to a reactive,
dynamic environment, so no hardware result in either chapter tests
performance against a field that changes during a trial
\cite{39}. Ground truth ties directly to the print, but the
trials do not exercise dynamic field behavior. The second-order trackers
of Chapter~\ref{ch:second_order} are evaluated on the steady double
gyre and on real, time-varying Santa Barbara Channel current data, but
every second-order result is simulated: no hardware trial of either
tracker exists, and the ocean run itself samples the gridded current
product without added measurement noise, so the noise-robustness
results of that chapter rest on the double gyre alone. More generally,
an indoor, small-scale testbed of this kind does not necessarily
represent the scale of a real oceanic or atmospheric deployment
\cite{39}.

\subsection{Sampling Rate and Communication}
% Ported from: cwaight_systems_paper.tex Sec. VI-C (lines 487-488);
% Draft_11.tex Sec. VII (centralized-communication clause,
% lines 1061-1065).

Both control architectures assume a sampling rate high enough that the
closed loop approximates continuous control, and both assume every
robot reports its position and field measurement synchronously each
cycle. Chapter~\ref{ch:first_order}'s hardware trials run this loop at
10~Hz over WiFi. Robustness to lower update rates, to communication
dropouts, to partial measurement sets, or to asynchronous reporting was
not tested in either chapter. Every robot also samples only a single
point in the field at its current location, with no look-ahead; this
limits testing of any predictive or path-planning extension to the
purely reactive primitives studied here \cite{39}.

\subsection{Physical Testbed Constraints}
% Ported from: idetc2025_testbed.txt Sec. 5 (System Limitations).

The Decabot testbed of Chapter~\ref{ch:tools} that hosts every hardware
result in this dissertation carries its own physical constraints,
independent of estimation order. The workspace is bounded between
roughly 1~m~$\times$~1~m and 3~m~$\times$~6~m, battery life limits a
continuous trial to about one hour per robot, and the commandable
velocity range is 0.05 to 0.3~m/s, set by the stiction floor and the
motor's rated maximum \cite{39}. The fleet is homogeneous and
carries no obstacle-avoidance sensing, so heterogeneous multi-robot
behavior and obstacle-rich environments are outside every hardware
result reported here. The printed floor map itself wears under repeated
robot traffic and needs periodic replacement, an operational cost
distinct from any algorithmic limitation \cite{39}.

\subsection{Estimator Conditioning and Tracker Robustness Near
Degeneracies}
\label{sec:ch9:estimator_conditioning}
% Ported from: Draft_11.tex Sec. VII (lines 1067-1085);
% cwaight_systems_paper.tex Sec. VI-A (condition number discussion,
% line 458).

Chapter~\ref{ch:estimation}'s formation matrix $\boldsymbol{\Phi}$ is
singular when the robots lie on a common curve of the fit's own degree,
a line for the first-order estimator and a conic for the second-order
estimator. The first-order estimator's degradation as the three-robot
formation approaches collinearity is discussed in
Chapter~\ref{ch:first_order}; monitoring $\kappa(\boldsymbol{\Phi})$
online and correcting for it is not implemented in either chapter. No
field tested in Chapter~\ref{ch:second_order} drives the six-robot
formation toward a common conic, so the second-order estimator is
untested in the regime its own conditioning limit describes.

Two further robustness gaps are specific to the second-order trackers.
The eigenvector comparison the $s_1$ tracker uses for tangent selection
weakens near saddle points; the capture latch of
Chapter~\ref{ch:second_order} damps, but does not exclude, the residual
risk of a wrong branch under heavy noise. The $D$ tracker's traversal
past a saddle is untested outside the double-gyre benchmark, and on the
benchmark it already runs outside that tracker's own stated premise,
since the fitted Hessian is locally indefinite where the true Hessian
is not.

\subsection{Second-Order Hardware Verification}
\label{sec:ch9:second_order_hw}
% Ported from: Draft_11.tex Sec. VII ("All results are simulated,"
% line 1088); thesis.tex Ch.5 Limitations (lines 4637-4638); parallel
% noted per Hart's Sec. 7.2 preliminary primitives.

Every second-order result in Chapter~\ref{ch:second_order}, on the
double gyre and on the Santa Barbara Channel record alike, is
simulation only. No hardware trial of the $D$ tracker or the $s_1$
tracker exists on the Decabot testbed of Chapter~\ref{ch:tools}, and
the calibration and dynamics identification that testbed provides have
so far validated only the zeroth- and first-order primitives of
Chapters~\ref{ch:zeroth} and~\ref{ch:first_order}. The second-order
trackers are, in this sense, in the same position the ridge, trench,
and saddle-point primitives of an earlier reactive-swarm architecture
occupied: developed and verified in simulation, with hardware
validation still to come \cite{59}.

\section{Future Work}

\subsection{Third-Order Estimation from Ten Robots}
% Ported from: Draft_11.tex Sec. VII, item c) (lines 1096-1100),
% generalizing the minimality argument of Chapter~\ref{ch:estimation}
% (Draft_11.tex Sec. II-B, lines 208-233).

The estimator of Chapter~\ref{ch:estimation} generalizes to a cubic
local model by the same counting argument that fixes the first- and
second-order minimums. A planar cubic basis has ten terms per velocity
component, $1, x, y, x^2, xy, y^2, x^3, x^2y, xy^2, y^3$, so ten robots
is the minimum formation size, and the formation matrix is singular
exactly when the ten sampling points lie on a common cubic curve, the
third-order analogue of the collinear and common-conic degeneracies of
Chapters~\ref{ch:first_order} and~\ref{ch:second_order}. A cubic fit
recovers the third derivatives of the field that the quadratic model of
Chapter~\ref{ch:second_order} truncates away. Those third derivatives
are exactly what is missing from the true Hessian of $D$ and from the
transverse curvature of an $s_1$ trench; recovering them would remove
the structural bias in $\hat{\mathbf{H}}_D$ noted in
Section~\ref{sec:ch9:ow_corner_branch} below, rather than only bound
it.

\subsection{Online Formation Reshaping}
% Ported from: Draft_11.tex Sec. VII, item b) (lines 1096-1098);
% thesis.tex Ch.5 Future Work (lines 4665-4668); cwaight_systems_paper.tex
% Sec. VI-A (fallback strategies, line 458).

Adapting a formation's shape online, the pair lengths and SAS
parameters of Chapter~\ref{ch:estimation}'s formations, to preserve the
conditioning of $\boldsymbol{\Phi}$ as a cluster nears a degenerate
configuration, is future work at every fit order. For the first-order
triangle of Chapter~\ref{ch:first_order}, this means resizing or
reshaping as the formation approaches collinearity; for the
second-order pentagon of Chapter~\ref{ch:second_order}, it means the
same as the formation approaches a common conic. Neither chapter
implements the fix; both leave it as the direct answer to the
conditioning gap of Section~\ref{sec:ch9:estimator_conditioning} above.

\subsection{Advection-Aware Control}
\label{sec:ch9:advection_control}
% Ported from: cwaight_systems_paper.tex Sec. VI-D (Future Work,
% line 496), paired with the No-Advection limitation above
% (Draft_11.tex Sec. VII, lines 1091-1094).

Extending the robot dynamics model of Chapter~\ref{ch:tools} to include
a direct advective term, so that a physical current displaces the
platform rather than only the sensed field, is a priority for future
work in both estimation chapters. This closes the gap identified in
Section~\ref{sec:ch9:no_advection} above: it is the difference between
the Decabot testbed, where no physical flow acts on the robot, and a
genuine ocean or atmospheric deployment, where it does.

A complementary route attacks the same sim-to-reality gap from the
robot side rather than the field side: a digital twin of the Decabot
platform, built by capturing each robot's velocity and acceleration
response directly rather than assuming the first-order momentum model
of Chapter~\ref{ch:tools}, would let the simulator used throughout
Chapters~\ref{ch:first_order} and~\ref{ch:second_order} reproduce
hardware dynamics more faithfully, independent of whether an advective
term is also added \cite{39}.

\subsection{Three-Dimensional Fields}
% Ported from: cwaight_systems_paper.tex Sec. VI-C. The second-order
% count and degeneracy condition are new text, the three-variable form of
% Chapter~\ref{ch:estimation}'s minimality and degeneracy results.
%% NOTE(new-math, checked 2026-09-29): not in any source paper.

A polynomial of total degree $k$ in three variables has
$\binom{k+3}{3} = (k+1)(k+2)(k+3)/6$ coefficients per field component,
so the minimality argument of Chapter~\ref{ch:estimation} carries over
with 4, 10, and 20 unknowns per component for $k = 1, 2, 3$. A
first-order estimate in three dimensions therefore needs at least four
robots, and it is degenerate exactly when they are coplanar. A
second-order estimate needs at least ten robots; with exactly ten, the
formation matrix is singular exactly when the robots lie on a common
quadric surface, since a null vector of the matrix is the coefficient
vector of a quadric through all ten positions, the three-dimensional
counterpart of the common-conic condition. No formation avoiding this
set has yet been designed, and the second-order surrogates themselves
change in three dimensions, where the Okubo-Weiss partition generalizes
to the $Q$-criterion and the strain tensor has three eigenvalues. The
zeroth-order primitives of Chapter~\ref{ch:zeroth} were also developed
and tested only in two dimensions.

\subsection{Additional Topological Features}
% Ported from: cwaight_systems_paper.tex Sec. VI-D (line 498);
% thesis.tex Ch.4 Future Work (lines 2375-2386).

Chapter~\ref{ch:second_order} tracks two Eulerian surrogates for a
separating boundary, the Jacobian determinant and the rate-of-strain
eigenvalue. Chapter~\ref{ch:preliminaries} introduces further
candidates, the Q-criterion and the Hua-Klein criterion among them,
that neither tracker in this dissertation follows. Heteroclinic orbits
connecting distinct critical points through their manifolds, offering
transit routes through a complex field, are also untouched; the
eigenvector-alignment argument Chapter~\ref{ch:first_order} sketches
for detecting such connections between neighboring critical points was
never implemented or tested. More generally, the feature taxonomy of
Chapters~\ref{ch:zeroth} through~\ref{ch:second_order} is not
exhaustive, and other regions of interest in a vector field remain to
be catalogued.

\subsection{Deployment on Surface Vessels}
% Ported from: Draft_11.tex Sec. VII, item d) (lines 1096-1100),
% paired with the Second-Order Hardware Verification limitation above;
% idetc2025_testbed.txt Sec. 6 (Future Work, outdoor surface-vessel and
% UAV testing).

Testing both second-order trackers on the Decabot testbed of
Chapter~\ref{ch:tools}, and then on surface vessels, is the direct
extension the simulation-only status of
Section~\ref{sec:ch9:second_order_hw} above calls for. A Decabot
deployment would use a field emitted by an onboard computer and sampled
by the robots' own sensors, following the pattern
Chapter~\ref{ch:tools} already establishes for the zeroth- and
first-order primitives. A surface-vessel deployment would additionally
confront the advection gap of
Section~\ref{sec:ch9:advection_control} directly, since a vessel's own
motion is no longer decoupled from the current it senses. Outdoor
testing was already proposed independently, for the zeroth- and
first-order primitives, as a way to examine transferability from the
indoor testbed's small scale to a real deployment, and it named both
autonomous surface vessels and UAVs as candidate platforms
\cite{39}; the second-order trackers would follow the same path.

\subsection{Scalar-Field Calibration and a Unified Testbed}
% Ported from: idetc2025_testbed.txt Sec. 6 (Future Work, greyscale
% sensor calibration for scalar fields).

The neural-network calibration technique of Chapter~\ref{ch:tools} was
developed for the hue-and-saturation encoding used throughout this
dissertation's vector-field trials. Extending the same supervised
calibration approach to greyscale sensors would support monochrome
colormaps encoding a scalar field, letting a single testbed and a
single calibration methodology serve both the vector-field primitives
of Chapters~\ref{ch:zeroth} through~\ref{ch:second_order} and the
scalar-field extremum-, contour-, and front-following primitives that
motivate them \cite{39}. No scalar-field calibration or hardware
trial is reported in this dissertation.

\subsection{Okubo-Weiss Corner Branch Selection (Deferred 2026-07-02)}
\label{sec:ch9:ow_corner_branch}

% Research note carried over unchanged from phd_thesis_template.tex's
% prior Chapter 5 (git history, this file). Relocated here because it
% is a specific, evaluable future-work item, not a general limitation.
% Sources: experiments/ow_clean_runs.py and
% experiments/outputs/ow_clean/ under
% trunk/Python_Simulations/Vector_Fields/VF_Robot/, git commits 7dc01a8
% and c19732f, and plan.md (repo root) Phase 3 and Phase 7 entries.

\textbf{Context.} The Okubo-Weiss boundary tracker of the second paper
drives $\hat{D} = \det(\hat{\mathbf{J}})$ to zero and drifts along the
level-set tangent $\mathrm{perp}(\nabla\hat{D})$, sign-aligned with the
ambient flow. On the steady double gyre the $D = 0$ set is not a pair of
closed diamonds but a network of straight lines,
$x_f \pm y_f \in \tfrac{1}{2} + \mathbb{Z}$, which cross at the diamond
corners where $\nabla D = \mathbf{0}$. Noise-free experiments (four
starts) showed lock-on within 13 to 16 control steps and mean $|D|$
between 0.003 and 0.007 on the boundary, 0.3 to 0.8 percent of the well
depth $\pi^4 A^2 = 0.974$, but the tangent follows its line straight
through every corner crossing: the gradient-floor fallback carries the
team across the degenerate point and the Newton action resumes on the
continuing line. Circulation around a single vortex core therefore never
occurs with the published law. The paper presents this straight-through
behavior honestly; the wording of Problem 2 was changed from ``circulate
along'' to ``travel along'' the boundary.

\textbf{The deferred idea.} At an X-crossing the field is locally
$D(\mathbf{p}^* + \mathbf{t}) \approx \tfrac{1}{2}\mathbf{t}^\top
\mathbf{H}_D \mathbf{t}$ with $\mathbf{H}_D$ indefinite, so the zero set
consists of the two asymptote directions $\mathbf{t}$ satisfying
$\mathbf{t}^\top \mathbf{H}_D \mathbf{t} = 0$. In the eigenbasis of
$\mathbf{H}_D$ with $\lambda_1 < 0 < \lambda_2$ these lie at angles
$\pm\arctan\sqrt{-\lambda_1/\lambda_2}$ from the $\lambda_2$
eigenvector. A corner rule can therefore be computed onboard from
$\hat{\mathbf{H}}_D$ alone: inside the gradient-floor ball, evaluate
both estimated asymptotes, discard the one parallel to the incoming
tangent, and exit along the turning branch, sign-chosen by the flow.
Selecting the turning branch at every corner closes the circuit around
one diamond; always selecting the continuing branch reproduces the
published straight-through behavior, so the rule is a second
one-parameter behavior switch, mirroring the park-versus-traverse knob
of the separatrix tracker.

\textbf{Known risk.} The $\hat{\mathbf{H}}_D$ estimate carries a
formation-radius-independent structural bias (third-derivative terms
absent from the quadratic model; see the estimator accuracy results of
the second paper). Eigen-directions survive that bias to within a few
tenths of a degree on the separatrix, but the accuracy of the estimated
asymptote directions at the corners specifically is untested, and near
the corners the eigenvalues themselves are strongly biased. This is the
first thing to measure.

\textbf{Validation plan when revisited.} Implement as a new primitive
(working name \textit{logic\_g\_fable\_step}) alongside the untouched
\textit{logic\_g\_newton\_contour\_pentagon}, following the pattern of
\textit{separatrix\_logic\_fable\_step} (bit-exact equivalence to the
parent law when the knob is off, then the new branch enabled).
Head-to-head at all eight diamond corners plus the domain-center
crossing: success is loop closure around one diamond; metrics are
circumnavigation time and mean $|D|$ per circuit, against the
straight-through baseline of \textit{ow\_clean\_runs.py}.

% ===================================================================
% CHAPTER 10: CONCLUSION
% ===================================================================
\chapter{Conclusion}
\label{ch:conclusion}

\section{Summary of Contributions}

\textbf{Chapter~\ref{ch:tools}} established the Decabot testbed used for
every hardware result in this dissertation: an HSV floor-map encoding,
direction in hue and magnitude in saturation, together with a
supervised neural-network calibration procedure that removed the
RGB-channel interference and the roughly 0.15~Hz discrete-palette
sampling limit of the prior testbed generation \cite{39}. The
calibration's first published iteration reached a cross-trained
accuracy of $R^2 = 0.96$ in direction and $R^2 = 0.91$ in magnitude
\cite{39}; a later recalibration, performed after preventive
maintenance and servicing for the first-order hardware campaign of
Chapter~\ref{ch:first_order}, reached hue RMSE of 0.17~rad
($9.7^\circ$, $R^2 = 0.991$) and saturation RMSE of 4.4\%
($R^2 = 0.951$). This chapter's testbed carried the vector-sum and
vector-to-scalar primitives of Chapter~\ref{ch:zeroth} through their
first hardware demonstrations and underlies every subsequent hardware
result in Chapters~\ref{ch:zeroth} and~\ref{ch:first_order}.

\textbf{Chapter~\ref{ch:preliminaries}} assembled the shared
mathematical background the rest of the dissertation builds on: the
velocity gradient tensor and its six canonical critical-point types,
the rotation-strain (Okubo-Weiss) partition of the plane, Lagrangian
and objective Eulerian coherent structures, and the steady double gyre
used as the running benchmark in Chapters~\ref{ch:second_order}
and~\ref{ch:crosscutting}.

\textbf{Chapter~\ref{ch:estimation}} stated once the machinery the
papers each re-derived: the three-layer control architecture, the
three-robot and six-robot cluster-space formations, and the local
polynomial estimator in general form. A fit of order $k$ has
$(k+1)(k+2)/2$ coefficients per field component, so it needs that many
robots, three for the affine fit and six for the quadratic fit, and it
is singular exactly when the robots lie on a common curve of degree $k$,
a line at first order and a conic at second. Its noise gains scale as
$\rho^{-q}$ for a coefficient of degree $q$, against a truncation error
that grows with the formation radius $\rho$.

\textbf{Chapter~\ref{ch:zeroth}} instantiated this framework at zeroth
order: a single scalar reduction of the sensed flow, direction alone
(vector-sum) or magnitude alone (vector-to-scalar), drives a cluster
toward a feature with no coordinate estimate. Hardware trials on fixed
and sinking vortex fields (three robots, $p=q=0.35$~m,
$\beta = 1.05$~rad, formation held within about 1~cm throughout)
demonstrated both primitives in continuous motion for the first time
\cite{39}. The vector-sum primitive settled within
approximately 0.1~m of the sinking vortex's center but showed a
consistent outward drift, stable at all angles, on the fixed vortex,
a momentum effect absent from the primitive's simulation model; the
vector-to-scalar primitive settled within approximately 0.2~m of the
fixed vortex's center, with performance essentially unchanged under
randomized starting positions \cite{39}. Neither primitive
recovers the critical point's coordinate or supports a commanded
orbital radius, the gap Chapter~\ref{ch:first_order} closes.

\textbf{Chapter~\ref{ch:first_order}} showed that three robots in a
non-collinear triangle are both necessary and sufficient to recover a vector field's
local Jacobian from simultaneous measurements alone. This estimate
supports a proportional attraction law and an orbital law, both
validated across eight simulated fields (six canonical, two
sensor-reconstructed) with 100\% convergence, and validated on
hardware with a 100\% success rate over 157 convergence trials at
position errors of 0.012~m or less. Orbital control held commanded
radii within 0.25~m error at small radii and 0.1~m at larger radii in
hardware. Recomputed offline from the logs, the same estimates named the
class of both printed fields correctly in effectively every cycle, once
each map's sensor frame was registered to the motion-capture frame. A
direct comparison showed both zeroth-order primitives of
Chapter~\ref{ch:zeroth} drift outward under orbital command, since
neither has a mechanism to correct radial error against a specific
point; the first-order estimate is what supplies that mechanism.

\textbf{Chapter~\ref{ch:second_order}} extended this to second order:
six robots in a pentagon-plus-center formation, the minimum for the
quadratic fit, recover a flow's local Jacobian, its determinant $D$, and its
Hessian $\mathbf{H}_D$ from one synchronized measurement set. Two
adaptive navigation primitives built on this estimate, the $D$ tracker
and the $s_1$ (rate-of-strain) tracker, acquire and ride a separating
boundary from instantaneous measurements alone, with no pre-straddling
and no integration horizon. Monte Carlo sweeps on the double gyre
showed the $D$ tracker withstands roughly four times the measurement
noise the $s_1$ tracker does before its 50\% success crossing, while
under a rotating observer the ranking reverses: the $s_1$ tracker is
objective and keeps its straddle at any heading drift tested, while the
$D$ tracker's formation loses the straddle once the drift exceeds
roughly $7^\circ$/h. On real Santa Barbara Channel current data, both
trackers followed the same dominant transport corridor over a 28-hour
window to its bifurcation, at mean distances of 1.5~km and 2.1~km from
an FTLE ridge neither tracker could see.

\textbf{Chapter~\ref{ch:crosscutting}} compared the two estimation
orders directly. A feature located as the zero of a quantity built from
$m$-th field derivatives is placed with an error that scales as
$\rho^{-m}\sigma$ over that quantity's transverse slope. For the critical
point $m = 0$, which recovers
$\sigma_p \approx \|\mathbf{J}^{-1}\|\sigma_v/\sqrt{3}$ independent of
formation radius; for both second-order trackers $m = 2$, so their
terminal error grows as $\rho^{-2}$. Under a rotating observer the trace
of the Jacobian is unchanged while its determinant shifts, so the
first-order sink-versus-source distinction survives a frame rotation but
the saddle-versus-node distinction and the critical point location do
not; among the trackers only the $s_1$ tracker is objective. What each
order certifies at the terminal step also differs: the first-order law
converges to its own estimate, limited by stiction, while the
second-order trackers certify only a vanishing gradient of their
surrogate.

\section{Final Thoughts}

The chapters above share one method applied at rising fit order: a
multirobot team recovers local structure in an unmapped field from its
own simultaneous measurements, with no forecast and no prior model, and
uses that structure directly as a control input. Each order needs
exactly as many robots as its fit has coefficients per component, three
for the affine fit and six for the quadratic fit, placed off the fit's
degenerate curve, a line or a conic. The zeroth-order
primitives of Chapter~\ref{ch:zeroth} show what a team gains by
crossing into this framework at all: direction alone drifts on an
orbit it cannot correct, and only an explicit coordinate estimate
supplies the mechanism to hold one.

Between the two estimation chapters, the choice is not which method is
better but which the platform can support. A cluster with an absolute
heading reference, a compass, GNSS heading, or overhead tracking,
should carry the $D$ tracker of Chapter~\ref{ch:second_order}, which
tolerates a noisier sensor. A cluster whose heading is dead-reckoned,
and so drifts, should carry the $s_1$ tracker, which stays objective
past its one seed step.


% ===================================================================
% REFERENCES
% ===================================================================
\clearpage
\singlespacing
\addcontentsline{toc}{chapter}{References}

\begin{thebibliography}{99}

% Consolidated bibliography, IEEE style, numbered in order of first citation.
% Generated 2026-09-29 from bib_sources.tex and bib_idetc.tex (source-prefixed
% keys). The old-key to number map is in sources/bib_key_map.txt. To add a
% reference, append an entry with the next integer key.

\bibitem{1} D. Kularatne and M. A. Hsieh, ``Tracking attracting manifolds in flows,'' \textit{Auton. Robots}, vol. 41, no. 8, pp. 1575--1588, 2017.

\bibitem{2} X. Tricoche, G. Scheuermann, and H. Hagen, "Continuous topology simplification of planar vector fields," in \textit{Proc. IEEE Visualization}, pp. 159--166, 2001.

\bibitem{3} W. G. Pichel et al., "Marine debris collects within the North Pacific subtropical convergence zone," \textit{Marine Pollution Bulletin}, vol. 54, no. 8, pp. 1207--1211, 2007.

\bibitem{4} M. Mateus, R. Canelas, L. Pinto, and N. Vaz, "When tragedy strikes: potential contributions from ocean observation to search and rescue operations after drowning accidents," \textit{Frontiers in Marine Science}, vol. 7, p. 55, 2020.

\bibitem{5} P. Keramea, K. Spanoudaki, G. Zodiatis, G. Gikas, and G. Sylaios, "Oil spill modeling: A critical review on current trends, perspectives, and challenges," \textit{Journal of Marine Science and Engineering}, vol. 9, no. 2, p. 181, 2021.

\bibitem{6} J. Wang, Y. Lin, R. Liu, and J. Fu, "Odor source localization of multi-robots with swarm intelligence algorithms: A review," \textit{Frontiers in Neurorobotics}, vol. 16, p. 949888, 2022.

\bibitem{7} C. Dong et al., ``Oceanic mesoscale eddies,'' \textit{Ocean-Land-Atmos. Res.}, vol. 4, p. 0081, 2025.

\bibitem{8} Z. Li, X. Wang, J. Hu, F. P. Andutta, and Z. Liu, "A study on an anticyclonic-cyclonic eddy pair off Fraser Island, Australia," \textit{Frontiers in Marine Science}, vol. 7, p. 594358, 2020.

\bibitem{9} R. Deogharia, H. Gupta, S. Sil, A. Gangopadhyay, and A. Shee, "On the evidence of helico-spiralling recirculation within coherent cores of eddies using Lagrangian approach," \textit{Scientific Reports}, vol. 14, no. 1, p. 11014, 2024.

\bibitem{10} T. Peacock and G. Haller, ``Lagrangian coherent structures: The hidden skeleton of fluid flows,'' \textit{Phys. Today}, vol. 66, no. 2, pp. 41--47, 2013.

\bibitem{11} M. Serra et al., ``Search and rescue at sea aided by hidden flow structures,'' \textit{Nat. Commun.}, vol. 11, no. 1, p. 2525, 2020.

\bibitem{12} J. Das et al., ``Coordinated sampling of dynamic oceanographic features with underwater vehicles and drifters,'' \textit{Int. J. Robot. Res.}, vol. 31, no. 5, pp. 626--646, 2012.

\bibitem{13} S. McCammon et al., ``Ocean front detection and tracking using a team of heterogeneous marine vehicles,'' \textit{J. Field Robot.}, vol. 38, no. 6, pp. 854--881, 2021.

\bibitem{14} S. R. Ramp et al., ``Preparing to predict: the second autonomous ocean sampling network (AOSN-II) experiment in the Monterey Bay,'' \textit{Deep-Sea Res. II}, vol. 56, no. 3-5, pp. 68--86, 2009.

\bibitem{15} L. Briñón-Arranz, A. Renzaglia, and L. Schenato, ``Multirobot symmetric formations for gradient and Hessian estimation with application to source seeking,'' \textit{IEEE Trans. Robot.}, vol. 35, no. 3, pp. 782--789, 2019.

\bibitem{16} P. Ögren, E. Fiorelli, and N. E. Leonard, ``Cooperative control of mobile sensor networks: Adaptive gradient climbing in a distributed environment,'' \textit{IEEE Trans. Autom. Control}, vol. 49, no. 8, pp. 1292--1302, 2004.

\bibitem{17} C. A. Kitts, R. T. McDonald, and M. A. Neumann, ``Adaptive navigation control primitives for multirobot clusters: Extrema finding, contour following, ridge/trench following, and saddle point station keeping,'' \textit{IEEE Access}, vol. 6, pp. 17625--17642, 2018.

\bibitem{18} R. T. McDonald et al., ``Navigation of scalar fronts with multirobot clusters in simulation and experiment,'' \textit{IEEE Syst. J.}, vol. 14, no. 3, pp. 3755--3766, 2020.

\bibitem{19} R. T. McDonald, C. A. Kitts, and M. A. Neumann, ``Experimental implementation and verification of scalar field ridge, trench, and saddle point maneuvers using multirobot adaptive navigation,'' \textit{IEEE Access}, vol. 7, pp. 62950--62961, 2019.

\bibitem{20} T. Adamek, C. A. Kitts, and I. Mas, ``Gradient-based cluster space navigation for autonomous surface vessels,'' \textit{IEEE/ASME Trans. Mechatronics}, vol. 20, no. 2, pp. 506--518, 2015.

\bibitem{21} F. Pagano, N. De Carli, E. Restrepo, A. Marino, and P. Robuffo Giordano, "Distributed multi-robot active-sensing of a diffusive source," \textit{IEEE Robotics and Automation Letters}, vol. 10, no. 10, pp. 10807--10814, 2025.

\bibitem{22} S. Li, R. Kong, and Y. Guo, "Cooperative distributed source seeking by multiple robots: Algorithms and experiments," \textit{IEEE/ASME Transactions on Mechatronics}, vol. 19, no. 6, pp. 1810--1820, 2014.

\bibitem{23} Z. Deng and J. Luo, "Fully distributed algorithms for constrained nonsmooth optimization problems of general linear multiagent systems and their application," \textit{IEEE Transactions on Automatic Control}, vol. 69, no. 2, pp. 1377--1384, 2023.

\bibitem{24} W. Li, J. A. Farrell, S. Pang, and R. M. Arrieta, "Moth-inspired chemical plume tracing on an autonomous underwater vehicle," \textit{IEEE Transactions on Robotics}, vol. 22, no. 2, pp. 292--307, 2006.

\bibitem{25} J. Wang et al., "Dynamic plume tracking by cooperative robots," \textit{IEEE/ASME Transactions on Mechatronics}, vol. 24, no. 2, pp. 609--620, 2019.

\bibitem{26} Y. Zhang, M. A. Godin, J. G. Bellingham, and J. P. Ryan, "Using an autonomous underwater vehicle to track a coastal upwelling front," \textit{IEEE Journal of Oceanic Engineering}, vol. 37, no. 3, pp. 338--347, 2012.

\bibitem{27} D. R. Nelson, D. B. Barber, T. W. McLain, and R. W. Beard, "Vector field path following for miniature air vehicles," \textit{IEEE Transactions on Robotics}, vol. 23, no. 3, pp. 519--529, 2007.

\bibitem{28} W. Yao, H. G. de Marina, B. Lin, and M. Cao, "Singularity-free guiding vector field for robot navigation," \textit{IEEE Transactions on Robotics}, vol. 37, no. 4, pp. 1206--1221, 2021.

\bibitem{29} T. Liddy, T.-F. Lu, P. Lozo, and D. Harvey, "Obstacle avoidance using complex vector fields," in \textit{Proc. Australasian Conference on Robotics and Automation}, 2008.

\bibitem{30} D. Panagou, "Motion planning and collision avoidance using navigation vector fields," in \textit{Proc. IEEE International Conference on Robotics and Automation}, 2014, pp. 2513--2518.

\bibitem{31} X. Wang et al., "Adaptive vector field guidance without a priori knowledge of course dynamics and wind," \textit{IEEE/ASME Transactions on Mechatronics}, vol. 27, no. 6, pp. 4597--4607, 2022.

\bibitem{32} T. Inanc, S. C. Shadden, and J. E. Marsden, "Optimal trajectory generation in ocean flows," in \textit{Proc. American Control Conference}, pp. 674--679, 2005.

\bibitem{33} R. N. Smith, M. Schwager, S. L. Smith, B. H. Jones, D. Rus, and G. S. Sukhatme, "Persistent ocean monitoring with underwater gliders: Adapting sampling resolution," \textit{Journal of Field Robotics}, vol. 28, no. 5, pp. 714--741, 2011.

\bibitem{34} M. Michini, M. A. Hsieh, E. Forgoston, and I. B. Schwartz, ``Robotic tracking of coherent structures in flows,'' \textit{IEEE Trans. Robot.}, vol. 30, no. 3, pp. 593--603, 2014.

\bibitem{35} M. Michini, H. Rastgoftar, M. A. Hsieh, and S. Jayasuriya, ``Distributed formation control for collaborative tracking of manifolds in flows,'' in \textit{Proc. Amer. Control Conf.}, pp. 3874--3880, 2014.

\bibitem{36} M. Michini, M. A. Hsieh, E. Forgoston, and I. B. Schwartz, ``Experimental validation of robotic manifold tracking in gyre-like flows,'' in \textit{Proc. IEEE/RSJ Int. Conf. Intell. Robots Syst.}, pp. 2306--2311, 2014.

\bibitem{37} G. Knizhnik, P. Li, X. Yu, and M. A. Hsieh, "Flow-based control of marine robots in gyre-like environments," in \textit{Proc. IEEE International Conference on Robotics and Automation (ICRA)}, pp. 3047--3053, 2022.

\bibitem{38} R. Cooper et al., "Initial Study of Multirobot Adaptive Navigation for Exploring Environmental Vector Fields," in \textit{Proc. IEEE Aerospace Conference}, 2019.

\bibitem{39} C. Waight and C. A. Kitts, ``A functional indoor testbed for multirobot adaptive navigation in vector field environments,'' in \textit{Proc. ASME Int. Des. Eng. Tech. Conf.}, vol. 89251, 2025.

\bibitem{40} C. Waight and C. A. Kitts, ``Adaptive navigation of multirobot systems to critical points in 2D vector fields,'' submitted for publication.

\bibitem{41} S. Tomer et al., ``A low-cost indoor testbed for multirobot adaptive navigation research,'' in \textit{Proc. IEEE Aerosp. Conf.}, pp. 1--12, 2018.

\bibitem{42} D. Rajabhandharaks, R. T. McDonald, M. A. Neumann, and C. A. Kitts, "Indoor testbed for vector field multirobot adaptive navigation: Feasibility study," in \textit{Proc. ASME Int. Design Engineering Technical Conf. Computers and Information in Engineering Conf.}, vol. 59292, p. V009T12A043, 2019.

\bibitem{43} A. Okubo, ``Horizontal dispersion of floatable particles in the vicinity of velocity singularities such as convergences,'' \textit{Deep-Sea Res.}, vol. 17, no. 3, pp. 445--454, 1970.

\bibitem{44} J. Weiss, ``The dynamics of enstrophy transfer in two-dimensional hydrodynamics,'' \textit{Physica D}, vol. 48, no. 2-3, pp. 273--294, 1991.

\bibitem{45} J. C. R. Hunt, A. A. Wray, and P. Moin, ``Eddies, streams, and convergence zones in turbulent flows,'' in \textit{Proc. Summer Program, Center Turbul. Res.}, Rep. CTR-S88, pp. 193--208, 1988.

\bibitem{46} B.-L. Hua and P. Klein, ``An exact criterion for the stirring properties of nearly two-dimensional turbulence,'' \textit{Physica D}, vol. 113, no. 1, pp. 98--110, 1998.

\bibitem{47} G. Haller and G. Yuan, ``Lagrangian coherent structures and mixing in two-dimensional turbulence,'' \textit{Physica D}, vol. 147, no. 3-4, pp. 352--370, 2000.

\bibitem{48} G. Haller, ``Lagrangian coherent structures,'' \textit{Annu. Rev. Fluid Mech.}, vol. 47, pp. 137--162, 2015.

\bibitem{49} S. C. Shadden, F. Lekien, and J. E. Marsden, ``Definition and properties of Lagrangian coherent structures from finite-time Lyapunov exponents in two-dimensional aperiodic flows,'' \textit{Physica D}, vol. 212, no. 3-4, pp. 271--304, 2005.

\bibitem{50} P. J. Nolan, M. Serra, and S. D. Ross, ``Finite-time Lyapunov exponents in the instantaneous limit and material transport,'' \textit{Nonlinear Dyn.}, vol. 100, no. 4, pp. 3825--3852, 2020.

\bibitem{51} T. D. Harms, S. L. Brunton, and B. J. McKeon, ``Lagrangian gradient regression for the detection of coherent structures from sparse trajectory data,'' \textit{R. Soc. Open Sci.}, vol. 11, no. 9, p. 240586, 2024.

\bibitem{52} M. Serra and G. Haller, ``Objective Eulerian coherent structures,'' \textit{Chaos}, vol. 26, no. 5, p. 053110, 2016.

\bibitem{53} C. A. Kitts and I. Mas, ``Cluster space specification and control of mobile multirobot systems,'' \textit{IEEE/ASME Trans. Mechatronics}, vol. 14, no. 2, pp. 207--218, 2009.

\bibitem{54} I. Mas and C. Kitts, "Obstacle avoidance policies for cluster space control of nonholonomic multirobot systems," \textit{IEEE/ASME Transactions on Mechatronics}, vol. 17, no. 6, pp. 1068--1079, 2012.

\bibitem{55} K. Garg and D. Panagou, "Finite-time estimation and control for multi-aircraft systems under wind and dynamic obstacles," \textit{Journal of Guidance, Control, and Dynamics}, vol. 42, no. 7, pp. 1489--1505, 2019.

\bibitem{56} Y. Jiao, H. Hang, J. Merel, and E. Kanso, "Sensing flow gradients is necessary for learning autonomous underwater navigation," \textit{Nature Communications}, vol. 16, no. 1, p. 3044, 2025.

\bibitem{57} D. Eberly et al., ``Ridges for image analysis,'' \textit{J. Math. Imaging Vis.}, vol. 4, no. 4, pp. 353--373, 1994.

\bibitem{58} NOAA Nat. Centers Environ. Inf., ``Near-real-time surface ocean velocities derived from HF-radar stations,'' NCEI Accession gov.noaa.nodc:IOOS-HFRadarRTVector.

\bibitem{59} S. T. Hart and C. A. Kitts, "Unifying control architecture for reactive particle swarms," \textit{IEEE/ASME Transactions on Mechatronics}, vol. 28, no. 2, pp. 873--883, 2022.


\end{thebibliography}

% ===================================================================
% APPENDICES
% ===================================================================
\appendix

\chapter{Canonical Field Definitions}
\label{app:fields}

\section{Six Canonical Critical Point Fields}
\label{sec:appA:canonical}

% Ported from: cwaight_systems_paper.tex
Six vector field environments centered at the origin were used to evaluate the first-order primitives of Chapter~\ref{ch:first_order}. Let $r = \sqrt{x^2 + y^2}$ and $\theta = \text{atan2}(y, x)$. The base fields are the vortex (pure rotational flow), the sink (radial inward flow) whose negation $\mathbf{v}_{\text{source}} = -\mathbf{v}_{\text{sink}}$ is the source (radial outward flow), and the saddle, an axis-aligned saddle rotated by $45^\circ$ so that its stable and unstable directions lie along the diagonals $y = x$ and $y = -x$:
\begin{equation}\label{eq:appA:14}
\mathbf{v}_{\text{vortex}} = \begin{bmatrix} -r \sin(\theta) \\ r \cos(\theta) \end{bmatrix}, \;\;
\mathbf{v}_{\text{sink}} = \begin{bmatrix} -x \\ -y \end{bmatrix}, \;\;
\mathbf{v}_{\text{saddle}} = \begin{bmatrix} -y \\ -x \end{bmatrix}
\end{equation}

% Ported from: cwaight_systems_paper.tex
The sinking vortex (spiraling inward) and spewing vortex (spiraling outward) are linear combinations of the base fields:
\begin{equation}\label{eq:appA:15}
\begin{aligned}
\mathbf{v}_{\text{sink-vortex}} &= 0.4\,\mathbf{v}_{\text{vortex}} + 0.15\,\mathbf{v}_{\text{sink}} \\
\mathbf{v}_{\text{spew-vortex}} &= -0.4\,\mathbf{v}_{\text{vortex}} - 0.15\,\mathbf{v}_{\text{sink}}
\end{aligned}
\end{equation}

\section{Testbed Vortex Fields}
\label{sec:appA:testbed}

% New text
The fixed and sinking vortex maps of Chapter~\ref{ch:zeroth} are given by
(\ref{eq:ch5:rtheta})--(\ref{eq:ch5:sinking}). The fixed vortex there is
the linear clockwise rotation $(y - c_y,\, -(x - c_x))$, the vortex of
(\ref{eq:appA:14}) with the opposite sense. Its sinking component falls
off as $1/r$, so it differs from the linear sinking vortex of
(\ref{eq:appA:15}) and is singular at the center apart from the
$10^{-6}$ regularization.

\section{Steady Double Gyre}
\label{sec:appA:dg}

% New text
The steady double gyre of Chapter~\ref{ch:second_order} is defined in
(\ref{eq:ch3:dg_u}), with its determinant, Hessian, and smaller strain
eigenvalue in closed form in Section~\ref{sec:ch3:dg_example}.

\begin{table}[htbp]
\centering
\caption{Fields used in this dissertation.}
\label{tab:appA:summary}
\begin{tabular}{lll}
\hline
Field & Critical points & Chapter \\
\hline
Vortex, sink, source, saddle & center, stable node, unstable node, saddle & \ref{ch:first_order} \\
Sinking and spewing vortex (linear) & stable and unstable spiral & \ref{ch:first_order} \\
Fixed vortex (testbed) & center & \ref{ch:zeroth} \\
Sinking vortex (testbed, $1/r$ sink) & singular sink at the center & \ref{ch:zeroth} \\
Steady double gyre & two saddles, two centers & \ref{ch:second_order} \\
Santa Barbara Channel record & measured, time-varying & \ref{ch:second_order} \\
\hline
\end{tabular}
\end{table}

\chapter{Three-Robot Cluster Kinematics}
\label{app:kin3}

% Ported from: cwaight_systems_paper.tex
\section{Forward Kinematics}

The forward kinematics transforms robot positions $(x_1, y_1)$, $(x_2, y_2)$, $(x_3, y_3)$ to the SAS parameters and centroid coordinates. The centroid position is:
\begin{equation}\label{eq:appB:20}
x_c = \frac{x_1 + x_2 + x_3}{3}, \quad y_c = \frac{y_1 + y_2 + y_3}{3}
\end{equation}
The side lengths, with $\mathbf{p}_i = (x_i, y_i)$ the position of robot $i$, are:
\begin{equation}\label{eq:appB:21}
p = \lVert \mathbf{p}_2 - \mathbf{p}_1 \rVert, \quad q = \lVert \mathbf{p}_3 - \mathbf{p}_2 \rVert
\end{equation}
The interior angle at robot 2 is:
\begin{equation}\label{eq:appB:23}
\beta = \arccos\left(\frac{(\mathbf{p}_1 - \mathbf{p}_2) \cdot (\mathbf{p}_3 - \mathbf{p}_2)}{pq}\right)
\end{equation}

The orientation angle is:
\begin{equation}\label{eq:appB:24}
\theta_c = \text{atan2}(y_1 - y_2, x_1 - x_2)
\end{equation}

\section{Inverse Kinematics}

Given the SAS parameters $(p, \beta, q)$ and centroid parameters $(x_c, y_c, \theta_c)$, the robot positions are recovered by constructing the triangle in a local frame and then applying rotation and translation.

In the local frame with robot 2 at the origin, the triangle is constructed as:
\begin{equation}\label{eq:appB:25}
\begin{aligned}
(x_1, y_1)^{\text{local}} &= (p,\, 0), \qquad (x_2, y_2)^{\text{local}} = (0,\, 0), \\
(x_3, y_3)^{\text{local}} &= (q\cos\beta,\, q\sin\beta)
\end{aligned}
\end{equation}

This construction assumes robot 3 lies counterclockwise from the robot 1--robot 2 edge; since $\beta$ from (\ref{eq:appB:23}) lies in $(0, \pi)$, the mirror configuration (robot 3 clockwise of the edge) maps to the same SAS parameters.

The local centroid is:
\begin{equation}\label{eq:appB:28}
x_c^{\text{local}} = \frac{x_1^{\text{local}} + x_2^{\text{local}} + x_3^{\text{local}}}{3}, \quad y_c^{\text{local}} = \frac{y_1^{\text{local}} + y_2^{\text{local}} + y_3^{\text{local}}}{3}
\end{equation}

After centering the triangle at the local origin and applying rotation $\theta_c$ and translation to $(x_c, y_c)$, the global positions become:
\begin{equation}\label{eq:appB:29}
\begin{bmatrix} x_i \\ y_i \end{bmatrix} = \begin{bmatrix} \cos\theta_c & -\sin\theta_c \\ \sin\theta_c & \cos\theta_c \end{bmatrix} \begin{bmatrix} x_i^{\text{local}} - x_c^{\text{local}} \\ y_i^{\text{local}} - y_c^{\text{local}} \end{bmatrix} + \begin{bmatrix} x_c \\ y_c \end{bmatrix}
\end{equation}
for $i \in \{1, 2, 3\}$.

\section{Inverse Jacobian}

The inverse Jacobian maps shape parameter velocities to robot velocities:
\begin{equation}\label{eq:appB:30}
\begin{bmatrix}
\delta x_1 \\ \delta y_1 \\ \delta x_2 \\ \delta y_2 \\ \delta x_3 \\ \delta y_3
\end{bmatrix} =
\mathbf{J}_c^{-1}
\begin{bmatrix}
\delta x_c \\ \delta y_c \\ \delta \theta_c \\ \delta p \\ \delta \beta \\ \delta q
\end{bmatrix}
\end{equation}

where $\mathbf{J}_c^{-1}$ is a $6 \times 6$ matrix containing partial derivatives $\frac{\partial x_i}{\partial \xi_j}$ and $\frac{\partial y_i}{\partial \xi_j}$ for robot positions and shape parameters $\xi_j \in \{x_c, y_c, \theta_c, p, \beta, q\}$. This matrix is computed analytically by taking partial derivatives of the inverse kinematics equations (\ref{eq:appB:29}) with respect to each shape parameter. The explicit entries are not listed here. The simulation software computes them by forward finite differences of the inverse kinematics (\texttt{compute\_inverse\_jacobian} in \texttt{src/control/kinematics.py}), and the \texttt{cluster\_builder} tool generates them in closed symbolic form (\texttt{clusterbuilder.py} with the \texttt{--symbolic} option).

\chapter{Six-Robot Cluster Kinematics}
\label{app:kin6}

% New text. Checked against src/control/pentagon_kinematics.py
% (VF_Robot), which implements this parameterization.
The six-robot cluster of Chapter~\ref{ch:second_order} is a cluster of
clusters. The robots form three pairs, and the pair midpoints form a
triangle described with the same SAS variables as the three-robot cluster
of Appendix~\ref{app:kin3}. The estimator of
Chapter~\ref{ch:second_order} reads only the six measured positions, so
this parameterization matters to the formation controller alone.

\section{Forward Kinematics}
\label{sec:appC:forward}

Pair A holds robots 1 and 2, pair B robots 3 and 4, and pair C robots 5
and 6. The pair midpoints are
\begin{equation}
    \mathbf{m}_A = \tfrac{1}{2}(\mathbf{p}_1 + \mathbf{p}_2), \quad
    \mathbf{m}_B = \tfrac{1}{2}(\mathbf{p}_3 + \mathbf{p}_4), \quad
    \mathbf{m}_C = \tfrac{1}{2}(\mathbf{p}_5 + \mathbf{p}_6),
    \label{eq:appC:midpoints}
\end{equation}
and the cluster centroid is their mean,
$\mathbf{p}_c = \tfrac{1}{3}(\mathbf{m}_A + \mathbf{m}_B + \mathbf{m}_C)$.
The midpoint triangle has sides
$p_1 = \lVert\mathbf{m}_B - \mathbf{m}_A\rVert$ and
$q_1 = \lVert\mathbf{m}_C - \mathbf{m}_B\rVert$, with $\beta_1$ the
interior angle at $\mathbf{m}_B$, and the heading $\theta_c$ is the
direction from $\mathbf{m}_A$ to $\mathbf{m}_B$. Each pair has a length
and an orientation,
\begin{equation}
    L_2 = \lVert\mathbf{p}_2 - \mathbf{p}_1\rVert, \qquad
    \theta_2 = \operatorname{atan2}(y_2 - y_1,\, x_2 - x_1),
    \label{eq:appC:pair}
\end{equation}
and likewise $(L_3, \theta_3)$ for pair B and $(L_4, \theta_4)$ for pair C.
The twelve cluster variables are
\begin{equation}
    (x_c,\, y_c,\, \theta_c,\, p_1,\, \beta_1,\, q_1,\,
     L_2,\, \theta_2,\, L_3,\, \theta_3,\, L_4,\, \theta_4).
    \label{eq:appC:state}
\end{equation}

\section{Inverse Kinematics}
\label{sec:appC:inverse}

The inverse kinematics builds the midpoint triangle from
$(p_1, \beta_1, q_1)$ in a local frame, as in
Appendix~\ref{app:kin3}, then rotates it by the heading and translates it
to $(x_c, y_c)$. Each pair's two robots are then placed at its midpoint
plus and minus $\tfrac{1}{2}L_k(\cos\theta_k, \sin\theta_k)$.

\section{Inverse Jacobian}
\label{sec:appC:jacobian}

The $12 \times 12$ inverse Jacobian $\mathbf{J}_c^{-1}$ maps the rates of
the twelve cluster variables to the velocities of the six robots. Its
entries are the partial derivatives of the inverse kinematics and are not
listed here. The simulation software computes them analytically
(\texttt{compute\_inverse\_jacobian} in
\texttt{src/control/pentagon\_kinematics.py}), and the
\texttt{cluster\_builder} tool generates them in closed symbolic form
(\texttt{clusterbuilder.py} with the \texttt{--symbolic} option). The
cluster's heading is not regulated, since the formation controller
commands zero spin throughout.

\chapter{Frame Equivariance of the $s_1$ Tracker}
\label{app:frame_equivariance}
% Ported from: Draft_11.tex Sec. "Relation Between the Fields" (D is not
% objective; s1 is objective) and Sec. "Behavior Under a Rotating
% Observer" (the double-gyre rotating-frame experiment and its result).

This appendix collects the statement, and what supporting derivation
exists in the source paper, for why the $s_1$ tracker of
Chapter~\ref{ch:second_order} is objective (frame-indifferent) while the
$D$ tracker is not.

\section{The Determinant Field Is Not Objective}
\label{sec:appD:D_not_objective}

Both the determinant field $D=\det\mathbf{J}$ and the smaller strain
eigenvalue $s_1$ are read from the same fitted Jacobian $\hat{\mathbf{J}}$.
Where the vorticity $\omega = \partial v/\partial x - \partial u/\partial y$
vanishes, $\mathbf{J}$ equals its symmetric part $\mathbf{S}$ and
$D = s_1s_2$, which is $-s_1^2$ in incompressible flow. Along such a
curve the two fields share their trenches. Elsewhere the spin
$\mathbf{W} = \mathbf{J}-\mathbf{S}$, the antisymmetric part of
$\mathbf{J}$ that carries the vorticity, separates them, and it also
separates them across observers.

A frame rotating clockwise at rate $\Omega$, with rotation $\mathbf{Q}$,
adds a uniform spin that raises $\omega$ to $\omega+2\Omega$. The strain
tensor only co-rotates, so $s_1$ is objective \cite{52}, while the
determinant keeps the spin,
\begin{equation}
D' = D + \Omega\,\omega + \Omega^2.
\label{eq:appD:D_not_objective}
\end{equation}
Relatively rotating observers therefore disagree about where the
trenches of $D$ lie.
%% NOTE(new-math, checked 2026-09-29): Draft_11 states
%% eq:appD:D_not_objective without the algebra; the derivation below is new.

To see where Equation~\eqref{eq:appD:D_not_objective} comes from, split
the Jacobian into strain and spin,
\begin{equation}
\mathbf{S} = \mu\mathbf{I} + \begin{bmatrix} s_n & s_s \\ s_s & -s_n\end{bmatrix},
\qquad
\mathbf{W} = \frac{\omega}{2}\begin{bmatrix} 0 & -1 \\ 1 & 0\end{bmatrix},
\label{eq:appD:split}
\end{equation}
with $\mu = \tfrac12\operatorname{div}\mathbf{v}$ and
$r = \sqrt{s_n^2+s_s^2}$, so the strain eigenvalues are $s_{1,2} = \mu \mp r$.
Expanding the determinant,
\begin{equation}
D = \det(\mathbf{S}+\mathbf{W}) = \mu^2 - r^2 + \frac{\omega^2}{4}
  = s_1 s_2 + \frac{\omega^2}{4}.
\label{eq:appD:D_split}
\end{equation}
The rotating observer of Equation~\eqref{eq:appD:rotating_field} adds the
field $\Omega(-y',x')^{\mathsf T}$, whose gradient is purely antisymmetric,
so it leaves $\mathbf{S}$ unchanged up to the co-rotation
$\mathbf{S}' = \mathbf{Q}\mathbf{S}\mathbf{Q}^{\mathsf T}$, which preserves
$\mu$ and $r$, and replaces $\omega/2$ by $\omega/2 + \Omega$. Substituting
into Equation~\eqref{eq:appD:D_split},
\begin{equation}
D' = \mu^2 - r^2 + \Bigl(\frac{\omega}{2}+\Omega\Bigr)^2
   = D + \Omega\,\omega + \Omega^2 ,
\end{equation}
while $s_1' = s_1$ and $s_2' = s_2$. The $s_1$ tracker's gradient and
Hessian are those of an observer-invariant scalar field, and its tangent
reference $\mathbf{t}_{\mathrm{ref}}$ is carried in state rather than read
from a world axis, so its commanded velocity co-rotates with the frame.

\section{Behavior Under a Rotating Observer}
\label{sec:appD:rotating}

Both trackers were run from $(0.05,0.40)$, just inside the upper saddle
of the double gyre, on the same physical flow in two observer frames: the
inertial one and a frame rotating clockwise at $\Omega=0.23$, a heading
drift of about $8^\circ$/h at the ocean time unit of
Chapter~\ref{ch:second_order}, where the field observed in the rotating
frame is
\begin{equation}
\mathbf{v}'(\mathbf{p}',t) = \mathbf{Q}\mathbf{v}(\mathbf{Q}^\top\mathbf{p}') + \Omega\,(-y',x')^\top.
\label{eq:appD:rotating_field}
\end{equation}
The rate keeps the structure's apparent speed along the tracked path,
$\Omega\lVert\mathbf{p}\rVert \leq 0.11$, inside the command authority
$k\,c_{\max}=0.12$, the condition required by
Appendix~\ref{app:stability}. Rotating-frame trajectories are pulled back
to inertial coordinates and compared with the inertial run of the same
primitive.

The trajectories are shown in Figure~\ref{fig:ch7:oecs_objectivity}.

Between acquisition and the far saddle, the $s_1$ tracker stays within
0.016 of the separatrix in the rotating frame, against 0.010 in the
inertial frame, keeps the straddle throughout, and both runs capture at
the same bottom saddle. The $D$ tracker's rotating-frame run is displaced
by up to 0.084, against 0.010, and loses the straddle for 36 steps across
the crest, with all six robots on one side of the separatrix. Sweeping
the rate, the $D$ tracker first loses the straddle at
$\Omega \approx 0.21$, where $\Omega\lVert\mathbf{p}\rVert$ at the crest
is a sixth of the command authority.

The frame's solid-body vorticity shifts the $D'$ landscape
(Equation~\eqref{eq:appD:D_not_objective}) and adds a transport term to
the measured flow, so both terms of the $D$ tracker are steered by a
frame artifact. On the separatrix $\omega=0$ but
$\partial_x\omega = 2\pi^3A\sin\pi y_f$, so $\Omega\omega$ tilts $D'$
across the trench. Against the transverse curvature $2\pi^6A^2$, it moves
the trench by $\Omega\sin(\pi y_f)/(\pi^3A)$, largest at the origin and
zero at the saddles, so the run rejoins the separatrix near
$\mathbf{p}_1^*$.

\section{Why the $s_1$ Command Is Equivariant}
\label{sec:appD:s1_equivariant}

The $s_1$ tracker's tests are scalar inequalities in quantities a
rotation leaves invariant, so its command co-rotates with the frame, and
its residual error is the frame-transport term bounded in
Appendix~\ref{app:stability}. Its one non-objective input is the tangent
seed taken from the measured flow, which fixes the ride direction on the
first ride step (Equation~\eqref{eq:ch7:tangent_select}); the frames agree
while
\begin{equation}
\Omega\lVert\mathbf{p}_{\text{seed}}\rVert < \bigl\lvert\hat{\mathbf{v}}_0^\top\mathbf{t}_{\text{raw}}\bigr\rvert.
\label{eq:appD:seed_condition}
\end{equation}
Past that one seed step, the tangent is carried in state through
$\mathbf{t}_{\mathrm{ref}}$ (Equation~\eqref{eq:ch7:tangent_select}), which is
itself built from objective quantities, so the $s_1$ tracker is objective
except for this single, bounded, first-step dependence on the measured
flow.
%% NOTE: \ref{eq:ch7:tangent_select} matches the eq:ch7:-prefixed label
%% currently defined in ch07_second_order.tex for the $s_1$ tracker's
%% tangent-selection rule; re-check if that chapter's labels change.

\chapter{Transverse Stability Proofs}
\label{app:stability}
% Ported from: Draft_11.tex Appendix "Stability of the Trackers" (its
% only stability appendix), in full. thesis.tex's older Sec. "Stability
% Analysis" (a superset of proofs for an earlier, three-mode controller
% design that Draft_11 replaced with the simpler D/s1 tracker pair) is
% NOT ported here; see the task report for the list of proofs it drops.

Both arguments assume the cluster realizes $\dot{\mathbf{p}}_c$
\cite{17,53} and neglect robot lag. Near a trench, with $n$ the
distance across it and $\ell$ the arclength along it, the surrogate is
$f = h(\ell) + \tfrac12\kappa_\perp(\ell)n^2$ to second order, with
$\kappa_\perp > 0$.

On a quadratic model with Hessian eigenpair $(\lambda_j,\mathbf{w}_j)$,
the gradient component $g_j = \mathbf{w}_j^\top\nabla f$ obeys
$\dot g_j = \lambda_j\mathbf{w}_j^\top\dot{\mathbf{p}}_c$, so the
saturated Newton channel
$\mathbf{w}_j^\top\dot{\mathbf{p}}_c = -k\,\operatorname{sat}(x_j)$,
$x_j = g_j/\lambda_j$, gives
\begin{equation}
\frac{d}{dt}\,\tfrac12g_j^2 = -k\,\lambda_j^2\,x_j\,\operatorname{sat}(x_j) \leq 0
\label{eq:appE:newton_lyap}
\end{equation}
for either sign of $\lambda_j$, with equality only at $g_j=0$. A full
Newton step would therefore halt at the crest, so the $D$ tracker keeps
only the across-trench channel.

\begin{proposition}
\label{prop:appE:transverse}
Let $\dot n = -k\,\operatorname{sat}(a_\perp n + e)$ with $a_\perp > 0$
and $|e| \leq \delta$. Then $\limsup_{t\to\infty}|n| \leq \delta/a_\perp$.
If $e=0$, $|n|$ falls below $c_{\max}/a_\perp$ in finite time, then decays
at rate at least $k\,a_\perp\tanh 1$.
\end{proposition}

\begin{proof}
With $V = \tfrac12n^2$, $\dot V = -k\,n\,\operatorname{sat}(a_\perp n+e) <
0$ for $|n| > \delta/a_\perp$, where the argument of $\operatorname{sat}$
has the sign of $n$. With $e=0$, $|\dot n| \geq k\,c_{\max}\tanh 1$ while
$a_\perp|n| \geq c_{\max}$, and $|\tanh x| \geq |x|\tanh 1$ for $|x|\leq1$
gives the rate.
\end{proof}

Let the fitted transverse gradient be $\kappa_\perp n + \epsilon$,
$|\epsilon| \leq \delta_g$. In the $D$ tracker's law (Chapter
\ref{ch:second_order}, Eq.~\ref{eq:ch7:d_law}), $a_\perp =
\kappa_\perp/\hat\kappa_\perp$ and $e = \epsilon/\hat\kappa_\perp$. In the
$s_1$ tracker's law (Eq.~\ref{eq:ch7:s1_law}), $a_\perp = g_\perp\kappa_\perp$
and $e = g_\perp\epsilon$. Both bounds equal $\delta_g/\kappa_\perp$,
independent of gain and fitted curvature, but only the $D$ tracker needs
$\hat\kappa_\perp > 0$. Where $\kappa_\perp$ vanishes over a length
$\Delta\ell$, as at the fold, the $s_1$ ride crosses it in time
$\Delta\ell/(k\,c_{\max}\tanh 1)$ and $|n|$ grows by at most
$\Delta\ell/\tanh 1$. A frame rotating at $\Omega$ perturbs its $\dot n$
by at most $\Omega\lVert\mathbf{p}\rVert$, leaving $n$ bounded while
$\Omega\lVert\mathbf{p}\rVert < k\,c_{\max}$
(Appendix~\ref{app:frame_equivariance}).

Both branches of the $D$ tracker's along-trench speed law
(Eq.~\ref{eq:ch7:vpar}) are nonnegative, so the $D$ tracker never reverses
along the trench. If $\operatorname{sat}(v_\parallel) \geq v_{\min} > 0$
outside the flow-saddle neighborhoods where its capture test
(Eq.~\ref{eq:ch7:d_capture}) applies, it crosses a segment of length
$L_\Gamma$ within $L_\Gamma/(k\,v_{\min})$.
%% NOTE: the \ref{eq:ch7:d_law}, \ref{eq:ch7:s1_law}, \ref{eq:ch7:vpar},
%% and \ref{eq:ch7:d_capture} cross-references above match the
%% eq:ch7:-prefixed labels currently defined in ch07_second_order.tex;
%% re-check them if that chapter's labels change.

\chapter{Statistical Methods}
\label{app:statistics}

% New text
\section{Bias and Precision}
Chapter~\ref{ch:first_order} summarizes each set of convergence trials by
its final centroid positions. The systematic bias is the distance from
the mean final position to the true critical point, and it measures what
averaging cannot remove. The precision is the standard deviation of the
final positions, and it measures repeatability. Orbital trials are
summarized by the mean and standard deviation of the radial error.

\section{Circular Statistics for Hue}
Hue is an angle, so a reading near 0 and a reading near 1 are the same
direction. The hue network of Section~\ref{sec:ch2:calibration} predicts
the sine and cosine of hue for this reason, and its accuracy is reported
with circular statistics.

\section{Success Rates}
The noise sweeps of Chapter~\ref{ch:second_order} report success rates
over $N = 10^4$ trials per noise level. Each trial succeeds or fails
independently, so the standard error of a rate $\hat{p}$ is
$\sqrt{\hat{p}(1-\hat{p})/N} \leq 1/(2\sqrt{N}) = 0.005$.


\end{document}
